\documentclass[11pt]{article}
% Shared preamble for all solution files.
\usepackage[margin=1.1in]{geometry}
\usepackage{enumitem}
\usepackage{tikz}
\usepackage[hidelinks]{hyperref}
\usepackage{amsmath,amssymb,amsthm}
\newcommand{\fa}{\mathfrak a}
\newcommand{\fb}{\mathfrak b}
\newcommand{\fc}{\mathfrak c}
\newcommand{\fp}{\mathfrak p}
\newcommand{\fq}{\mathfrak q}
\newcommand{\fm}{\mathfrak m}
\newcommand{\fn}{\mathfrak n}
\newcommand{\fr}{\mathfrak r}
\newcommand{\Nil}{\mathfrak N}
\newcommand{\Jac}{\mathfrak R}
\newcommand{\rad}{r}
\newcommand{\Spec}{\operatorname{Spec}}
\newcommand{\Max}{\operatorname{Max}}
\newcommand{\Ann}{\operatorname{Ann}}
\newcommand{\Ker}{\operatorname{Ker}}
\newcommand{\Z}{\mathbf Z}
\newcommand{\Q}{\mathbf Q}
\newcommand{\R}{\mathbf R}
\newcommand{\C}{\mathbf C}
\newcommand{\ci}{\mathrm{c}}
\newcommand{\ex}{\mathrm{e}}

\setlength{\parindent}{0pt}
\setlength{\parskip}{0.55em}
% \exercise{label}{statement}  -- followed by the solution text
\newcommand{\exercise}[2]{\par\bigskip\noindent{\large\textbf{Exercise #1}}\par\smallskip
  \begin{quote}\small\textbf{Problem.}\ #2\end{quote}\par}
\newcommand{\pf}{\noindent\textbf{Solution.}\ }

\title{Liu, Chapter 2: General properties of schemes\\[2pt]\large Solutions to the exercises}
\author{}
\date{}
\begin{document}
\maketitle
\noindent\textbf{Conventions.} Rings are commutative with $1$. Results of the text are cited by their numbers in Liu (``Prop.~3.14'', ``Lemma~2.7'', \dots). Exercise labels are those of the book, numbered within each section of the chapter: for instance Exercise~3.12 is the twelfth exercise at the end of Section 2.3. Each exercise is restated before its solution; sub-items that the book displays as nested lists are written as running prose.

Honest notes. Several printed statements need a correction or an interpretation, and each solution says what is proved: Ex.~1.4 (a nonnoetherian counterexample to nilpotency of $\mathfrak pA_{\mathfrak p}$), Ex.~2.3(c) (false for a non-separated presheaf, proved for separated ones, with a counterexample), Ex.~3.3(b) and Ex.~3.10 (a closure bar and the case $A=0$), Section 2.4 (the empty open set in a statement with $U=\emptyset$, and the meaning of ``indecomposable''), Ex.~5.5 (false if the ideal contains nonzero elements of degree $0$). Some solutions rely on standard facts named where used.

\section*{2.1 Spectrum of a ring}

Throughout, rings are commutative with $1$. Numbers such as ``Lemma~1.6'', ``Cor.~1.12'' refer to the results of Section 2.1 itself (Chapter 2). We use freely that a finitely generated algebra over a field is noetherian (Hilbert basis theorem), and that the integral elements of a ring extension form a subring.

\exercise{1.1}{Let $A=k[[T]]$ be the ring of formal power series with coefficients in a field $k$. Determine $\Spec A$.}
\pf
\emph{Answer.} $\Spec A=\{0,\ TA\}$ has exactly two points: the generic point $0$ and the closed point $TA$; the closed subsets are $\emptyset$, $\{TA\}$ and $\Spec A$ (so the nonempty open sets are $\{0\}$ and $\Spec A$).

\emph{Units.} Let $a=\sum a_iT^i\in A$ with $a_0\neq0$. We solve $ab=1$ for $b=\sum b_iT^i$: put $b_0=a_0^{-1}$ and, recursively, $b_n=-a_0^{-1}\sum_{i=1}^{n}a_ib_{n-i}$. Thus $a$ is a unit. Conversely a unit has nonzero constant term. Hence the non-units are exactly the series with $a_0=0$, i.e. the elements of $TA$.

\emph{Nonzero elements.} Every $a\neq0$ can be written $a=T^nu$ with $n\ge0$ and $u$ a unit (take $n$ the smallest index with $a_n\neq0$ and divide by $T^n$). So every ideal $I\neq0$ equals $T^nA$ with $n$ the minimal order of its elements: indeed $T^nu\in I$ gives $T^n\in I$, and every element of $I$ is a multiple of $T^n$.

\emph{Primes.} $0$ is prime since $A$ is a domain. $TA$ is prime (maximal), since $A/TA\simeq k$. If $n\ge2$, then $T^nA$ is not prime ($T\cdot T^{n-1}\in T^nA$ but $T,T^{n-1}\notin T^nA$), and $A=T^0A$ is not prime by convention. So $\Spec A=\{0,TA\}$.

\emph{Topology.} $V(0)=\Spec A$, $V(A)=\emptyset$, and $V(T^nA)=\{TA\}$ for $n\ge1$ because every prime containing $T^n$ contains $T$. Hence the closed sets are $\emptyset,\{TA\},\Spec A$, and $\overline{\{0\}}=\Spec A$ while $\{TA\}$ is closed. \hfill$\square$

\exercise{1.2}{Let $\varphi:A\to B$ be a homomorphism of finitely generated algebras over a field. Show that the image of a closed point under $\Spec\varphi$ is a closed point.}
\pf
\emph{Remark on the hypothesis.} ``Homomorphism of algebras over a field'' must mean a homomorphism of algebras over the same field $k$, i.e. a $k$-linear ring homomorphism. For a mere ring homomorphism the statement is false: $\Q[x]$ and $\Q(t)$ are finitely generated algebras over the fields $\Q$ and $\Q(t)$ respectively, and the ring homomorphism $\Q[x]\to\Q(t)$, $x\mapsto t$, sends the closed point $0$ of $\Spec\Q(t)$ to the non-closed point $0$ of $\Spec\Q[x]$.

So let $\varphi$ be a $k$-algebra homomorphism, let $\fm\subset B$ be a maximal ideal (a closed point, Remark~1.3) and $\fp=\varphi^{-1}(\fm)$, which is a prime ideal. By Cor.~1.12, $B/\fm$ is a finite extension of $k$, so $\dim_kB/\fm<\infty$. The map $\varphi$ induces an injective $k$-algebra homomorphism $A/\fp\hookrightarrow B/\fm$. Hence $D:=A/\fp$ is a domain which is a finite dimensional $k$-vector space. For $0\neq d\in D$, multiplication by $d$ is an injective $k$-linear endomorphism of $D$, hence surjective, so $de=1$ for some $e$. Thus $D$ is a field and $\fp$ is maximal, i.e. $\fp$ is a closed point of $\Spec A$. \hfill$\square$

\exercise{1.3}{Let $k=\R$ and $A=k[X,Y]/(X^2+Y^2+1)$. Let $x,y$ be the images of $X,Y$ in $A$.
(a) Let $\fm$ be a maximal ideal of $A$. Show that there exist $a,b,c,d\in k$ such that $x^2+ax+b,\ y^2+cy+d\in\fm$. Using $x^2+y^2+1=0$, show that $\fm$ contains an element $f=\alpha x+\beta y+\gamma$ with $(\alpha,\beta)\neq(0,0)$. Deduce $\fm=fA$.
(b) Show that $(\alpha,\beta,\gamma)\mapsto(\alpha x+\beta y+\gamma)A$ is a bijection from $\mathbf P(k^3)\setminus\{(0,0,1)\}$ onto the set of maximal ideals of $A$.
(c) Let $\fp$ be a non-maximal prime ideal of $A$. Show that $k[X]\to A$ is finite and injective, and deduce $\fp\cap k[X]=0$. For $g\in\fp$ with integral equation $g^n+a_{n-1}g^{n-1}+\dots+a_0=0$, $a_i\in k[X]$, show $a_0=0$. Conclude $\fp=0$.}
\pf
Write $K=A/\fm$ for a maximal ideal $\fm$; by Cor.~1.12, $\dim_\R K\le2$ (so $K\simeq\R$ or $\C$). Since $X^2+Y^2+1$ is irreducible in $\R[X,Y]$ (see (c)), $A$ is a domain, and $A\neq 0$.

\emph{(a)} Since $\dim_\R K\le2$, the elements $1,\bar x,\bar x^2$ of $K$ are linearly dependent over $\R$: there are real numbers $(\lambda_0,\lambda_1,\lambda_2)\neq0$ with $\lambda_0+\lambda_1\bar x+\lambda_2\bar x^2=0$. If $\lambda_2\ne0$ we divide by it and obtain $a,b$ with $x^2+ax+b\in\fm$. If $\lambda_2=0$ then $\bar x\in\R$ (when $\lambda_1\neq0$; and $\lambda_1=\lambda_2=0$ forces $\lambda_0=0$), say $\bar x=r$, and $x^2-rx=x(x-r)\in\fm$, which has the required form with $a=-r$, $b=0$. The same argument for $y$ gives $c,d$.

Now $1,\bar x,\bar y\in K$ are linearly dependent over $\R$ because $\dim_\R K\le2$: there is $(\alpha,\beta,\gamma)\neq(0,0,0)$ with $\alpha\bar x+\beta\bar y+\gamma=0$. If $\alpha=\beta=0$ then $\gamma=0$, a contradiction. So $f=\alpha x+\beta y+\gamma\in\fm$ with $(\alpha,\beta)\neq(0,0)$. (This is the stated claim; the relation $x^2+y^2+1=0$ would give $ax+cy+(b+d-1)\in\fm$, but this element could be a constant, so we use the linear dependence directly.)

\emph{$fA$ is maximal.} Let $f=\alpha x+\beta y+\gamma$ be any element with $(\alpha,\beta)\ne(0,0)$. By symmetry between $x$ and $y$ assume $\alpha\neq0$. In $A/fA=\R[X,Y]/(X^2+Y^2+1,\alpha X+\beta Y+\gamma)$ we can eliminate $X=-(\beta Y+\gamma)/\alpha$, so $A/fA\simeq\R[Y]/(g)$ with
\[ g=\Bigl(\tfrac{\beta^2}{\alpha^2}+1\Bigr)Y^2+\tfrac{2\beta\gamma}{\alpha^2}Y+\tfrac{\gamma^2}{\alpha^2}+1. \]
This is a polynomial of degree $2$ with discriminant
\[ \tfrac{4\beta^2\gamma^2}{\alpha^4}-4\,\tfrac{(\alpha^2+\beta^2)(\alpha^2+\gamma^2)}{\alpha^4}=-\tfrac{4(\alpha^2+\beta^2+\gamma^2)}{\alpha^2}<0, \]
so $g$ is irreducible over $\R$ and $A/fA\simeq\R[Y]/(g)\simeq\C$ is a field. Thus $fA$ is maximal. Also $1,\bar y$ form an $\R$-basis of $A/fA$.

For our $\fm$ we now have $fA\subseteq\fm$ with $fA$ maximal and $\fm\neq A$, so $\fm=fA$.

\emph{(b)} The point $(\alpha:\beta:\gamma)$ with $(\alpha,\beta)=(0,0)$ is $(0:0:1)$, and $f=1$ then generates $A$. For $(\alpha:\beta:\gamma)\in\mathbf P(k^3)\setminus\{(0:0:1)\}$ we have $(\alpha,\beta)\neq(0,0)$, and the ideal $fA$ does not change if $(\alpha,\beta,\gamma)$ is multiplied by a nonzero scalar; by (a), $fA$ is a maximal ideal. So the map is well defined, and surjective by (a).

\emph{Injectivity.} Let $f=\alpha x+\beta y+\gamma$ and $f'=\alpha'x+\beta'y+\gamma'$ with $f'A=fA$ (both $(\alpha,\beta)$, $(\alpha',\beta')$ nonzero). Consider the $\R$-linear map $\theta:\R^3\to A/fA$, $(u,v,w)\mapsto u\bar x+v\bar y+w$. Assume $\alpha\neq0$ (the case $\beta\ne0$ is symmetric). Then $\bar x\in\R+\R\bar y$ in $A/fA$ and $\{1,\bar y\}$ is a basis of $A/fA$, so $\theta$ is surjective onto a $2$-dimensional space and $\dim\Ker\theta=1$. Since $(\alpha,\beta,\gamma)\in\Ker\theta$ is nonzero, $\Ker\theta=\R(\alpha,\beta,\gamma)$. As $f'\in fA$, $(\alpha',\beta',\gamma')\in\Ker\theta$, so it equals $\lambda(\alpha,\beta,\gamma)$ for some $\lambda\neq0$. The two points of $\mathbf P(k^3)$ coincide. So the map is a bijection.

\emph{(c)} \emph{Finite and injective.} In $A$ we have $y^2+(x^2+1)=0$, so $y$ is integral over $\R[x]$ (the image of $\R[X]$) and $A=\R[x]+\R[x]y$ is generated by $1,y$ as an $\R[X]$-module: finite. Injective: if $P(X)\in\R[X]$ lies in $(X^2+Y^2+1)\R[X,Y]$, say $P=(X^2+Y^2+1)Q$ with $Q\neq0$, then the $Y$-degree of the right side is $\ge2$ while that of $P$ is $0$; so $Q=0$ and $P=0$.

\emph{$A$ is a domain.} The polynomial $Y^2+(X^2+1)\in\R[X][Y]$ is monic of degree $2$ in $Y$; it is reducible iff $-(X^2+1)$ is a square in $\R[X]$, which is impossible (the leading coefficient of a square is positive). Hence it is irreducible, so it is a prime element of the UFD $\R[X,Y]$, and $A$ is a domain.

\emph{$\fp\cap\R[X]=0$.} Let $\fq=\fp\cap\R[X]$, a prime ideal of the PID $\R[X]$. Suppose $\fq\neq0$. Then $\fq$ is maximal and $\R[X]/\fq$ is a finite extension of $\R$. The ring $A/\fp$ is generated as an $\R[X]/\fq$-module by $1,\bar y$, so $\dim_\R A/\fp<\infty$; it is a domain, hence a field (as in Ex.~1.2), and $\fp$ is maximal, contrary to assumption. So $\fq=0$.

\emph{Conclusion.} Let $g\in\fp$. Since $A$ is finite over $\R[X]$, $g$ satisfies an integral equation $g^n+a_{n-1}g^{n-1}+\dots+a_0=0$ with $a_i\in\R[X]$ (images in $A$); choose one with $n\ge1$ minimal. Then $a_0=-g(g^{n-1}+a_{n-1}g^{n-2}+\dots+a_1)\in\fp\cap\R[X]=0$. If $n=1$ then $g=-a_0=0$. If $n\ge2$, then $g\,(g^{n-1}+\dots+a_1)=0$ in the domain $A$; if $g\neq0$ this gives $g^{n-1}+a_{n-1}g^{n-2}+\dots+a_1=0$, an integral equation of smaller degree, contradicting minimality; hence $g=0$. Therefore $\fp=0$.

\emph{Summary.} $\Spec A=\{0\}\cup\{(\alpha x+\beta y+\gamma)A\}$, the nonzero points being the maximal ideals in bijection with $\mathbf P(\R^3)\setminus\{(0:0:1)\}$ as in (b); $0$ is prime since $A$ is a domain. \hfill$\square$

\exercise{1.4}{Let $A$ be a ring.
(a) Let $\fp$ be a minimal prime ideal of $A$. Show that $\fp A_\fp$ is nilpotent. Deduce that every element of $\fp$ is a zero divisor in $A$.
(b) Show that if $A$ is reduced, then any zero divisor in $A$ lies in a minimal prime ideal. Show by an example (use Lemma~1.6) that this is false if $A$ is not reduced.}
\pf
\emph{(a)} The primes of $A_\fp$ correspond to the primes of $A$ contained in $\fp$ (Lemma~1.7(c)); by minimality of $\fp$ the only one is $\fp$. So $A_\fp$ has the unique prime ideal $\fp A_\fp$, and by Lemma~1.6(a) (applied to $A_\fp$ and $I=0$) $\fp A_\fp$ is the nilradical of $A_\fp$. Thus \emph{every element} of $\fp A_\fp$ is nilpotent.

\emph{Correction.} The ideal $\fp A_\fp$ itself need not be nilpotent when $A$ is not noetherian: take $A=k[x_1,x_2,\dots]/(x_1,x_2^2,x_3^3,\dots)$. Every $x_i$ is nilpotent, so the only prime is $\fm=(x_i)_i$ (with $A/\fm=k$), $A=A_\fm$, and $\fm^N\ni x_{N+1}^N\neq0$ for all $N$. What is true is: $\fp A_\fp$ is a nil ideal, and it is nilpotent if it is finitely generated (e.g. if $A$ is noetherian): if $\fp A_\fp=(g_1,\dots,g_r)$ with $g_i^{n_i}=0$, then $(\fp A_\fp)^{N}=0$ for $N=\sum(n_i-1)+1$, since each monomial of degree $N$ in the $g_i$ contains some $g_i^{n_i}$.

\emph{Zero divisors.} Let $a\in\fp$. Then $a/1\in\fp A_\fp$ is nilpotent, so $a^n/1=0$ for some $n\ge1$, i.e. there is $s\in A\setminus\fp$ with $sa^n=0$. Choose $n\ge1$ minimal with this property (for some $s$ fixed). Then $b:=sa^{n-1}\neq0$: for $n=1$, $b=s\neq0$ because $0\in\fp$ and $s\notin\fp$; for $n\ge2$ by minimality of $n$. And $ab=sa^n=0$. So $a$ is a zero divisor.

\emph{(b)} Let $A$ be reduced and $ab=0$ with $b\neq0$ (the zero ring has no zero divisors). First, every prime $\fq$ contains a minimal prime: the set of primes contained in $\fq$, ordered by reverse inclusion, is inductive, since the intersection of a chain $(\fp_i)$ of primes is prime (if $uv\in\bigcap\fp_i$ and $u\notin\fp_i$, then $u\notin\fp_j$ for all $\fp_j\subseteq\fp_i$, so $v\in\fp_j$ for those, and $v\in\fp_i\subseteq\fp_j$ for the others); Zorn's lemma gives a minimal element. Hence the intersection of all primes equals the intersection of the minimal primes; by Lemma~1.6(a) it is $\sqrt0=0$. Since $b\neq0$, some minimal prime $\fp$ does not contain $b$. From $ab=0\in\fp$ we get $a\in\fp$.

\emph{Counterexample in the non-reduced case.} Let $A=k[X,Y]/(X^2,XY)$ and $x,y$ the images. Then $x\neq0$ (the ideal $(X^2,XY)$ consists of polynomials whose terms have degree $\ge2$), $x^2=0$, and $A/xA\simeq k[Y]$ is a domain, so $xA$ is prime. Every prime contains the nilpotent $x$, hence contains $xA$; so $xA$ is the unique minimal prime (equivalently $\sqrt 0=xA$, Lemma~1.6(a)). The element $y$ is a zero divisor ($yx=0$, $x\ne0$), but $y\notin xA$ (its image in $A/xA=k[Y]$ is $Y\neq0$). So $y$ lies in no minimal prime. \hfill$\square$

\exercise{1.5}{Let $k$ be a field and $\fm$ a maximal ideal of $k[T_1,\dots,T_n]$.
(a) Let $P_1(T_1)$ be a generator of $\fm\cap k[T_1]$ and $k_1=k[T_1]/(P_1)$. Show that there is an exact sequence $0\to P_1(T_1)k[T_1,\dots,T_n]\to k[T_1,\dots,T_n]\to k_1[T_2,\dots,T_n]\to0$.
(b) Show that there exist $n$ polynomials $P_1(T_1),P_2(T_1,T_2),\dots,P_n(T_1,\dots,T_n)$ such that $k[T_1,\dots,T_i]\cap\fm=(P_1,\dots,P_i)$ for all $i\le n$ (the ideal being taken in $k[T_1,\dots,T_i]$). In particular $\fm$ is generated by $n$ elements.}
\pf
\emph{(a)} Let $R=k[T_1,\dots,T_n]$. The ideal $\fm\cap k[T_1]$ is prime and proper ($1\notin\fm$). It is nonzero: by Cor.~1.12, $R/\fm$ is a finite, hence algebraic, extension of $k$, so the image of $T_1$ satisfies a nonzero polynomial over $k$, which is then in $\fm\cap k[T_1]$. So $P_1\neq0$ is irreducible (a nonzero prime of the PID $k[T_1]$), and $k_1$ is a field. Write $R=k[T_1][T_2,\dots,T_n]$ and let $\pi:R\to k_1[T_2,\dots,T_n]$ reduce coefficients modulo $P_1$. It is a surjective ring homomorphism, and $\Ker\pi$ is the set of polynomials all of whose coefficients lie in $P_1k[T_1]$, i.e. $\Ker\pi=P_1R$. Multiplication by $P_1$ is injective since $R$ is a domain. This is the exact sequence.

\emph{(b)} We prove by induction on $n$, for all fields $k$ simultaneously, that such $P_i$ exist. For $n=1$ take $P_1$ as in (a): $\fm=\fm\cap k[T_1]=(P_1)$. Let $n\ge2$ and take $P_1$, $k_1$, $\pi$ as in (a). Since $P_1\in\fm=\pi^{-1}(\pi(\fm))$ (because $\Ker\pi=P_1R\subseteq\fm$), the ideal $\fm'=\pi(\fm)$ is a maximal ideal of $k_1[T_2,\dots,T_n]$ (it corresponds to $\fm/P_1R$ in $R/P_1R$), and $\fm=\pi^{-1}(\fm')$. By the induction hypothesis (field $k_1$, $n-1$ variables $T_2,\dots,T_n$) there are $Q_2(T_2),\dots,Q_n(T_2,\dots,T_n)$ with $k_1[T_2,\dots,T_j]\cap\fm'=(Q_2,\dots,Q_j)$ for $2\le j\le n$. Lift each $Q_j$ to $P_j\in k[T_1,\dots,T_j]$ with $\pi(P_j)=Q_j$ (lift each coefficient in $k_1$ to $k[T_1]$). Then $P_j\in\pi^{-1}(\fm')=\fm$.

Fix $i\le n$ and put $R_i=k[T_1,\dots,T_i]$. The inclusion $(P_1,\dots,P_i)R_i\subseteq\fm\cap R_i$ is clear. Conversely let $F\in\fm\cap R_i$. Then $\pi(F)\in\fm'\cap k_1[T_2,\dots,T_i]$, hence $\pi(F)=\sum_{j=2}^iC_jQ_j$ with $C_j\in k_1[T_2,\dots,T_i]$ (for $i=1$ the sum is empty). Lift $C_j$ to $D_j\in R_i$. Then $F-\sum_jD_jP_j\in R_i$ is killed by $\pi$, so it lies in $P_1R\cap R_i=P_1R_i$ (a polynomial in $T_1,\dots,T_i$ divisible by $P_1$ in $R$ is divisible in $R_i$, by comparing coefficients of the monomials in $T_2,\dots,T_n$). Hence $F\in(P_1,\dots,P_i)R_i$.

For $i=n$ this gives $\fm=(P_1,\dots,P_n)$. \hfill$\square$

\exercise{1.6}{A topological space $X$ is quasi-compact if every open covering has a finite subcovering. Show that any closed subspace of a quasi-compact space is quasi-compact. Show that $\Spec A$ is quasi-compact.}
\pf
\emph{Closed subspaces.} Let $X$ be quasi-compact, $Z\subseteq X$ closed, and $Z\subseteq\bigcup_{i}U_i$ be a covering of $Z$ by sets $U_i\cap Z$ with $U_i$ open in $X$ (every open of $Z$ has this form). Then $\{U_i\}_i\cup\{X\setminus Z\}$ is an open covering of $X$; take a finite subcovering. Discarding $X\setminus Z$ (which does not meet $Z$) we get finitely many $U_i$ covering $Z$.

\emph{$\Spec A$.} Let $\Spec A=\bigcup_iU_i$ with $U_i$ open. By Prop.~1.1 each $U_i$ is a union of principal open sets $D(f)$, so the family $\mathcal F$ of all $D(f)$ contained in some $U_i$ covers $\Spec A$. It suffices to find finitely many members of $\mathcal F$ covering $\Spec A$: each lies in some $U_i$, and those finitely many $U_i$ then cover.

Write $\mathcal F=\{D(f_j)\}_{j\in J}$ and let $I$ be the ideal generated by the $f_j$. Then $V(I)=\bigcap_jV(f_jA)=\Spec A\setminus\bigcup_jD(f_j)=\emptyset$ (Prop.~1.1(b)). So $I$ is contained in no prime ideal, hence in no maximal ideal, so $I=A$ (otherwise $I$ lies in a maximal ideal by Zorn's lemma). Thus $1=\sum_{l=1}^ma_lf_{j_l}$ for finitely many indices. The ideal generated by $f_{j_1},\dots,f_{j_m}$ is $A$, so no prime contains all of them, i.e. $\Spec A=\bigcup_{l}D(f_{j_l})$. \hfill$\square$

\exercise{1.7}{Describe the closed subsets of $\Spec\Z[T]$. Show that they are not all principal.}
\pf
\emph{Step 1: the points of $\Spec\Z[T]$.} Let $P$ be a prime of $\Z[T]$ and $P\cap\Z=p\Z$ with $p=0$ or a prime number.

If $p\neq0$: $P/p\Z[T]$ is a prime of $\Z[T]/(p)=\mathbf F_p[T]$, so it is $0$ or $(\bar g)$ with $\bar g$ irreducible. Thus $P=(p)$ or $P=(p,g)$ with $g\in\Z[T]$ a lift of an irreducible $\bar g\in\mathbf F_p[T]$; the latter are maximal because $\mathbf F_p[T]/(\bar g)$ is a field.

If $p=0$: $P\cap(\Z\setminus0)=\emptyset$, so by Lemma~1.7(c) $P$ corresponds to the prime $P\Q[T]$ of $\Q[T]$, and $P=P\Q[T]\cap\Z[T]$. Now $P\Q[T]=0$ or $=h\Q[T]$ with $h$ irreducible; scaling, take $h\in\Z[T]$ primitive; by Gauss's lemma $h$ is irreducible in $\Z[T]$ and $h\Q[T]\cap\Z[T]=h\Z[T]$. So $P=0$ or $P=(h)$ with $h$ a primitive irreducible polynomial of degree $\ge1$.

So the points are: $0$; the primes $(h)$, $h$ primitive irreducible of positive degree; the primes $(q)$, $q$ a prime number (these two kinds are the ``height one'' points); and the maximal ideals $(q,g)$.

\emph{Step 2: irreducible closed sets.} For a point $P$ the closure of $\{P\}$ is $V(P)$. We have $V(0)=\Spec\Z[T]$. For $P=(h)$ or $P=(q)$ we claim that $V(P)=\{P\}\cup\{\text{maximal ideals containing }P\}$. Indeed let $P'\supsetneq P$ be prime. If $P=(h)$: $P'\cap\Z=0$ would force $P'=(h')$ (it is nonzero since it contains $h$), and $(h)\subseteq(h')$ gives $h'\mid h$, so $h'$ is an associate of $h$ (as $h$ is irreducible), contradicting $P'\ne P$; so $P'\cap\Z=q\Z$ with $q$ prime, and $P'=(q)$ or $(q,g)$; but $(q)\supseteq(h)$ would give $h=qu$ with $h$ primitive, so $q$ a unit, impossible; so $P'=(q,g)$ is maximal. If $P=(q)$: $P'\ni q$ so $P'\cap\Z=q\Z$ and $P'=(q)$ or $(q,g)$, i.e. $P'$ is maximal. For $P$ maximal, $V(P)=\{P\}$.

\emph{Step 3: all closed sets.} $\Z[T]$ is noetherian (Hilbert basis), so every radical ideal is a finite intersection of primes $P_1\cap\dots\cap P_r$; by Lemma~1.6(b) and Prop.~1.1(a), $V(I)=V(\sqrt I)=V(P_1\cap\dots\cap P_r)=V(P_1)\cup\dots\cup V(P_r)$ (if $\sqrt I=A$ then $V(I)=\emptyset$, the empty union). Conversely a finite union of closed sets is closed. Hence:

\emph{The closed subsets of $\Spec\Z[T]$ are exactly the finite unions of sets $V(P)$ with $P$ prime.} Concretely: $\Spec\Z[T]$ itself, and the proper closed subsets, which are the sets
\[ V(h_1)\cup\dots\cup V(h_r)\cup V(q_1)\cup\dots\cup V(q_s)\cup\{\fm_1,\dots,\fm_t\} \]
with $r,s,t\ge0$ finite, $h_i$ primitive irreducible nonconstant polynomials, $q_j$ prime numbers and $\fm_l$ maximal ideals. Here $V(h)$ consists of the point $(h)$ and the maximal ideals $(q,g)$ with $h\in(q,g)$, and $V(q)$ of $(q)$ and all $(q,g)$. In words: a proper closed subset is a finite set of closed points together with finitely many ``curves'' (horizontal $V(q)$, vertical $V(h)$, each with its generic point).

\emph{Not all closed sets are principal.} Consider $C=\{\fm\}$ with $\fm=(2,T)$; it is closed, $C=V(2,T)$ (indeed $\fm$ is maximal, Remark~1.3). Suppose $C=V(f)$. Then $f\neq0$ (else $V(f)=\Spec\Z[T]$), and $f$ is not a unit (else $V(f)=\emptyset$). Take a prime element $\pi$ dividing $f$ in the UFD $\Z[T]$. Then $(\pi)$ is a prime containing $f$, so $(\pi)\in V(f)$, but $(\pi)$ is not maximal: if $\pi$ is a constant prime number $q$, then $(q)\subsetneq(q,T)$; if $\deg\pi\ge1$, choose a prime number $q$ not dividing the leading coefficient of $\pi$; then the image of $\pi$ in $\mathbf F_q[T]$ has the same positive degree, so $(\pi,q)\ne\Z[T]$ (as $\Z[T]/(\pi,q)\simeq\mathbf F_q[T]/(\bar\pi)\neq0$), and $q\notin(\pi)$ (a nonzero multiple of $\pi$ has positive degree); hence $(\pi)\subsetneq(\pi,q)$. Thus $V(f)$ contains a non-closed point, so $V(f)\neq C$. Hence $C$ is a closed subset that is not principal. \hfill$\square$

\exercise{1.8}{Let $\varphi:A\to B$ be an integral homomorphism.
(a) Show that $\Spec\varphi:\Spec B\to\Spec A$ maps a closed point to a closed point, and that any preimage of a closed point is a closed point.
(b) Let $\fp\in\Spec A$. Show that the canonical homomorphism $A_\fp\to B\otimes_AA_\fp$ is integral.
(c) Let $T=\varphi(A\setminus\fp)$. Suppose $\varphi$ is injective. Show that $T$ is a multiplicative subset of $B$, and $B\otimes_AA_\fp=T^{-1}B\neq0$. Deduce that $\Spec\varphi$ is surjective if $\varphi$ is integral and injective.}
\pf
\emph{(a)} Let $\fq\in\Spec B$ and $\fp=\varphi^{-1}(\fq)=(\Spec\varphi)(\fq)$. The map $\varphi$ induces an injective homomorphism of domains $\bar\varphi:D:=A/\fp\hookrightarrow E:=B/\fq$, and $E$ is integral over $D$ (reduce an integral equation modulo $\fq$). We show that $D$ is a field iff $E$ is a field; by Remark~1.3 this is exactly the claim (the image of a closed point is closed, and a point whose image is closed is closed).

($E$ field $\Rightarrow D$ field.) Let $0\neq a\in D$. Then $a^{-1}\in E$ is integral over $D$: $a^{-n}+c_{n-1}a^{-(n-1)}+\dots+c_0=0$ with $c_i\in D$. Multiplying by $a^{n-1}$: $a^{-1}=-(c_{n-1}+c_{n-2}a+\dots+c_0a^{n-1})\in D$.

($D$ field $\Rightarrow E$ field.) Let $0\neq b\in E$ and take an integral equation $b^n+c_{n-1}b^{n-1}+\dots+c_0=0$ over $D$ of minimal degree $n\ge1$. Then $c_0\neq0$: otherwise $b(b^{n-1}+\dots+c_1)=0$ gives (as $E$ is a domain) either $n=1$ and $b=0$, or an equation of degree $n-1$. So $b\cdot(b^{n-1}+\dots+c_1)=-c_0$ with $c_0$ invertible in $D$, and $b$ is invertible in $E$.

\emph{(b)} Put $S=A\setminus\fp$ and $C=B\otimes_AA_\fp$, with the ring homomorphisms $\beta:B\to C$, $b\mapsto b\otimes1$, and $\iota:A_\fp\to C$, $a/s\mapsto1\otimes a/s$. Every element of $C$ is a finite sum of elements $b\otimes\frac as=\beta(b)\,\iota(a/s)$. For $b\in B$ choose $b^n+\sum_{i<n}\varphi(a_i)b^i=0$; applying $\beta$ gives $\beta(b)^n+\sum_{i<n}\iota(a_i/1)\beta(b)^i=0$ (since $\beta\varphi(a)=a\cdot(1\otimes 1)=\iota(a/1)$). So $\beta(b)$ is integral over $\iota(A_\fp)$, and $\iota(a/s)$ trivially is. The integral elements form a subring, so every element of $C$ is integral over $A_\fp$.

\emph{(c)} $T$ is multiplicative: $1=\varphi(1)\in T$, and $\varphi(s)\varphi(t)=\varphi(st)$ with $st\in A\setminus\fp$ since $\fp$ is prime. Moreover $0\notin T$: $\varphi(s)=0$ with $\varphi$ injective gives $s=0\in\fp$. The canonical map $B\otimes_AS^{-1}A\to T^{-1}B$, $b\otimes a/s\mapsto\varphi(a)b/\varphi(s)$, is the standard isomorphism $B\otimes_AS^{-1}A\simeq T^{-1}B$ (localization of the $A$-module $B$ at $S$, which is the localization of the ring $B$ at $T$). Since $0\notin T$, $1/1\neq0$ in $T^{-1}B$ (a relation $t\cdot1=0$ with $t\in T$ would give $0\in T$). So $C=T^{-1}B\neq0$.

\emph{Surjectivity.} Let $\fp\in\Spec A$. As $C\ne0$, $C$ has a maximal ideal $\fn$. By (b), $\iota:A_\fp\to C$ is integral, so by (a) the prime $\iota^{-1}(\fn)$ is a maximal ideal of the local ring $A_\fp$, i.e. $\iota^{-1}(\fn)=\fp A_\fp$. Let $\fq=\beta^{-1}(\fn)\in\Spec B$. Since $\beta\circ\varphi=\iota\circ(A\to A_\fp)$, we get
\[ \varphi^{-1}(\fq)=\varphi^{-1}\beta^{-1}(\fn)=(A\to A_\fp)^{-1}\iota^{-1}(\fn)=(A\to A_\fp)^{-1}(\fp A_\fp)=\fp. \]
So $\fp$ is in the image of $\Spec\varphi$. \hfill$\square$

\exercise{1.9}{Let $A$ be a finitely generated algebra over a field $k$.
(a) Suppose $A$ is finite over $k$. Show that $\Spec A$ is a finite set of cardinality at most $\dim_kA$ (construct a strictly descending chain of ideals by taking intersections of maximal ideals). Show that every prime ideal of $A$ is maximal.
(b) Show that $\Spec k[T_1,\dots,T_d]$ is infinite if $d\ge1$.
(c) Show that $\Spec A$ is finite iff $A$ is finite over $k$.}
\pf
\emph{(a)} Let $n=\dim_kA<\infty$. For a prime $\fp$, $A/\fp$ is a domain of finite dimension over $k$, hence a field (as in Ex.~1.2), so $\fp$ is maximal: $\Spec A=\Max A$.

Let $\fm_1,\dots,\fm_r$ be distinct maximal ideals and $I_0=A$, $I_i=\fm_1\cap\dots\cap\fm_i$. We claim $I_{i}\supsetneq I_{i+1}$ for $0\le i<r$. Otherwise $I_i\subseteq\fm_{i+1}$; for $i=0$ this contradicts $\fm_1\ne A$. For $i\ge1$, $\fm_1\cdots\fm_i\subseteq I_i\subseteq\fm_{i+1}$, and as $\fm_{i+1}$ is prime, some $\fm_j\subseteq\fm_{i+1}$ with $j\le i$, so $\fm_j=\fm_{i+1}$ by maximality, contradicting distinctness. The $I_i$ are $k$-subspaces, so $n=\dim I_0>\dim I_1>\dots>\dim I_r\ge0$ forces $r\le n$. Hence $\#\Spec A\le n$.

\emph{(b)} Let $d\ge1$ and let $P\in k[T_1]$ be monic irreducible. The ideal $\fm_P=(P(T_1),T_2,\dots,T_d)$ has quotient $k[T_1]/(P)$, a field, so $\fm_P$ is maximal. Moreover $\fm_P\cap k[T_1]=(P)$: the composite $k[T_1]\to k[T_1,\dots,T_d]\to k[T_1,\dots,T_d]/\fm_P\simeq k[T_1]/(P)$ is the canonical surjection, with kernel $(P)$. Hence distinct monic irreducible $P$ give distinct $\fm_P$. There are infinitely many monic irreducible polynomials in $k[T_1]$ for any field $k$ (Euclid): given a finite list $P_1,\dots,P_r$, the polynomial $P_1\cdots P_r+1$ has degree $\ge1$, hence a monic irreducible factor, which is none of the $P_i$. So $\Spec k[T_1,\dots,T_d]$ is infinite.

\emph{(c)} ($\Leftarrow$) is (a). ($\Rightarrow$) Suppose $\Spec A$ finite. If $A=0$ there is nothing to show, so let $A\neq0$. Let $\fm_1,\dots,\fm_r$ be the maximal ideals (finitely many). By Lemma~1.18, $\Nil:=\sqrt0=\bigcap_i\fm_i$. Each $A/\fm_i$ is finite over $k$ by Cor.~1.12, and the natural map $A/\Nil\to\prod_iA/\fm_i$ is injective, so $\dim_kA/\Nil<\infty$.

Now $A$ is noetherian (Hilbert basis), so $\Nil$ is finitely generated, say by $g_1,\dots,g_s$ with $g_j^{n_j}=0$; as in Ex.~1.4(a), $\Nil^N=0$ for $N=\sum(n_j-1)+1$. Consider the filtration $A\supseteq\Nil\supseteq\Nil^2\supseteq\dots\supseteq\Nil^N=0$. Each $\Nil^i$ is a finitely generated $A$-module (noetherian), hence $\Nil^i/\Nil^{i+1}$ is a finitely generated module over $A/\Nil$, which is finite dimensional over $k$; therefore $\dim_k\Nil^i/\Nil^{i+1}<\infty$. Summing, $\dim_kA=\dim_kA/\Nil+\sum_{i=1}^{N-1}\dim_k\Nil^i/\Nil^{i+1}<\infty$. So $A$ is finite over $k$. \hfill$\square$

\section*{2.2 Ringed topological spaces}

Throughout, sheaves and presheaves are of abelian groups on a topological space $X$, and results of the text are cited by their numbers in Liu (``Def.~2.2'', ``Prop.~2.15'', \dots). A morphism of presheaves is injective if all $\alpha(U)$ are, and surjective if all stalk maps $\alpha_x$ are (Def.~2.10).

\exercise{2.1}{Determine the sheaf associated to the presheaf $A_X$ of Example 2.4 (a non-trivial abelian group $A$; $A_X(U)=A$ for $U\neq\emptyset$, $A_X(\emptyset)=0$, restrictions $\operatorname{Id}_A$ between non-empty open sets). Such a sheaf is called a constant sheaf. Show that $A_X$ is a sheaf if and only if every non-empty open subset of $X$ is connected.}
\pf
\emph{The stalks.} For $x\in X$ the direct limit of $A_X(U)=A$ over the neighbourhoods $U\ni x$ with identity transition maps is $A$ (all $U\ni x$ are non-empty); we write $(A_X)_x=A$, and the germ at $x$ of $a\in A_X(U)$ is $a$.

\emph{The associated sheaf.} By the construction in Prop.~2.15, $A_X^\dagger(U)$ is the group of functions $f:U\to\bigsqcup_{x\in U}(A_X)_x=U\times A$, i.e.\ functions $f:U\to A$, such that every $x\in U$ has an open neighbourhood $V\subseteq U$ and a section $a\in A_X(V)=A$ with $f(y)=a$ for all $y\in V$. These are exactly the locally constant functions $U\to A$ (equivalently the continuous maps $U\to A$ when $A$ has the discrete topology). Restrictions are restrictions of functions, and $\theta(U):A_X(U)\to A_X^\dagger(U)$ sends $a$ to the constant function with value $a$ (for $U\neq\emptyset$). So \emph{the sheaf associated to $A_X$ is the sheaf $U\mapsto\{\text{locally constant functions }U\to A\}$}. (It is a sheaf, as Prop.~2.15 asserts: being locally constant and equality of functions are local conditions, so uniqueness and glueing of Def.~2.2 hold for it as a subpresheaf of the sheaf of all functions into $A$.) Its stalk at every $x$ is $A$ and $\theta_x=\operatorname{Id}_A$.

\emph{Criterion.} Let $U\neq\emptyset$ be open. The map $\theta(U)$ is injective, because a constant function with value $a$ on $U\ne\emptyset$ is zero only if $a=0$. We claim $\theta(U)$ is surjective if and only if $U$ is connected.

If $U$ is connected and $f:U\to A$ is locally constant, then for each $a\in A$ the set $f^{-1}(a)$ is open (local constancy), and $U$ is the disjoint union of these sets. Since $U\neq\emptyset$ is connected, only one of them is non-empty, so $f$ is constant, $f=\theta(U)(a)$. If $U$ is not connected, write $U=U_1\sqcup U_2$ with $U_1,U_2$ non-empty, open and disjoint. Choose $a\neq0$ in $A$ and let $f=a$ on $U_1$, $f=0$ on $U_2$. Then $f$ is locally constant but not constant, so it is not in the image of $\theta(U)$.

For $U=\emptyset$ both groups are $0$. Hence $\theta:A_X\to A_X^\dagger$ is an isomorphism of presheaves iff every non-empty open set is connected.

\emph{Conclusion.} If every non-empty open set is connected, $A_X\simeq A_X^\dagger$ is a sheaf. Conversely, if $A_X$ is a sheaf, then $\theta:A_X\to A_X^\dagger$ is a morphism of sheaves which is an isomorphism on all stalks (Prop.~2.15), hence an isomorphism by Prop.~2.12; so every non-empty open set is connected by the criterion. \hfill$\square$

\exercise{2.2}{Let $\mathcal F$ be a sheaf on $X$ and $s,t\in\mathcal F(X)$. Show that the set of $x\in X$ such that $s_x=t_x$ is open in $X$.}
\pf
Put $u=s-t\in\mathcal F(X)$. Since taking germs is a group homomorphism, $s_x=t_x$ if and only if $u_x=0$. Let $x$ be such that $u_x=0$. By the definition of the direct limit $\mathcal F_x=\varinjlim_{U\ni x}\mathcal F(U)$, the germ of $u$ at $x$ is zero iff there is an open neighbourhood $V$ of $x$ with $u|_V=0$. Then for every $y\in V$ we have $u_y=(u|_V)_y=0$. Thus $V$ is contained in the set $\{x: s_x=t_x\}$, which is therefore a neighbourhood of each of its points, i.e.\ open. (Only the presheaf property was used.) \hfill$\square$

\exercise{2.3}{(Sheaf associated to a presheaf) Fix a topological space $X$ and a presheaf $\mathcal F$ on $X$; keep the notation before Lemma 2.7. (a) Let $\mathcal U=\{U_i\}_i$ be an open covering of $X$ and $\mathcal F_{\mathcal U}(X)=\Ker d^1$. For any open $W\subseteq X$ define in the same way a group $\mathcal F_{\mathcal U}(W)$ relative to the covering $\{W\cap U_i\}_i$ of $W$. Show that $\mathcal F_{\mathcal U}$ is a presheaf and that there is a morphism of presheaves $\mathcal F\to\mathcal F_{\mathcal U}$. (b) Let $\mathcal V=\{V_k\}_k$ be another open covering of $X$; write $\mathcal V\preceq\mathcal U$ if every $V_k$ is contained in some $U_i$. Show that there is a canonical homomorphism $\mathcal F_{\mathcal U}(W)\to\mathcal F_{\mathcal V}(W)$ for every open $W$ if $\mathcal V\preceq\mathcal U$. (c) Show that $\preceq$ makes the set of open coverings of $X$ into a direct system. Let $\mathcal F'(W)=\varinjlim_{\mathcal V}\mathcal F_{\mathcal V}(W)$, the limit over the open coverings of $W$. Show that $\mathcal F'$ is the sheaf associated to $\mathcal F$.}
\pf
\textbf{Honest note.} Parts (a), (b) are proved as stated. Part (c) is \emph{false as stated for an arbitrary presheaf}: $\mathcal F'$ (the ``plus construction'') is always a separated presheaf, and it is a sheaf, hence the sheaf associated to $\mathcal F$, \emph{if $\mathcal F$ satisfies the uniqueness axiom (4)}; for general $\mathcal F$ one must apply the construction twice. A counterexample is given at the end. I prove: (c1) the direct system; (c2) $\theta_x:\mathcal F_x\to\mathcal F'_x$ is an isomorphism for all $x$; (c3) $\mathcal F'$ is separated for every $\mathcal F$, and is a sheaf if $\mathcal F$ is separated; (c4) the corrected statement that $\mathcal F'$ is the associated sheaf iff it is a sheaf (in particular if $\mathcal F$ is separated), and $(\mathcal F')'$ is always the associated sheaf.

\textbf{Conventions.} An open covering of an open set $W$ is a \emph{set} $\mathcal U$ of open subsets of $W$ with union $W$ (so the coverings of $W$ form a set, and we index a covering by its members). For such $\mathcal U$ put
\[ C_W(\mathcal U)=\Bigl\{(s_A)_{A\in\mathcal U}\in\prod_{A\in\mathcal U}\mathcal F(A)\ :\ s_A|_{A\cap B}=s_B|_{A\cap B}\ \text{for all }A,B\in\mathcal U\Bigr\}=\Ker d^1 .\]
For a covering $\mathcal U$ of $X$ and open $W\subseteq X$, let $\mathcal U|_W=\{W\cap A:A\in\mathcal U\}$, a covering of $W$; then $\mathcal F_{\mathcal U}(W)$ of the exercise is $C_W(\mathcal U|_W)$. (Using indexed families $\{U_i\}_i$ as in the book changes nothing below, since only the sets $U_i$ and the compatibility of the $s_i$ matter.)

\textbf{(a)} For $W'\subseteq W$ and $s=(s_B)_{B\in\mathcal U|_W}\in\mathcal F_{\mathcal U}(W)$ define $\rho(s)\in\prod_{B'\in\mathcal U|_{W'}}\mathcal F(B')$ as follows: a member of $\mathcal U|_{W'}$ has the form $B'=W'\cap A$ with $A\in\mathcal U$, and we put $\rho(s)_{B'}=s_{W\cap A}|_{B'}$. This is well defined: if also $B'=W'\cap A_1$, then $B'\subseteq W\cap A\cap A_1$, and $s_{W\cap A}$, $s_{W\cap A_1}$ agree on $W\cap A\cap A_1$ by compatibility. The family $\rho(s)$ is compatible because restriction is transitive. Functoriality $\rho_{W''W'}\rho_{W'W}=\rho_{W''W}$ and $\rho_{WW}=\operatorname{Id}$ is clear from transitivity of restriction in $\mathcal F$. Moreover $\mathcal F_{\mathcal U}(\emptyset)=0$ because $\mathcal F(\emptyset)=0$. So $\mathcal F_{\mathcal U}$ is a presheaf. The map $\theta_{\mathcal U}(W):\mathcal F(W)\to\mathcal F_{\mathcal U}(W)$, $s\mapsto(s|_{W\cap A})_A$, lands in $\Ker d^1$ since $(s|_{W\cap A})|_{W\cap A\cap B}=s|_{W\cap A\cap B}$ is symmetric in $A,B$, i.e.\ $d^1d^0=0$; it is a group homomorphism and commutes with restrictions by transitivity. This is the morphism of presheaves $\mathcal F\to\mathcal F_{\mathcal U}$.

\textbf{(b)} Work with coverings $\mathcal U,\mathcal V$ of an open set $W$ (apply to $\mathcal U|_W$, $\mathcal V|_W$ for the statement as given; $\mathcal V\preceq\mathcal U$ implies $\mathcal V|_W\preceq\mathcal U|_W$). Let $\mathcal V\preceq\mathcal U$ and choose a map $\tau:\mathcal V\to\mathcal U$ with $V\subseteq\tau(V)$ for all $V\in\mathcal V$. Define
\[ c_\tau:C_W(\mathcal U)\to C_W(\mathcal V),\qquad (s_A)_A\mapsto(s_{\tau(V)}|_V)_{V\in\mathcal V}. \]
It is a homomorphism, and its image is compatible: for $V,V'\in\mathcal V$ both $s_{\tau(V)}|_{V\cap V'}$ and $s_{\tau(V')}|_{V\cap V'}$ are restrictions to $V\cap V'\subseteq\tau(V)\cap\tau(V')$ of two elements which agree on $\tau(V)\cap\tau(V')$. \emph{Independence of $\tau$:} if $\tau'$ is another choice, then $V\subseteq\tau(V)\cap\tau'(V)$ and $s_{\tau(V)}|_V=s_{\tau'(V)}|_V$ by the compatibility of $(s_A)$. So $c_{\mathcal V\mathcal U}:=c_\tau$ is canonical. It is compatible with restriction to a smaller open set $W'$ (both ways around the square give $(s_{\tau(V)}|_{W'\cap V})_V$, by the independence just proved), and with the maps $\theta$ of (a): $c(\theta(s))=(s|_{V})_V=\theta_{\mathcal V}(s)$. Also $c_{\mathcal U\mathcal U}=\operatorname{Id}$ (take $\tau=\operatorname{Id}$) and, for $\mathcal V'\preceq\mathcal V\preceq\mathcal U$, $c_{\mathcal V'\mathcal V}c_{\mathcal V\mathcal U}=c_{\mathcal V'\mathcal U}$ (compose the choice functions and use independence).

\textbf{(c1) The direct system.} The relation $\preceq$ is reflexive and transitive (a preorder), and any two coverings $\mathcal U,\mathcal V$ of $W$ have the common refinement $\{A\cap B:A\in\mathcal U,B\in\mathcal V\}\preceq\mathcal U,\mathcal V$. So the coverings of $W$ form a directed (pre)ordered set and, by (b), $\mathcal U\mapsto C_W(\mathcal U)$ with the maps $c_{\mathcal V\mathcal U}$ is a direct system. Put $\mathcal F'(W)=\varinjlim_{\mathcal U}C_W(\mathcal U)$ (this is the limit meant in (c)). Write $[\mathcal U;(s_A)]$ for the class of $(s_A)\in C_W(\mathcal U)$. Two elements $(s)\in C_W(\mathcal U)$, $(s')\in C_W(\mathcal U')$ define the same class iff they have the same image in $C_W(\mathcal V)$ for some common refinement $\mathcal V$; in particular $[\mathcal U;(s)]=0$ iff $c_{\mathcal V\mathcal U}(s)=0$ for some $\mathcal V\preceq\mathcal U$. For $W'\subseteq W$ the maps $\rho:C_W(\mathcal U)\to C_{W'}(\mathcal U|_{W'})$, $(s_A)\mapsto(s_A|_{W'\cap A})$ commute with the $c$'s, hence induce $\mathcal F'(W)\to\mathcal F'(W')$; these are functorial, and $\mathcal F'(\emptyset)=0$. So $\mathcal F'$ is a presheaf, and $\theta(W):\mathcal F(W)\to\mathcal F'(W)$, $s\mapsto[\{W\};(s)]$, is a morphism of presheaves. Note: for $t=[\mathcal U;(s_A)]\in\mathcal F'(W)$ and $A_0\in\mathcal U$ we have
\[ t|_{A_0}=\theta(s_{A_0})\ \text{ in }\mathcal F'(A_0)\qquad(*) \]
Indeed $t|_{A_0}=[\mathcal U|_{A_0};(s_A|_{A_0\cap A})_A]$, while $\{A_0\cap A\}_A\preceq\{A_0\}$ and the image of $s_{A_0}$ in $C_{A_0}(\mathcal U|_{A_0})$ is $(s_{A_0}|_{A_0\cap A})_A$, which equals $(s_A|_{A_0\cap A})_A$ by compatibility.

\textbf{(c2) $\theta_x:\mathcal F_x\to\mathcal F'_x$ is an isomorphism for every $x$.} \emph{Surjective:} a germ in $\mathcal F'_x$ is $t_x$ with $t=[\mathcal U;(s_A)]\in\mathcal F'(W)$, $x\in W$. Pick $A_0\in\mathcal U$ containing $x$. By $(*)$, $t_x=(t|_{A_0})_x=\theta(s_{A_0})_x=\theta_x((s_{A_0})_x)$. \emph{Injective:} let $\sigma=s_x$ with $s\in\mathcal F(V)$ and suppose $\theta_x(\sigma)=\theta(s)_x=0$. Then there is an open $V'\ni x$, $V'\subseteq V$, with $\theta(s|_{V'})=\theta(s)|_{V'}=0$ in $\mathcal F'(V')$, i.e.\ some covering $\mathcal V\preceq\{V'\}$ of $V'$ satisfies $s|_{B}=0$ for all $B\in\mathcal V$. Choose $B\in\mathcal V$ containing $x$: then $\sigma=(s|_B)_x=0$.

\textbf{(c3) $\mathcal F'$ is separated; it is a sheaf if $\mathcal F$ is separated.}

\emph{Separated.} Let $W=\bigcup_jW_j$ be an open covering and $t=[\mathcal U;(s_A)]\in\mathcal F'(W)$ with $t|_{W_j}=0$ for all $j$. Then for each $j$ there is a covering $\mathcal V_j\preceq\mathcal U|_{W_j}$ of $W_j$ whose image of $(s_A|_{W_j\cap A})$ is $0$: for every $B\in\mathcal V_j$ and any $A\in\mathcal U$ with $B\subseteq A$ we get $s_A|_B=0$. Let $\mathcal V=\bigcup_j\mathcal V_j$, a covering of $W$ refining $\mathcal U$. For $B\in\mathcal V$ and $A\in\mathcal U$ with $B\subseteq A$ we have $s_A|_B=0$ (choose $j$ with $B\in\mathcal V_j$; any choice of $A$ gives the same restriction to $B$ by compatibility). Hence $c_{\mathcal V\mathcal U}(s)=0$ and $t=0$. This is exactly axiom (4) for $\mathcal F'$.

\emph{Glueing when $\mathcal F$ is separated.} Let $W=\bigcup_jW_j$ and $t_j=[\mathcal U_j;(s_{j,A})_{A\in\mathcal U_j}]\in\mathcal F'(W_j)$ with $t_j|_{W_{jl}}=t_l|_{W_{jl}}$, $W_{jl}=W_j\cap W_l$.

\emph{Claim.} For $j,l$, $A\in\mathcal U_j$, $B\in\mathcal U_l$: $s_{j,A}|_{A\cap B}=s_{l,B}|_{A\cap B}$. Indeed, $t_j|_{W_{jl}}=t_l|_{W_{jl}}$ gives a covering $\mathcal Z$ of $W_{jl}$, refining both $\mathcal U_j|_{W_{jl}}$ and $\mathcal U_l|_{W_{jl}}$, such that for every $Z\in\mathcal Z$ and all $A'\in\mathcal U_j$, $B'\in\mathcal U_l$ with $Z\subseteq A'$, $Z\subseteq B'$ we have $s_{j,A'}|_Z=s_{l,B'}|_Z$. Now the sets $Z\cap A\cap B$ ($Z\in\mathcal Z$) cover $A\cap B\subseteq W_{jl}$; for each of them, by compatibility of $(s_{j,\cdot})$ and $(s_{l,\cdot})$, and using $A',B'$ with $Z\subseteq A',B'$,
\[ s_{j,A}|_{Z\cap A\cap B}=s_{j,A'}|_{Z\cap A\cap B}=s_{l,B'}|_{Z\cap A\cap B}=s_{l,B}|_{Z\cap A\cap B}. \]
Since $\mathcal F$ is separated, $s_{j,A}|_{A\cap B}=s_{l,B}|_{A\cap B}$. This proves the claim.

Now let $\mathcal U=\bigcup_j\mathcal U_j$, a covering of $W$; for each $A\in\mathcal U$ choose $j(A)$ with $A\in\mathcal U_{j(A)}$ and put $\sigma_A=s_{j(A),A}$. By the claim $(\sigma_A)\in C_W(\mathcal U)$; let $t=[\mathcal U;(\sigma_A)]$. To see $t|_{W_j}=t_j$: let $\mathcal Y=\{A\cap A':A\in\mathcal U,A'\in\mathcal U_j\}$, a covering of $W_j$ refining $\mathcal U|_{W_j}$ and $\mathcal U_j$. The image of $t|_{W_j}$ in $C_{W_j}(\mathcal Y)$ is $(\sigma_A|_Y)$ and that of $t_j$ is $(s_{j,A'}|_Y)$ (for $Y\subseteq A\cap A'$); they agree by the claim ($\sigma_A=s_{j(A),A}$, applied to $j(A)$ and $j$). So $t|_{W_j}=t_j$. Hence $\mathcal F'$ satisfies (5), and is a sheaf.

\textbf{(c4) Conclusion.} Suppose $\mathcal F'$ is a sheaf. Let $(\mathcal F^\dagger,\theta^\dagger)$ be the associated sheaf (Prop.~2.15). By the universal property there is a morphism $\psi:\mathcal F^\dagger\to\mathcal F'$ with $\psi\theta^\dagger=\theta$. On stalks $\theta_x=\psi_x\theta^\dagger_x$ and $\theta_x$ (by (c2)) and $\theta^\dagger_x$ (Prop.~2.15) are isomorphisms, so each $\psi_x$ is an isomorphism, and $\psi$ is an isomorphism by Prop.~2.12. Thus $(\mathcal F',\theta)$ is the sheaf associated to $\mathcal F$. This applies when $\mathcal F$ is separated, by (c3). Conversely, if $\mathcal F'$ is the associated sheaf, it is a sheaf. For arbitrary $\mathcal F$, $\mathcal F'$ is separated by (c3), so $(\mathcal F')'$ is a sheaf by (c3) again, and $\mathcal F\to\mathcal F'\to(\mathcal F')'$ induces isomorphisms on all stalks by (c2) applied twice; the argument above, applied to the morphism $\mathcal F\to(\mathcal F')'$ (a composite with stalk isomorphisms into a sheaf), shows that $(\mathcal F')'$ is the sheaf associated to $\mathcal F$.

\textbf{Counterexample to (c) for non-separated $\mathcal F$.} Let $X=\{a,b,p,q\}$ with the open sets $\emptyset,\{a\},\{b\},W=\{a,b\},U_p=\{a,b,p\},U_q=\{a,b,q\},X$ (closed under unions and intersections). Define $\mathcal F(W)=\mathcal F(U_p)=\mathcal F(U_q)=\Z$, $\mathcal F(\emptyset)=\mathcal F(\{a\})=\mathcal F(\{b\})=\mathcal F(X)=0$, restrictions among the three copies of $\Z$ the identity, all others zero (a presheaf, not separated, since $\mathcal F(W)\to\mathcal F(\{a\})\times\mathcal F(\{b\})$ is not injective). Every covering of $U_p$ contains $U_p$, so $\{U_p\}$ is equivalent to every covering of $U_p$ and $\mathcal F'(U_p)=\Z$; likewise $\mathcal F'(U_q)=\Z$. The covering $\{\{a\},\{b\}\}$ of $W$ refines every covering of $W$, so $\mathcal F'(W)=\mathcal F(\{a\})\times\mathcal F(\{b\})=0$. Every covering of $X$ contains $X$ or both $U_p,U_q$, and $\{U_p,U_q\}$ refines every covering of $X$; so $\mathcal F'(X)=\{(m,n)\in\Z^2: m=n\}\simeq\Z$, mapping onto $\mathcal F'(U_p)$ and $\mathcal F'(U_q)$ by the diagonal. The sections $1\in\mathcal F'(U_p)$ and $2\in\mathcal F'(U_q)$ agree on $W$ (both restrict to $0\in\mathcal F'(W)$) but do not glue, so $\mathcal F'$ is not a sheaf. (The associated sheaf has $\Z\times\Z$ as sections over $X$.) \hfill$\square$

\exercise{2.4}{Let $0\to\mathcal F'\to\mathcal F\to\mathcal F''$ be an exact sequence of sheaves on $X$. Show that $0\to\mathcal F'(X)\to\mathcal F(X)\to\mathcal F''(X)$ is an exact sequence of abelian groups.}
\pf
Call the maps $\alpha:\mathcal F'\to\mathcal F$ and $\beta:\mathcal F\to\mathcal F''$. By Prop.~2.18 the sequence of stalks $0\to\mathcal F'_x\to\mathcal F_x\to\mathcal F''_x$ is exact for every $x\in X$, i.e.\ $\alpha_x$ is injective and $\Ker\beta_x=\Imm\alpha_x$.

\emph{Exact at $\mathcal F'(X)$.} If $s\in\mathcal F'(X)$ and $\alpha(X)(s)=0$, then $\alpha_x(s_x)=(\alpha(X)(s))_x=0$, so $s_x=0$ for all $x$, and $s=0$ by Lemma~2.9.

\emph{$\beta\alpha=0$ on $X$.} For $s\in\mathcal F'(X)$ the germs of $\beta(X)\alpha(X)(s)$ are $\beta_x\alpha_x(s_x)=0$ (since $\Imm\alpha_x\subseteq\Ker\beta_x$), so it vanishes by Lemma~2.9.

\emph{$\Ker\beta(X)\subseteq\Imm\alpha(X)$.} Let $u\in\mathcal F(X)$ with $\beta(X)(u)=0$. For each $x\in X$, $\beta_x(u_x)=0$, so $u_x=\alpha_x(\sigma)$ for some $\sigma\in\mathcal F'_x$. Represent $\sigma=(s_x)_x$ by a section $s_x\in\mathcal F'(U_x)$ on an open neighbourhood $U_x$ of $x$. The germs of $\alpha(U_x)(s_x)$ and $u|_{U_x}$ at $x$ coincide, hence after shrinking $U_x$ (and replacing $s_x$ by its restriction) we may assume $\alpha(U_x)(s_x)=u|_{U_x}$. On $U_x\cap U_y$,
\[ \alpha(U_x\cap U_y)\bigl(s_x|_{U_x\cap U_y}-s_y|_{U_x\cap U_y}\bigr)=u|_{U_x\cap U_y}-u|_{U_x\cap U_y}=0. \]
By injectivity of $\alpha$ on stalks, hence on every open set (the same argument as in the first step with Lemma~2.9 on $U_x\cap U_y$), $s_x$ and $s_y$ agree on $U_x\cap U_y$. As $\mathcal F'$ is a sheaf they glue to $s\in\mathcal F'(X)$ with $s|_{U_x}=s_x$. Then $\alpha(X)(s)$ and $u$ agree on every $U_x$, hence are equal by (4). So $u\in\Imm\alpha(X)$. \hfill$\square$

\exercise{2.5}{(Supports of sheaves) Let $\mathcal F$ be a sheaf on $X$ and $\Supp\mathcal F=\{x\in X\mid\mathcal F_x\neq0\}$. We want to show that in general $\Supp\mathcal F$ is not a closed subset. Fix a sheaf $\mathcal G$ on $X$ and a closed point $x_0\in X$. Define a presheaf $\mathcal F$ by $\mathcal F(U)=\mathcal G(U)$ if $x_0\notin U$ and $\mathcal F(U)=\{s\in\mathcal G(U)\mid s_{x_0}=0\}$ otherwise. Show that $\mathcal F$ is a sheaf and that $\Supp\mathcal F=\Supp\mathcal G\setminus\{x_0\}$.}
\pf
\emph{Presheaf.} $\mathcal F(U)$ is a subgroup of $\mathcal G(U)$, and $\mathcal F(\emptyset)=\mathcal G(\emptyset)=0$. Let $V\subseteq U$ and $s\in\mathcal F(U)$. If $x_0\notin V$, then $s|_V\in\mathcal G(V)=\mathcal F(V)$. If $x_0\in V$, then $x_0\in U$, so $s_{x_0}=0$ and $(s|_V)_{x_0}=s_{x_0}=0$, hence $s|_V\in\mathcal F(V)$. So the restrictions of $\mathcal G$ induce restrictions on $\mathcal F$, and $\mathcal F$ is a subpresheaf of $\mathcal G$.

\emph{Sheaf.} Axiom (4) holds for $\mathcal F$ because it holds for $\mathcal G$ and $\mathcal F(U)\subseteq\mathcal G(U)$. For (5), let $U=\bigcup U_i$ and $s_i\in\mathcal F(U_i)$ with $s_i|_{U_i\cap U_j}=s_j|_{U_i\cap U_j}$. Since $\mathcal G$ is a sheaf, there is $s\in\mathcal G(U)$ with $s|_{U_i}=s_i$. If $x_0\notin U$, then $s\in\mathcal G(U)=\mathcal F(U)$. If $x_0\in U$, pick $i_0$ with $x_0\in U_{i_0}$; then $s_{x_0}=(s|_{U_{i_0}})_{x_0}=(s_{i_0})_{x_0}=0$ because $s_{i_0}\in\mathcal F(U_{i_0})$. So $s\in\mathcal F(U)$, and $s|_{U_i}=s_i$ in $\mathcal F$.

\emph{Stalks.} Let $x\neq x_0$. Since $x_0$ is closed, $X\setminus\{x_0\}$ is an open neighbourhood of $x$, and the open neighbourhoods $U\subseteq X\setminus\{x_0\}$ of $x$ are cofinal among all neighbourhoods. For such $U$, $\mathcal F(U)=\mathcal G(U)$ with the same restrictions. Hence $\mathcal F_x=\mathcal G_x$. At $x_0$: a germ in $\mathcal F_{x_0}$ is $s_{x_0}$ with $s\in\mathcal F(U)$, $x_0\in U$, hence $s_{x_0}=0$ in $\mathcal G_{x_0}$, so there is an open $V\ni x_0$, $V\subseteq U$, with $s|_V=0$ in $\mathcal G(V)$, hence in $\mathcal F(V)$. Thus the germ of $s$ at $x_0$ in $\mathcal F_{x_0}$ is $0$, and $\mathcal F_{x_0}=0$.

Therefore $\Supp\mathcal F=\{x\neq x_0:\mathcal G_x\neq0\}=\Supp\mathcal G\setminus\{x_0\}$. For example, if $\mathcal G$ is the sheaf of continuous real functions on $X=\R$ and $x_0=0$, then $\Supp\mathcal G=\R$ and $\Supp\mathcal F=\R\setminus\{0\}$ is not closed. \hfill$\square$

\exercise{2.6}{Let $f:X\to Y$, $g:Y\to Z$ be continuous maps, $\mathcal F$ a sheaf on $X$ and $\mathcal H$ a sheaf on $Z$. Show that $f^{-1}(g^{-1}\mathcal H)=(g\circ f)^{-1}\mathcal H$ and $g_*(f_*\mathcal F)=(g\circ f)_*\mathcal F$.}
\pf
\textbf{Direct images.} For $W\subseteq Z$ open, $(g_*(f_*\mathcal F))(W)=(f_*\mathcal F)(g^{-1}W)=\mathcal F(f^{-1}(g^{-1}W))=\mathcal F((g\circ f)^{-1}W)=((g\circ f)_*\mathcal F)(W)$, and the restriction maps of both sheaves are those of $\mathcal F$ (for $W'\subseteq W$ they are the restrictions $\mathcal F(f^{-1}g^{-1}W)\to\mathcal F(f^{-1}g^{-1}W')$). So the two sheaves are equal.

\textbf{Inverse images.} Since $f^{-1}$ is defined by a sheafification, the equality means a canonical isomorphism. For a presheaf $\mathcal P$ on $Y$ let $P_f\mathcal P$ be the presheaf $U\mapsto\varinjlim_{V\supseteq f(U)}\mathcal P(V)$ on $X$ ($V$ open in $Y$); so $f^{-1}\mathcal G=(P_f\mathcal G)^\dagger$. It is functorial in $\mathcal P$.

\emph{Stalks.} For every presheaf $\mathcal P$ on $Y$ and $x\in X$ there is an isomorphism $\varepsilon:(P_f\mathcal P)_x\to\mathcal P_{f(x)}$ sending the germ of the class $[V,s]$ ($s\in\mathcal P(V)$, $V\supseteq f(U)$, $x\in U$) to $s_{f(x)}$ (note $f(x)\in V$). It is well defined by the compatibility of the direct limits. It is surjective: given $s\in\mathcal P(V)$ with $f(x)\in V$, take $U=f^{-1}(V)\ni x$. It is injective: if $s_{f(x)}=0$, there is an open $V'\ni f(x)$, $V'\subseteq V$, with $s|_{V'}=0$; on $U'=U\cap f^{-1}(V')\ni x$ the class $[V,s]|_{U'}=[V',s|_{V'}]$ is $0$. These isomorphisms are natural in $\mathcal P$.

\emph{The comparison map.} Let $\theta:P_g\mathcal H\to g^{-1}\mathcal H$ be the canonical morphism. Define $\mu:P_{g\circ f}\mathcal H\to P_f(P_g\mathcal H)$ by sending the class $[W,s]$ ($s\in\mathcal H(W)$, $W\supseteq gf(U)$) to $[g^{-1}(W),[W,s]]$; here $V=g^{-1}(W)\supseteq f(U)$ and $[W,s]\in(P_g\mathcal H)(V)$ because $W\supseteq g(V)$. This is compatible with the identifications in the direct limits and with restrictions, so it is a morphism of presheaves on $X$. Let
\[ \varphi=P_f(\theta)\circ\mu:\ P_{g\circ f}\mathcal H\longrightarrow P_f(g^{-1}\mathcal H). \]
On stalks at $x$ we have, with $\varepsilon$ as above, $\varepsilon^{g\circ f}_x=\varepsilon^g_{f(x)}\circ\varepsilon^f_x\circ\mu_x$ (both send the germ of $[W,s]$ to $s_{gf(x)}$), hence $\mu_x$ is an isomorphism; and $\varepsilon^f_x\circ P_f(\theta)_x=\theta_{f(x)}\circ\varepsilon^f_x$, where $\theta_{f(x)}$ is an isomorphism by Prop.~2.15. So $\varphi_x$ is an isomorphism for all $x$.

\emph{Sheafifying.} Let $\vartheta_1:P_{g\circ f}\mathcal H\to(g\circ f)^{-1}\mathcal H$ and $\vartheta_2:P_f(g^{-1}\mathcal H)\to f^{-1}g^{-1}\mathcal H$ be the canonical maps. By the universal property (Def.~2.14) the morphism $\vartheta_2\varphi$ factors uniquely as $\Phi\vartheta_1$ with $\Phi:(g\circ f)^{-1}\mathcal H\to f^{-1}(g^{-1}\mathcal H)$ a morphism of sheaves. On stalks $\vartheta_{2,x}\varphi_x=\Phi_x\vartheta_{1,x}$, where $\vartheta_{1,x},\vartheta_{2,x},\varphi_x$ are isomorphisms (Prop.~2.15), so $\Phi_x$ is an isomorphism for every $x$, and $\Phi$ is an isomorphism by Prop.~2.12. This is the canonical identification $f^{-1}(g^{-1}\mathcal H)\simeq(g\circ f)^{-1}\mathcal H$ (compatible with the stalk identifications $\mathcal H_{gf(x)}$). \hfill$\square$

\exercise{2.7}{Let $\mathcal B$ be a base of open subsets of a topological space $X$ (stable under finite intersection, cf.\ Rem.~2.6). Let $\mathcal F,\mathcal G$ be two sheaves on $X$. Suppose that for every $U\in\mathcal B$ there is a homomorphism $\alpha(U):\mathcal F(U)\to\mathcal G(U)$ compatible with restrictions (between members of $\mathcal B$). Show that this extends in a unique way to a homomorphism of sheaves $\alpha:\mathcal F\to\mathcal G$. Show that if $\alpha(U)$ is surjective (resp.\ injective) for every $U\in\mathcal B$, then $\alpha$ is surjective (resp.\ injective).}
\pf
\textbf{Characterisation.} For an open $U$ and $s\in\mathcal F(U)$ we look for $t\in\mathcal G(U)$ with
\[ t|_V=\alpha(V)(s|_V)\qquad\text{for every }V\in\mathcal B,\ V\subseteq U.\qquad(\dagger) \]
\emph{Uniqueness.} The members of $\mathcal B$ contained in $U$ cover $U$ ($\mathcal B$ is a base), so by (4) at most one $t$ satisfies $(\dagger)$.

\emph{Existence.} Write $U=\bigcup_iU_i$ with $U_i\in\mathcal B$ and set $t_i=\alpha(U_i)(s|_{U_i})\in\mathcal G(U_i)$. Since $U_i\cap U_j\in\mathcal B$ and $\alpha$ commutes with restrictions between members of $\mathcal B$,
\[ t_i|_{U_i\cap U_j}=\alpha(U_i\cap U_j)(s|_{U_i\cap U_j})=t_j|_{U_i\cap U_j}. \]
By (5) there is $t\in\mathcal G(U)$ with $t|_{U_i}=t_i$. Let $V\in\mathcal B$, $V\subseteq U$. For every $i$, $V\cap U_i\in\mathcal B$ and
\[ t|_{V\cap U_i}=t_i|_{V\cap U_i}=\alpha(V\cap U_i)(s|_{V\cap U_i})=\bigl(\alpha(V)(s|_V)\bigr)|_{V\cap U_i}. \]
The sets $V\cap U_i$ cover $V$, so $t|_V=\alpha(V)(s|_V)$ by (4). Thus $t$ satisfies $(\dagger)$.

\textbf{Definition of $\alpha$.} Let $\alpha(U)(s)$ be this unique $t$. Since $(\dagger)$ is linear in $(s,t)$, $\alpha(U)$ is a group homomorphism. For $U'\subseteq U$ open, $\alpha(U)(s)|_{U'}$ satisfies $(\dagger)$ for $s|_{U'}$ (for $V\in\mathcal B$, $V\subseteq U'$: $(\alpha(U)(s)|_{U'})|_V=\alpha(U)(s)|_V=\alpha(V)(s|_V)$), so it equals $\alpha(U')(s|_{U'})$; hence $\alpha$ commutes with all restrictions and is a morphism of sheaves. If $U\in\mathcal B$, the given $\alpha(U)(s)$ satisfies $(\dagger)$ by the hypothesis that the given maps commute with restrictions between members of $\mathcal B$, so by uniqueness it equals the $t$ defined above. Thus $\alpha$ extends the given family.

\emph{Uniqueness of the extension.} If $\alpha'$ is a morphism of sheaves with $\alpha'(V)=\alpha(V)$ for $V\in\mathcal B$, then for every $U$ and $s\in\mathcal F(U)$ the section $t=\alpha'(U)(s)$ satisfies $t|_V=\alpha'(V)(s|_V)=\alpha(V)(s|_V)$ for $V\in\mathcal B$, $V\subseteq U$, i.e.\ $(\dagger)$. So $t=\alpha(U)(s)$ by uniqueness.

\textbf{Injectivity.} Suppose all $\alpha(V)$, $V\in\mathcal B$, are injective. Let $x\in X$ and $\sigma\in\mathcal F_x$ with $\alpha_x(\sigma)=0$. Write $\sigma=s_x$ with $s\in\mathcal F(U)$; since $\mathcal B$ is a base we may shrink $U$ and assume $U\in\mathcal B$. As $(\alpha(U)(s))_x=0$, there is an open $U'\ni x$, $U'\subseteq U$, with $\alpha(U)(s)|_{U'}=0$; choose $V\in\mathcal B$ with $x\in V\subseteq U'$. Then $\alpha(V)(s|_V)=\alpha(U)(s)|_V=0$ so $s|_V=0$ and $\sigma=0$. Hence every $\alpha_x$ is injective. Now if $U$ is open and $\alpha(U)(s)=0$, then $\alpha_x(s_x)=0$, so $s_x=0$ for all $x\in U$, and $s=0$ by Lemma~2.9. So $\alpha(U)$ is injective for every open $U$.

\textbf{Surjectivity.} Suppose all $\alpha(V)$, $V\in\mathcal B$, are surjective. Let $x\in X$ and $\tau\in\mathcal G_x$; write $\tau=t_x$ with $t\in\mathcal G(U)$ and shrink to $U\in\mathcal B$ containing $x$. Choose $s\in\mathcal F(U)$ with $\alpha(U)(s)=t$. Then $\tau=(\alpha(U)(s))_x=\alpha_x(s_x)$. So every $\alpha_x$ is surjective, i.e.\ $\alpha$ is surjective in the sense of Def.~2.10. (Surjectivity of $\alpha(U)$ for \emph{all} open $U$ cannot be expected, cf.\ Example~2.11.) \hfill$\square$

\emph{Conventions for this part.} Sheaves are sheaves of abelian groups unless said otherwise; all arguments work verbatim for sheaves of rings. For a morphism $f:X\to Y$ of ringed spaces and $V\subseteq Y$ open, $f^\#(V):\mathcal{O}_Y(V)\to\mathcal{O}_X(f^{-1}V)$ is the ring homomorphism given by $f^\#$. For morphisms $f:X\to Y$, $h:Y\to Z$ the composition is $(h\circ f)^\#=h_*(f^\#)\circ h^\#$, that is, $(h\circ f)^\#(W)=f^\#(h^{-1}W)\circ h^\#(W)$. For a morphism $g:U\to Y$ and an open set $U'\subseteq U$, the restriction $g|_{U'}$ is the composite of $g$ with the open immersion $U'\to U$, so $(g|_{U'})^\#(V)$ is $g^\#(V)$ followed by restriction $\mathcal{O}_X(g^{-1}V)\to\mathcal{O}_X(g^{-1}V\cap U')$.

\exercise{2.8}{(Glueing sheaves) Let $\{U_i\}_i$ be an open covering of $X$. Suppose given sheaves $\mathcal{F}_i$ on $U_i$ and isomorphisms $\varphi_{ij}:\mathcal{F}_i|_{U_i\cap U_j}\simeq\mathcal{F}_j|_{U_i\cap U_j}$ such that $\varphi_{ii}=\operatorname{Id}$ and $\varphi_{ik}=\varphi_{jk}\circ\varphi_{ij}$ on $U_i\cap U_j\cap U_k$. Show that there exists a unique sheaf $\mathcal{F}$ on $X$ with isomorphisms $\psi_i:\mathcal{F}|_{U_i}\simeq\mathcal{F}_i$ such that $\psi_j=\varphi_{ij}\circ\psi_i$ on $U_i\cap U_j$.}
\pf
``Unique'' means unique up to a unique isomorphism compatible with the $\psi_i$; this is what we prove. Taking $k=i$ in the cocycle condition gives $\varphi_{ji}\circ\varphi_{ij}=\varphi_{ii}=\operatorname{Id}$, so $\varphi_{ji}=\varphi_{ij}^{-1}$.

\emph{Construction.} For $U\subseteq X$ open, let $\mathcal{F}(U)$ be the set of families $s=(s_i)_i$ with $s_i\in\mathcal{F}_i(U\cap U_i)$ such that
\[ \varphi_{ij}\bigl(s_i|_{U\cap U_i\cap U_j}\bigr)=s_j|_{U\cap U_i\cap U_j}\qquad\text{for all } i,j. \]
This is a subgroup of $\prod_i\mathcal{F}_i(U\cap U_i)$ because the $\varphi_{ij}$ and the restrictions are homomorphisms. For $U'\subseteq U$ define the restriction componentwise, $s|_{U'}=(s_i|_{U'\cap U_i})_i$. It preserves the displayed condition because $\varphi_{ij}$ commutes with restrictions, and it is clearly functorial. Also $\mathcal{F}(\emptyset)=0$. So $\mathcal{F}$ is a presheaf.

\emph{$\mathcal{F}$ is a sheaf.} Let $U=\bigcup_a W_a$ be an open covering. (Uniqueness) If $s\in\mathcal{F}(U)$ and $s|_{W_a}=0$ for all $a$, then for every $i$ the sets $W_a\cap U_i$ cover $U\cap U_i$ and $s_i|_{W_a\cap U_i}=0$, so $s_i=0$ because $\mathcal{F}_i$ is a sheaf. (Glueing) Let $s^a\in\mathcal{F}(W_a)$ with $s^a|_{W_a\cap W_b}=s^b|_{W_a\cap W_b}$. Fix $i$. The sections $s^a_i\in\mathcal{F}_i(W_a\cap U_i)$ agree on the overlaps $W_a\cap W_b\cap U_i$, hence glue to $s_i\in\mathcal{F}_i(U\cap U_i)$ with $s_i|_{W_a\cap U_i}=s^a_i$. We check that $s=(s_i)$ lies in $\mathcal{F}(U)$. The two sections $\varphi_{ij}(s_i|_{U\cap U_i\cap U_j})$ and $s_j|_{U\cap U_i\cap U_j}$ of $\mathcal{F}_j$ over $U\cap U_i\cap U_j$ have the same restriction to each $W_a\cap U_i\cap U_j$, namely $\varphi_{ij}(s^a_i|)=s^a_j|$. By the uniqueness property of $\mathcal{F}_j$ they are equal. Clearly $s|_{W_a}=s^a$.

\emph{The isomorphisms $\psi_i$.} For $U\subseteq U_i$ open, put $\psi_i(U):\mathcal{F}(U)\to\mathcal{F}_i(U)$, $s\mapsto s_i$ (note $U\cap U_i=U$). It is a homomorphism compatible with restrictions. It is bijective: define $\chi:\mathcal{F}_i(U)\to\mathcal{F}(U)$ by $\chi(t)=(\varphi_{ij}(t|_{U\cap U_j}))_j$. This lies in $\mathcal{F}(U)$: for $j,k$, on $U\cap U_j\cap U_k$ we have $\varphi_{jk}(\varphi_{ij}(t|))=\varphi_{ik}(t|)$ by the cocycle condition, which is the required relation between the $j$-th and $k$-th components. Then $\psi_i\chi(t)=\varphi_{ii}(t)=t$. Conversely if $s\in\mathcal{F}(U)$ then the $j$-th component of $s$ is $s_j=\varphi_{ij}(s_i|_{U\cap U_j})$ by definition of $\mathcal{F}(U)$, so $\chi\psi_i(s)=s$. Hence $\psi_i:\mathcal{F}|_{U_i}\simeq\mathcal{F}_i$. Finally, for $U\subseteq U_i\cap U_j$ and $s\in\mathcal{F}(U)$ we have $\psi_j(s)=s_j=\varphi_{ij}(s_i)=\varphi_{ij}(\psi_i(s))$, i.e. $\psi_j=\varphi_{ij}\circ\psi_i$ on $U_i\cap U_j$.

\emph{Uniqueness.} Let $\mathcal{F}'$ be a sheaf with isomorphisms $\psi_i':\mathcal{F}'|_{U_i}\simeq\mathcal{F}_i$ such that $\psi'_j=\varphi_{ij}\psi'_i$ on $U_i\cap U_j$. Define $\Theta(U):\mathcal{F}'(U)\to\mathcal{F}(U)$ by $\Theta(U)(s)=(\psi'_i(s|_{U\cap U_i}))_i$. It lies in $\mathcal{F}(U)$ since $\varphi_{ij}\psi'_i(s|)=\psi_j'(s|)$ on $U\cap U_i\cap U_j$. It is a morphism of sheaves. It is injective: if all $\psi'_i(s|_{U\cap U_i})=0$ then $s|_{U\cap U_i}=0$ for all $i$, so $s=0$. It is surjective: given $(s_i)\in\mathcal{F}(U)$, put $t_i=\psi_i'^{-1}(s_i)\in\mathcal{F}'(U\cap U_i)$. On $U\cap U_i\cap U_j$ we have $\psi'_j(t_i|)=\varphi_{ij}\psi_i'(t_i|)=\varphi_{ij}(s_i|)=s_j|=\psi'_j(t_j|)$, so $t_i|=t_j|$ by injectivity of $\psi_j'$. The $t_i$ glue to $t\in\mathcal{F}'(U)$ and $\Theta(U)(t)=(s_i)$. So $\Theta:\mathcal{F}'\simeq\mathcal{F}$, and $\psi_i\circ\Theta|_{U_i}=\psi'_i$ by construction. Any isomorphism $\Theta'$ with $\psi_i\Theta'=\psi'_i$ on $U_i$ has $\Theta'(s)_i=\psi_i'(s|_{U\cap U_i})$, so $\Theta'=\Theta$. \hfill$\square$

\exercise{2.9}{Let $\mathcal{F}$ be a sheaf on $X$ whose support $Y=\operatorname{Supp}\mathcal{F}=\{x:\mathcal{F}_x\ne0\}$ is a finite set of closed points (a skyscraper sheaf). Show that for every open subset $U$ of $X$ we have $\mathcal{F}(U)=\bigoplus_{x\in U}\mathcal{F}_x$. In particular, $\mathcal{F}=0$ if and only if $\mathcal{F}(X)=0$.}
\pf
Since $\mathcal{F}_x=0$ for $x\notin Y$, the sum $\bigoplus_{x\in U}\mathcal{F}_x$ equals $\bigoplus_{y\in U\cap Y}\mathcal{F}_y$, a finite direct sum (equal to the product). The precise statement is that the canonical map
\[ \Phi_U:\mathcal{F}(U)\to\bigoplus_{x\in U}\mathcal{F}_x,\qquad s\mapsto(s_x)_{x\in U}, \]
is an isomorphism. It is a homomorphism, and $s_x=0$ for $x\notin Y$ so the image lies in the direct sum.

\emph{Injective.} If $s_y=0$ for all $y\in U\cap Y$, then $s_x=0$ for all $x\in U$, hence $s=0$ by Lemma~2.9 (applied to the sheaf $\mathcal{F}|_U$).

\emph{Surjective.} Let $(t_y)_{y\in U\cap Y}$ be given, $t_y\in\mathcal{F}_y$. For each $y$ choose an open neighbourhood $V'_y\subseteq U$ of $y$ and $s^y\in\mathcal{F}(V'_y)$ with germ $t_y$ at $y$. The set $Y\setminus\{y\}$ is finite, consisting of closed points, hence closed, so $V_y=V_y'\setminus(Y\setminus\{y\})$ is open, contains $y$, and $V_y\cap Y=\{y\}$. Replace $s^y$ by $s^y|_{V_y}$. Put $V_0=U\setminus Y$, open because $Y$ is closed (finite union of closed points), and $s^0=0\in\mathcal{F}(V_0)$. The sets $V_0$ and $V_y$ ($y\in U\cap Y$) cover $U$. For two distinct members of this family, the intersection contains no point of $Y$ (each $V_y$ meets $Y$ only in $y$, $V_0$ misses $Y$, and $V_y\cap V_z$ misses $y$ and $z$ for $y\neq z$). Since $\mathcal{F}_w=0$ for $w\notin Y$, the two sections have all germs zero on the intersection, and so coincide there by Lemma~2.9. Hence they glue to $s\in\mathcal{F}(U)$ with $s_y=s^y_y=t_y$ for $y\in U\cap Y$. So $\Phi_U(s)=(t_y)$.

\emph{The last statement.} If $\mathcal{F}=0$ then $\mathcal{F}(X)=0$. Conversely, if $\mathcal{F}(X)=0$, then $\bigoplus_{y\in Y}\mathcal{F}_y=0$ by the first part, so $\mathcal{F}_y=0$ for each $y\in Y$. But $\mathcal{F}_y\neq0$ for $y\in Y$ by definition of the support, so $Y=\emptyset$, all stalks vanish, and $\mathcal{F}=0$ by Prop.~2.12 (the morphism $\mathcal{F}\to0$ is an isomorphism on stalks). \hfill$\square$

\exercise{2.10}{Let $(X,\mathcal{O}_X)$ be a ringed topological space. Let $Y$ be a subset of $X$ endowed with the induced topology, and let $i:Y\to X$ be the canonical inclusion. Show that $(Y,i^{-1}\mathcal{O}_X)$ is a ringed topological space.}
\pf
Recall the construction: for $U\subseteq Y$ open, $P(U)=\varinjlim_{V\supseteq i(U)}\mathcal{O}_X(V)$ (limit over open $V\subseteq X$ containing $U$) is a presheaf on $Y$. Because $\mathcal{O}_X(V)$ are rings and restrictions are ring homomorphisms, $P(U)$ is a ring (a filtered direct limit of rings) and the restriction maps of $P$ are ring homomorphisms. So $P$ is a presheaf of rings. The sheaf $P^\dagger=i^{-1}\mathcal{O}_X$ of Prop.~2.15 consists of functions $f:U\to\coprod_{y\in U}P_y$ locally given by sections of $P$; the pointwise operations make $P^\dagger(U)$ a ring, since the stalks $P_y$ are rings (Def.~2.8) and the local description is stable under sums and products (take the intersection of two neighbourhoods). The restrictions are ring homomorphisms, and $\theta:P\to P^\dagger$ is a ring homomorphism. So $i^{-1}\mathcal{O}_X$ is a sheaf of rings.

It remains to show that the stalks are local rings. We have $(i^{-1}\mathcal{O}_X)_y\simeq P_y$ as rings (Prop.~2.15), and $P_y=\varinjlim_{U\ni y}\varinjlim_{V\supseteq U}\mathcal{O}_X(V)$. The pairs $(U,V)$ with $U$ open in $Y$, $V$ open in $X$, $y\in U\subseteq V$ form a directed system (order by reverse inclusion in both) and the pairs $(V\cap Y,V)$ are cofinal in it (given $(U,V)$, write $U=V'\cap Y$ with $V'$ open in $X$; then $(V\cap V'\cap Y,V\cap V')$ lies below $(U,V)$). Therefore $P_y=\varinjlim_{V\ni y}\mathcal{O}_X(V)=\mathcal{O}_{X,y}$, an isomorphism of rings. This is the formula $(i^{-1}\mathcal{O}_X)_y=\mathcal{O}_{X,i(y)}$ recalled in the text. Since $\mathcal{O}_{X,y}$ is local, so is $(i^{-1}\mathcal{O}_X)_y$. Hence $(Y,i^{-1}\mathcal{O}_X)$ is a ringed topological space. \hfill$\square$

\exercise{2.11}{Let $X,Y$ be two ringed topological spaces. Suppose given an open covering $\{U_i\}_i$ of $X$ and morphisms $f_i:U_i\to Y$ such that $f_i|_{U_i\cap U_j}=f_j|_{U_i\cap U_j}$ (as morphisms) for all $i,j$. Show that there exists a unique morphism $f:X\to Y$ such that $f|_{U_i}=f_i$ (the glueing of the $f_i$).}
\pf
\emph{Underlying map.} The hypothesis says in particular that the continuous maps $f_i$ and $f_j$ agree on $U_i\cap U_j$. So there is a unique map $f:X\to Y$ with $f|_{U_i}=f_i$. It is continuous: for $V\subseteq Y$ open, $f^{-1}V=\bigcup_i f_i^{-1}(V)$ with each $f_i^{-1}(V)$ open in $U_i$, hence open in $X$. Note $f_i^{-1}(V)=f^{-1}(V)\cap U_i$.

\emph{The sheaf map.} Let $V\subseteq Y$ be open and $a\in\mathcal{O}_Y(V)$. Put $b_i=f_i^\#(V)(a)\in\mathcal{O}_X(f^{-1}V\cap U_i)$ (recall $\mathcal{O}_{U_i}=\mathcal{O}_X|_{U_i}$). The equality $(f_i|_{U_i\cap U_j})^\#=(f_j|_{U_i\cap U_j})^\#$ at $V$, applied to $a$, says exactly that $b_i$ and $b_j$ have the same restriction to $f^{-1}V\cap U_i\cap U_j$. Since the $f^{-1}V\cap U_i$ cover $f^{-1}V$, there is a unique $f^\#(V)(a)\in\mathcal{O}_X(f^{-1}V)$ with restriction $b_i$ to each $f^{-1}V\cap U_i$. The map $f^\#(V)$ is a ring homomorphism, since sums, products and $1$ are checked after restricting to the $f^{-1}V\cap U_i$, where $f^\#(V)$ equals $f_i^\#(V)$. It is compatible with restriction in $V$ for the same reason. So $f^\#:\mathcal{O}_Y\to f_*\mathcal{O}_X$ is a morphism of sheaves of rings.

\emph{Locality.} Let $x\in X$, say $x\in U_i$. For a germ $\alpha=a_{f(x)}$ ($a\in\mathcal{O}_Y(V)$, $f(x)\in V$) the germ $f_x^\#(\alpha)$ is the germ at $x$ of $f^\#(V)(a)$, which coincides near $x$ with $f_i^\#(V)(a)$. So $f^\#_x=(f_i)^\#_x$, which is a local homomorphism. Thus $(f,f^\#)$ is a morphism of ringed spaces.

\emph{Properties.} By construction $f|_{U_i}$ has the underlying map $f_i$, and $(f|_{U_i})^\#(V)(a)=f^\#(V)(a)|_{f^{-1}V\cap U_i}=f_i^\#(V)(a)$. So $f|_{U_i}=f_i$.

\emph{Uniqueness.} If $g:X\to Y$ is a morphism with $g|_{U_i}=f_i$ for all $i$, then $g=f$ as maps, and for $a\in\mathcal{O}_Y(V)$ the section $g^\#(V)(a)$ restricts on $f^{-1}V\cap U_i$ to $(g|_{U_i})^\#(V)(a)=f_i^\#(V)(a)$. By the uniqueness in the sheaf property, $g^\#(V)(a)=f^\#(V)(a)$. \hfill$\square$

\exercise{2.12}{Let $f:X\to Y$ be a morphism of ringed topological spaces. Let $V$ be an open subset of $Y$ containing $f(X)$. Show that there exists a unique morphism $g:X\to V$ whose composition with the open immersion $V\to Y$ is $f$.}
\pf
Let $j:V\to Y$ be the inclusion; as a morphism of ringed spaces $(V,\mathcal{O}_Y|_V)\to(Y,\mathcal{O}_Y)$, $j^\#(W):\mathcal{O}_Y(W)\to\mathcal{O}_Y(W\cap V)$ is the restriction.

\emph{Existence.} Let $g:X\to V$ be the map $x\mapsto f(x)$; it is continuous because the topology of $V$ is induced from $Y$ and $g^{-1}(W\cap V)=f^{-1}(W)$. Since $f(X)\subseteq V$ we have $f^{-1}(W)=f^{-1}(W\cap V)=g^{-1}(W\cap V)$ for every open $W\subseteq Y$. For $W\subseteq V$ open (it is open in $V$ and in $Y$) define
\[ g^\#(W)=f^\#(W):\mathcal{O}_V(W)=\mathcal{O}_Y(W)\to\mathcal{O}_X(f^{-1}W)=\mathcal{O}_X(g^{-1}W). \]
This is a morphism of sheaves of rings on $V$ (compatible with restrictions between open subsets of $V$ since $f^\#$ is). For $x\in X$, $g^\#_x=f^\#_x$ since $\mathcal{O}_{V,g(x)}=\mathcal{O}_{Y,f(x)}$; this is local. So $g$ is a morphism of ringed spaces. For $W\subseteq Y$ open and $a\in\mathcal{O}_Y(W)$,
\[ (j\circ g)^\#(W)(a)=g^\#(W\cap V)(a|_{W\cap V})=f^\#(W\cap V)(a|_{W\cap V})=f^\#(W)(a)|_{f^{-1}(W\cap V)}=f^\#(W)(a), \]
using that $f^\#$ commutes with restriction and $f^{-1}(W\cap V)=f^{-1}(W)$. Hence $j\circ g=f$.

\emph{Uniqueness.} Let $g'$ satisfy $j\circ g'=f$. As $j$ is injective, $g'=g$ as maps. For $W\subseteq V$ open, $j^\#(W)$ is the identity of $\mathcal{O}_Y(W)$, so $f^\#(W)=(j\circ g')^\#(W)=g'^\#(W)\circ j^\#(W)=g'^\#(W)$. Thus $g'^\#=f^\#|_V=g^\#$. \hfill$\square$

\exercise{2.13}{Let $f:X\to Y$ be a continuous map of topological spaces. Let $\mathcal{F}$ be a sheaf on $X$ and $\mathcal{G}$ a sheaf on $Y$.
\begin{enumerate}[label=(\alph*),leftmargin=*]
\item Show that there exist canonical morphisms of sheaves $\mathcal{G}\to f_*f^{-1}\mathcal{G}$ and $f^{-1}f_*\mathcal{F}\to\mathcal{F}$. Moreover, if $f$ is a closed immersion, then the first morphism is surjective. If $f$ is an open immersion, then the second morphism is an isomorphism.
\item Using these maps, show that there exists a canonical bijection $\operatorname{Mor}_Y(\mathcal{G},f_*\mathcal{F})\simeq\operatorname{Mor}_X(f^{-1}\mathcal{G},\mathcal{F})$.
\end{enumerate}}
\pf
\emph{Remark on hypotheses.} Here ``closed immersion'' and ``open immersion'' are meant topologically (a homeomorphism onto a closed, resp.\ open, subset). In fact the proofs below only use that $f$ is a homeomorphism onto its image (an embedding), plus closedness of the image for the surjectivity claim.

\emph{Notation.} For $U\subseteq X$ open let $P_{\mathcal{G}}(U)=\varinjlim_{W\supseteq f(U)}\mathcal{G}(W)$ ($W$ open in $Y$), and write $[s]_W\in P_{\mathcal{G}}(U)$ for the class of $s\in\mathcal{G}(W)$ when $f(U)\subseteq W$. So $f^{-1}\mathcal{G}=P_{\mathcal{G}}^\dagger$ with $\theta=\theta_{\mathcal{G}}:P_{\mathcal{G}}\to f^{-1}\mathcal{G}$ as in Prop.~2.15.

\emph{Step 0: stalks.} For $x\in X$ and $y=f(x)$ the map $P_{\mathcal{G},x}\to\mathcal{G}_y$, sending the germ at $x$ of $[s]_W$ to $s_y$, is well defined and bijective. Surjective: a germ $s_y$ with $s\in\mathcal{G}(W)$ is the image of the germ of $[s]_W\in P_{\mathcal{G}}(f^{-1}W)$. Injective: if $s_y=0$ there is an open $W'\ni y$, $W'\subseteq W$ with $s|_{W'}=0$; then, for a representative $[s]_W\in P_{\mathcal{G}}(U)$ with $x\in U$, its restriction to $U\cap f^{-1}W'$ is $[s]_W=[s|_{W'}]_{W'}=0$ in $P_{\mathcal{G}}(U\cap f^{-1}W')$ (note $f(U\cap f^{-1}W')\subseteq W'$), so the germ at $x$ is $0$. Composing with $\theta_x$ (an isomorphism by Prop.~2.15) we get the isomorphism $\kappa:(f^{-1}\mathcal{G})_x\simeq\mathcal{G}_{f(x)}$ of the text, characterized by $\kappa\bigl(\theta([s]_W)_x\bigr)=s_{f(x)}$.

\emph{Step 1: the morphism $\eta:\mathcal{G}\to f_*f^{-1}\mathcal{G}$.} For $V\subseteq Y$ open we have $f(f^{-1}V)\subseteq V$, so $[s]_V\in P_{\mathcal{G}}(f^{-1}V)$ is defined for $s\in\mathcal{G}(V)$. Put $\eta(V)(s)=\theta(f^{-1}V)([s]_V)\in(f^{-1}\mathcal{G})(f^{-1}V)=(f_*f^{-1}\mathcal{G})(V)$. This is a homomorphism, and it is compatible with restrictions in $V$ because the classes $[s]_V$ and $\theta$ are.

\emph{Step 2: the morphism $\varepsilon:f^{-1}f_*\mathcal{F}\to\mathcal{F}$.} Here $P_{f_*\mathcal{F}}(U)=\varinjlim_{W\supseteq f(U)}\mathcal{F}(f^{-1}W)$. Define $\varepsilon'(U):P_{f_*\mathcal{F}}(U)\to\mathcal{F}(U)$ by $[t]_W\mapsto t|_U$ for $t\in\mathcal{F}(f^{-1}W)$, which makes sense as $U\subseteq f^{-1}W$. It is compatible with the transition maps of the direct limit (restriction in $W$ followed by restriction to $U$ equals restriction to $U$) and with restriction in $U$. So $\varepsilon'$ is a morphism of presheaves, and since $\mathcal{F}$ is a sheaf, the universal property of Def.~2.14 gives a unique $\varepsilon:f^{-1}f_*\mathcal{F}\to\mathcal{F}$ with $\varepsilon\circ\theta_{f_*\mathcal{F}}=\varepsilon'$.

\emph{Step 3: a lemma.} \textbf{Lemma.} Let $f$ be an embedding (homeomorphism onto its image with the induced topology), $\mathcal{H}$ a sheaf on $X$, $x\in X$, $y=f(x)$. Then $c:(f_*\mathcal{H})_y\to\mathcal{H}_x$, $t_y\mapsto t_x$ (for $t\in\mathcal{H}(f^{-1}V)$) is bijective. If moreover $f(X)$ is closed and $y\notin f(X)$ then $(f_*\mathcal{H})_y=0$.

\emph{Proof.} Surjective: a germ $h_x$ with $h\in\mathcal{H}(U)$; since $f$ is an embedding, $U=f^{-1}V$ for some open $V\subseteq Y$, so $h\in(f_*\mathcal{H})(V)$ and $c(h_y)=h_x$. Injective: if $t\in\mathcal{H}(f^{-1}V)$ and $t_x=0$, there is an open $U'\ni x$, $U'\subseteq f^{-1}V$ with $t|_{U'}=0$. Write $U'=f^{-1}V'$ and replace $V'$ by $V'\cap V$ (which does not change $f^{-1}V'$ because $U'\subseteq f^{-1}V$). Then $V'\subseteq V$ is a neighbourhood of $y$ and $t|_{V'}=t|_{f^{-1}V'}=0$ in $(f_*\mathcal{H})(V')$, so $t_y=0$. For the last claim, if $f(X)$ is closed and $y\notin f(X)$, then $y$ has an open neighbourhood $V$ disjoint from $f(X)$, and $(f_*\mathcal{H})(V)=\mathcal{H}(\emptyset)=0$; so the stalk is $0$. $\square$

\emph{Step 4: closed immersions.} Let $f$ be an embedding with closed image. We show $\eta_y:\mathcal{G}_y\to(f_*f^{-1}\mathcal{G})_y$ is surjective for every $y\in Y$. If $y\notin f(X)$, the target is $0$ by the Lemma. If $y=f(x)$, take $s\in\mathcal{G}(V)$ with $y\in V$. Then $\eta_y(s_y)$ is the germ at $y$ of $\theta([s]_V)\in(f^{-1}\mathcal{G})(f^{-1}V)$; applying $c$ gives the germ of this section at $x$, and then $\kappa$ gives $s_{f(x)}=s_y$. Hence $\kappa\circ c\circ\eta_y=\operatorname{Id}$ on $\mathcal{G}_y$; since $\kappa\circ c$ is bijective (Step 0 and the Lemma), $\eta_y$ is bijective, in particular surjective. So $\eta$ is surjective in the sense of the text (it is even an isomorphism on the stalks at points of $f(X)$).

\emph{Step 5: open immersions.} Let $f$ be an embedding (e.g.\ an open immersion). For $x\in X$ and $y=f(x)$, the stalk map $\varepsilon_x:(f^{-1}f_*\mathcal{F})_x\to\mathcal{F}_x$ sends $\theta([t]_W)_x$ (with $t\in\mathcal{F}(f^{-1}W)$) to $(\varepsilon'[t]_W)_x=t_x$. Under the isomorphism $\kappa:(f^{-1}f_*\mathcal{F})_x\simeq(f_*\mathcal{F})_y$ (Step 0 for the sheaf $f_*\mathcal{F}$, which sends that element to $t_y$) this says $\varepsilon_x\circ\kappa^{-1}=c$, which is bijective by the Lemma. So $\varepsilon_x$ is bijective for all $x$, and $\varepsilon$ is an isomorphism by Prop.~2.12. This proves (a).

\emph{(b) The bijection.} Define $\Phi:\operatorname{Mor}_X(f^{-1}\mathcal{G},\mathcal{F})\to\operatorname{Mor}_Y(\mathcal{G},f_*\mathcal{F})$ by $\Phi(\beta)=f_*(\beta)\circ\eta_{\mathcal{G}}$, i.e. $\Phi(\beta)(V)(s)=\beta(f^{-1}V)\bigl(\theta([s]_V)\bigr)$. Define also, for a morphism $\alpha:\mathcal{G}\to f_*\mathcal{F}$, the morphism $f^{-1}\alpha:f^{-1}\mathcal{G}\to f^{-1}f_*\mathcal{F}$ as the unique morphism with $f^{-1}\alpha\circ\theta_{\mathcal{G}}=\theta_{f_*\mathcal{F}}\circ P(\alpha)$, where $P(\alpha)(U)([s]_W)=[\alpha(W)(s)]_W$ (existence and uniqueness by Def.~2.14, since $f^{-1}f_*\mathcal{F}$ is a sheaf), and put $\Psi(\alpha)=\varepsilon_{\mathcal{F}}\circ f^{-1}\alpha$.

We show $\Phi$ is bijective with inverse $\Psi$. By Def.~2.14, $\beta\mapsto\beta\circ\theta$ is a bijection from $\operatorname{Mor}_X(f^{-1}\mathcal{G},\mathcal{F})$ to the set of presheaf morphisms $\gamma:P_{\mathcal{G}}\to\mathcal{F}$. Next, presheaf morphisms $\gamma$ correspond bijectively to morphisms $\alpha:\mathcal{G}\to f_*\mathcal{F}$ via
\[ \alpha(V)(s)=\gamma(f^{-1}V)([s]_V),\qquad \gamma(U)([s]_W)=\alpha(W)(s)|_U . \]
Indeed, given $\gamma$ the first formula defines a homomorphism compatible with restrictions (as $\gamma$ is and $[s|_{V'}]_{V'}=[s]_V$ in $P_{\mathcal{G}}(f^{-1}V')$). Given $\alpha$ the second formula is well defined on the direct limit (if $W'\subseteq W$, $\alpha(W')(s|_{W'})|_U=\alpha(W)(s)|_U$) and compatible with restriction in $U$. Starting from $\gamma$, then $\alpha$, then recomputing, $\alpha(W)(s)|_U=\gamma(f^{-1}W)([s]_W)|_U=\gamma(U)([s]_W|_U)=\gamma(U)([s]_W)$, and starting from $\alpha$ one gets $\gamma(f^{-1}V)([s]_V)=\alpha(V)(s)$; so the two constructions are inverse. Composing the two bijections, $\beta$ goes to $\alpha(V)(s)=\beta(f^{-1}V)(\theta([s]_V))$, which is $\Phi(\beta)$. Hence $\Phi$ is bijective.

Finally $\Psi=\Phi^{-1}$: let $\alpha$ be given and $\gamma_\alpha$ the presheaf morphism corresponding to it. Then $\Psi(\alpha)\circ\theta_{\mathcal{G}}=\varepsilon\circ\theta_{f_*\mathcal{F}}\circ P(\alpha)=\varepsilon'\circ P(\alpha)$, and $\varepsilon'P(\alpha)(U)([s]_W)=\alpha(W)(s)|_U=\gamma_\alpha(U)([s]_W)$. So $\Psi(\alpha)\circ\theta=\gamma_\alpha$, which by the above means $\Phi(\Psi(\alpha))=\alpha$. As $\Phi$ is bijective, $\Psi=\Phi^{-1}$. Thus $\operatorname{Mor}_Y(\mathcal{G},f_*\mathcal{F})\simeq\operatorname{Mor}_X(f^{-1}\mathcal{G},\mathcal{F})$ canonically, both directions being expressed through the maps $\eta$ and $\varepsilon$ of (2.1) (the same argument works for sheaves of rings and ring morphisms). \hfill$\square$

\exercise{2.14}{(Quotient space) Let a group $G$ act on a ringed topological space $(X,\mathcal{O}_X)$ (a homomorphism $G\to\operatorname{Aut}(X)$). A quotient of $X$ by $G$ is a ringed space $(Y,\mathcal{O}_Y)$ with a morphism $p:X\to Y$ such that $p=p\circ\sigma$ for every $\sigma\in G$, and every morphism of ringed spaces $f:X\to Z$ with $f=f\circ\sigma$ for all $\sigma\in G$ factors uniquely as $f=\tilde f\circ p$. Show the quotient exists: let $Y=X/G$ with the quotient topology and $p:X\to Y$ the projection.
\begin{enumerate}[label=(\alph*),leftmargin=*]
\item Show that $p$ is an open map.
\item Show that $G$ acts on the ringed space $p^{-1}(V)$ for $V\subseteq Y$ open. In particular $G$ acts on the ring $\mathcal{O}_X(p^{-1}(V))$.
\item For $V\subseteq Y$ open, let $\mathcal{O}_Y(V)=\mathcal{O}_X(p^{-1}(V))^G$ (the $G$-invariant elements). Show that $\mathcal{O}_Y$ is a sheaf of rings on $Y$.
\item Let $y\in Y$ and fix $x\in p^{-1}(y)$. Show that $\mathcal{O}_{Y,y}\subseteq\mathcal{O}_{X,x}$. Show that $\mathfrak{m}_y=\mathfrak{m}_x\cap\mathcal{O}_{Y,y}$ is the unique maximal ideal of $\mathcal{O}_{Y,y}$.
\item Show that $(Y,\mathcal{O}_Y)$ is a ringed space, that $p$ is a morphism of ringed spaces from $X$ to $Y$, and that $p$ has the universal property of the quotient.
\end{enumerate}}
\pf
For $\sigma\in G$ we write $\sigma=(\sigma,\sigma^\#)$ for the corresponding automorphism of $(X,\mathcal{O}_X)$. As $\sigma$ is invertible as a morphism, $\sigma$ is a homeomorphism and $\sigma^\#:\mathcal{O}_X\to\sigma_*\mathcal{O}_X$ is an isomorphism of sheaves, i.e.\ $\sigma^\#(W):\mathcal{O}_X(W)\to\mathcal{O}_X(\sigma^{-1}W)$ is a ring isomorphism for each open $W$; and $\sigma^\#_z:\mathcal{O}_{X,\sigma(z)}\to\mathcal{O}_{X,z}$ is an isomorphism of rings for every $z$ (the stalk of $\sigma_*\mathcal{O}_X$ at $\sigma(z)$ is $\mathcal{O}_{X,z}$ because $\sigma$ is a homeomorphism). The topology on $Y$ is such that $p$ is continuous, and $p\circ\sigma=p$ as maps.

\textbf{(a)} Let $U\subseteq X$ be open. Then $p^{-1}(p(U))=\bigcup_{\sigma\in G}\sigma(U)$: indeed $z$ lies in the left side iff $p(z)=p(u)$ for some $u\in U$ iff $z=\sigma u$ for some $\sigma\in G,u\in U$. Each $\sigma(U)$ is open since $\sigma$ is a homeomorphism. So $p^{-1}(p(U))$ is open, i.e.\ $p(U)$ is open in $Y$.

\textbf{(b)} Let $W=p^{-1}(V)$. Since $p\sigma=p$, $\sigma(W)\subseteq W$ for every $\sigma\in G$, and applying this to $\sigma^{-1}$ gives $\sigma(W)=W=\sigma^{-1}(W)$. Thus $\sigma$ restricts to a homeomorphism of $W$, and $\sigma^\#$ restricts to an isomorphism $\mathcal{O}_X|_W\to(\sigma|_W)_*(\mathcal{O}_X|_W)$; so $\sigma|_W$ is an automorphism of the ringed space $(W,\mathcal{O}_X|_W)$. In particular each $\sigma^\#(W)$ is a ring automorphism of $\mathcal{O}_X(W)$. By the composition rule, $(\sigma\tau)^\#(W)=\tau^\#(W)\circ\sigma^\#(W)$, so we get an action by ring automorphisms (a right action; the left action $s\mapsto(\sigma^{-1})^\#(W)(s)$ has the same invariants). The only thing we use is the invariant ring $\mathcal{O}_X(W)^G=\{s:\sigma^\#(W)(s)=s\ \forall\sigma\in G\}$, a subring of $\mathcal{O}_X(W)$.

\textbf{(c)} If $V'\subseteq V$ are open and $W'=p^{-1}V'\subseteq W=p^{-1}V$, then $\sigma^\#(W')(s|_{W'})=\sigma^\#(W)(s)|_{\sigma^{-1}W'}=\sigma^\#(W)(s)|_{W'}$ since $\sigma^\#$ is compatible with restrictions and $\sigma^{-1}W'=W'$. So the restriction of an invariant section is invariant, and $\mathcal{O}_Y$ is a presheaf of rings (with $\mathcal{O}_Y(\emptyset)=\mathcal{O}_X(\emptyset)=0$). Sheaf property: let $V=\bigcup_aV_a$ and put $W_a=p^{-1}V_a$, so $p^{-1}V=\bigcup W_a$ and $p^{-1}(V_a\cap V_b)=W_a\cap W_b$. If $s\in\mathcal{O}_Y(V)$ restricts to $0$ on each $V_a$, then $s|_{W_a}=0$ for all $a$ and $s=0$ in $\mathcal{O}_X(p^{-1}V)$. If $s_a\in\mathcal{O}_Y(V_a)$ agree on the overlaps $V_a\cap V_b$, they agree on $W_a\cap W_b$, so glue to $s\in\mathcal{O}_X(p^{-1}V)$ with $s|_{W_a}=s_a$. For $\sigma\in G$, $\sigma^\#(p^{-1}V)(s)|_{W_a}=\sigma^\#(W_a)(s_a)=s_a$ (as in the displayed formula), so $\sigma^\#(s)$ and $s$ have the same restriction to every $W_a$ and hence are equal. Thus $s\in\mathcal{O}_Y(V)$. So $\mathcal{O}_Y$ is a sheaf of rings.

\textbf{(d)} For $V\ni y$ open and $s\in\mathcal{O}_Y(V)\subseteq\mathcal{O}_X(p^{-1}V)$, since $x\in p^{-1}V$, we can take the germ of $s$ at $x$. This is compatible with restricting to smaller $V$, so it defines a ring homomorphism $\iota:\mathcal{O}_{Y,y}\to\mathcal{O}_{X,x}$ (it is $p^\#_x$ of part (e)).

\emph{$\iota$ is injective.} Suppose $s\in\mathcal{O}_Y(V)$ has $s_x=0$ in $\mathcal{O}_{X,x}$. Then there is an open $U\ni x$, $U\subseteq p^{-1}V$, with $s|_U=0$. For $\sigma\in G$, using $\sigma^\#(p^{-1}V)(s)=s$ and compatibility with restriction,
\[ \sigma^\#(\sigma U)\bigl(s|_{\sigma U}\bigr)=\sigma^\#(p^{-1}V)(s)|_{\sigma^{-1}\sigma U}=s|_U=0 . \]
As $\sigma^\#(\sigma U)$ is injective, $s|_{\sigma U}=0$. By part (a), $V'=p(U)$ is an open neighbourhood of $y$ contained in $V$, and $p^{-1}V'=\bigcup_\sigma\sigma U$ (proof of (a)). So $s|_{p^{-1}V'}=0$, hence $s|_{V'}=0$ in $\mathcal{O}_Y(V')$ and $s_y=0$.

So we identify $\mathcal{O}_{Y,y}$ with the subring $A=\iota(\mathcal{O}_{Y,y})\subseteq\mathcal{O}_{X,x}$ and $\mathfrak{m}_y:=\mathfrak{m}_x\cap A$, a proper ideal of $A$ (it does not contain $1$). We show every $a\in A\setminus\mathfrak{m}_x$ is a unit of $A$. Then the non-units of $A$ are exactly the elements of $\mathfrak{m}_x\cap A=\mathfrak{m}_y$ (elements of $\mathfrak{m}_x$ are not units of $\mathcal{O}_{X,x}$, hence not of $A$), so the non-units form an ideal and $A$ is local with maximal ideal $\mathfrak{m}_y$, as desired.

Let $a=s_x$ with $s\in\mathcal{O}_Y(V)$, $a\notin\mathfrak{m}_x$; thus $a$ is a unit of the local ring $\mathcal{O}_{X,x}$. Let $N=\{z\in p^{-1}V:\ s_z\in\mathcal{O}_{X,z}^\times\}$. If $s_z$ is a unit, there is an open $U_z\ni z$ and $t^z\in\mathcal{O}_X(U_z)$ with $s|_{U_z}t^z=1$; then every point of $U_z$ is in $N$, so $N$ is open and $x\in N$. Moreover $N$ is $G$-stable: from $\sigma^\#(p^{-1}V)(s)=s$ we get $\sigma^\#_z(s_{\sigma z})=s_z$, and $\sigma_z^\#$ is a ring isomorphism, so $s_{\sigma z}$ is a unit iff $s_z$ is. Hence $N=p^{-1}(V')$ with $V'=p(N)$ open (by (a)) and $y\in V'$. For $z,z'\in N$ the sections $t^z,t^{z'}$ coincide on $U_z\cap U_{z'}$ (both are the inverse of $s$ in the ring $\mathcal{O}_X(U_z\cap U_{z'})$), so they glue to $t\in\mathcal{O}_X(N)$ with $s|_Nt=1$. For $\sigma\in G$, $\sigma^\#(N)(t)$ is an inverse of $\sigma^\#(N)(s|_N)=s|_N$, so equals $t$ by uniqueness of inverses. Thus $t\in\mathcal{O}_Y(V')$, $s|_{V'}\,t=1$ in $\mathcal{O}_Y(V')$, and $s_y$ is a unit of $\mathcal{O}_{Y,y}$; i.e.\ $a$ is a unit of $A$.

\textbf{(e)} By (c) and (d), $(Y,\mathcal{O}_Y)$ is a ringed topological space. Let $p^\#(V):\mathcal{O}_Y(V)=\mathcal{O}_X(p^{-1}V)^G\hookrightarrow\mathcal{O}_X(p^{-1}V)=(p_*\mathcal{O}_X)(V)$ be the inclusion; this is a morphism of sheaves of rings. The stalk map $p^\#_x:\mathcal{O}_{Y,p(x)}\to\mathcal{O}_{X,x}$ is the map $\iota$ of (d), and $(p_x^\#)^{-1}(\mathfrak{m}_x)=\iota^{-1}(\mathfrak{m}_x\cap A)=\mathfrak{m}_{p(x)}$, so it is local. So $p:X\to Y$ is a morphism of ringed spaces. For $\sigma\in G$, $(p\sigma)^\#(V)(a)=\sigma^\#(p^{-1}V)(p^\#(V)(a))=a=p^\#(V)(a)$ because $a$ is invariant, so $p\circ\sigma=p$ as morphisms.

\emph{Universal property.} Let $f:X\to Z$ be a morphism with $f\circ\sigma=f$ for all $\sigma\in G$ (equality of morphisms, so the maps are equal and $(f\sigma)^\#=f^\#$). The map $f$ is constant on $G$-orbits, so $f=\tilde f\circ p$ for a unique map $\tilde f:Y\to Z$, and $\tilde f$ is continuous because $p^{-1}(\tilde f^{-1}W)=f^{-1}W$ is open for each open $W\subseteq Z$ and $Y$ has the quotient topology. For $W\subseteq Z$ open and $u\in\mathcal{O}_Z(W)$ put $b=f^\#(W)(u)\in\mathcal{O}_X(f^{-1}W)=\mathcal{O}_X(p^{-1}\tilde f^{-1}W)$. From $(f\sigma)^\#=f^\#$ we get $\sigma^\#(f^{-1}W)(b)=(f\sigma)^\#(W)(u)=b$, so $b$ is invariant, i.e.\ $b\in\mathcal{O}_Y(\tilde f^{-1}W)=(\tilde f_*\mathcal{O}_Y)(W)$. Define $\tilde f^\#(W)(u)=b$. This is a morphism of sheaves of rings $\mathcal{O}_Z\to\tilde f_*\mathcal{O}_Y$, since $f^\#$ is one and $\mathcal{O}_Y\subseteq p_*\mathcal{O}_X$ is a subsheaf of rings. For $x\in X$, $y=p(x)$, we have $\iota\circ\tilde f^\#_y=f^\#_x$ as maps $\mathcal{O}_{Z,f(x)}\to\mathcal{O}_{X,x}$ (both send the germ of $u$ to the germ at $x$ of $f^\#(W)(u)$). If $c\in\mathfrak{m}_{f(x)}$ then $\iota(\tilde f^\#_y(c))=f^\#_x(c)\in\mathfrak{m}_x$, and it lies in $A=\iota(\mathcal{O}_{Y,y})$, so $\iota(\tilde f_y^\#(c))\in\mathfrak{m}_x\cap A=\iota(\mathfrak{m}_y)$, and by injectivity of $\iota$, $\tilde f^\#_y(c)\in\mathfrak{m}_y$. So $\tilde f_y^\#$ is local and $\tilde f$ is a morphism of ringed spaces. Finally $(\tilde f\circ p)^\#(W)(u)=p^\#(\tilde f^{-1}W)(\tilde f^\#(W)(u))=f^\#(W)(u)$, so $\tilde f\circ p=f$.

\emph{Uniqueness.} Let $h:Y\to Z$ be a morphism with $h\circ p=f$. Since $p$ is surjective, $h=\tilde f$ as maps. For $u\in\mathcal{O}_Z(W)$ we have $f^\#(W)(u)=(hp)^\#(W)(u)=p^\#(h^{-1}W)(h^\#(W)(u))=h^\#(W)(u)$, viewed inside $\mathcal{O}_X(p^{-1}h^{-1}W)$, because $p^\#$ is an inclusion. Hence $h^\#(W)(u)=\tilde f^\#(W)(u)$ and $h=\tilde f$. Therefore $(Y,\mathcal{O}_Y)$ with $p$ is a quotient of $X$ by $G$. \hfill$\square$

\section*{2.3 Schemes}

Throughout, rings are commutative with $1$, and results of the text are cited by their numbers in Liu (``Prop.~3.25'', ``Lemma~3.40'', \dots). For a morphism of schemes $g:Z\to X$ the naturality (3.4) of the map $\rho$ of (3.3) reads
\[ \rho(f\circ g)=g^{\#}(X)\circ\rho(f)\qquad(f:X\to Y). \]
We refer to this identity as (N).

Honest notes. Ex.~3.3(b) is printed with a closure bar that is essential: the image of $f$ itself is in general not equal to $V(\operatorname{Ker}\varphi)$, and the solution proves the statement for $\overline{\operatorname{Im} f}$. Ex.~3.10 (``$X$ is affine iff $n=0$'') is false for $A=0$ (then $\mathbb{P}^n_A=\emptyset$ is affine), so the second half is proved for $A\neq0$. In Ex.~3.11(b) the hypothesis ``$B_e$ finitely generated over $A$'' is read as ``finitely generated $A$-module'' and ``$B_{ed}$ is generated by $B_e^d$'' as ``$B_{ed}$ is spanned over $A$ by products of $d$ elements of $B_e$'', where $B$ is a graded $A$-algebra.

\exercise{3.1}{Let $X$ be an open subscheme of an affine scheme $Y$. Show that the canonical morphism $X\to Y$ corresponds, by the map $\rho$ of Proposition 3.25, to the restriction $\mathcal{O}_Y(Y)\to\mathcal{O}_X(X)$.}
\pf
By definition (Def.~3.10) $X$ is an open subset of $Y$ with the structure sheaf $\mathcal{O}_Y|_X$, and the canonical morphism is $i=(i,i^{\#})$, where $i:X\to Y$ is the inclusion and $i^{\#}:\mathcal{O}_Y\to i_*(\mathcal{O}_Y|_X)$ is given, for each open $V\subseteq Y$, by the restriction map $\mathcal{O}_Y(V)\to\mathcal{O}_Y(V\cap X)=(i_*\mathcal{O}_X)(V)$. By definition of $\rho$ in (3.3),
\[ \rho(i)=i^{\#}(Y):\ \mathcal{O}_Y(Y)\longrightarrow (i_*\mathcal{O}_X)(Y)=\mathcal{O}_X(X), \]
and by the description of $i^{\#}$ this is the restriction map $\mathcal{O}_Y(Y)\to\mathcal{O}_Y(X)=\mathcal{O}_X(X)$. Since $Y$ is affine, Prop.~3.25 says that $i$ is the unique morphism $X\to Y$ with this image under $\rho$. \hfill$\square$

\exercise{3.2}{Let $U=\operatorname{Spec}B$ be an affine open subscheme of $X=\operatorname{Spec}A$. Show that the restriction $A\to B$ is a flat homomorphism.}
\pf
Here ``$U=\operatorname{Spec}B$'' means that we are given an isomorphism of schemes $\sigma:\operatorname{Spec}B\to(U,\mathcal{O}_X|_U)$; the restriction $A=\mathcal{O}_X(X)\to\mathcal{O}_X(U)\simeq B$ is the ring homomorphism $\varphi:=\tau^{\#}(X)=\rho(\tau)$, where $\tau:\operatorname{Spec}B\to X$ is $\sigma$ followed by the canonical open immersion $U\to X$ (use Ex.~3.1 and (N)). By Lemma~3.23, $\tau=f_\varphi$ is the morphism of Prop.~3.14 attached to $\varphi$. Let $\mathfrak{q}\in\operatorname{Spec}B$ and $\mathfrak{p}=\varphi^{-1}(\mathfrak{q})=\tau(\mathfrak{q})$. By Prop.~3.14 the stalk map $\tau^{\#}_{\mathfrak{q}}:\mathcal{O}_{X,\mathfrak{p}}=A_{\mathfrak{p}}\to\mathcal{O}_{\operatorname{Spec}B,\mathfrak{q}}=B_{\mathfrak{q}}$ is the localization of $\varphi$. Since $\tau$ is an isomorphism onto the open subscheme $U$, the stalk maps of $\tau$ are isomorphisms. Hence
\[ A_{\mathfrak{p}}\xrightarrow{\ \sim\ }B_{\mathfrak{q}}\quad\text{for every }\mathfrak{q}\in\operatorname{Spec}B,\ \mathfrak{p}=\varphi^{-1}(\mathfrak{q}). \]
Call this (1).

Now let $u:N'\to N$ be an injective homomorphism of $A$-modules, and let $K$ be the kernel of $u\otimes_A B:N'\otimes_AB\to N\otimes_AB$, a $B$-module. Fix $\mathfrak{q}$ and $\mathfrak{p}$ as above. Localization at $\mathfrak{q}$ is exact, so $K_{\mathfrak{q}}$ is the kernel of $N'\otimes_AB_{\mathfrak{q}}\to N\otimes_AB_{\mathfrak{q}}$. Elements of $A\setminus\mathfrak{p}$ are invertible in $B_{\mathfrak{q}}$, so $N'\otimes_AB_{\mathfrak{q}}=N'_{\mathfrak{p}}\otimes_{A_{\mathfrak{p}}}B_{\mathfrak{q}}$ and likewise for $N$. By (1), $B_{\mathfrak{q}}\simeq A_{\mathfrak{p}}$ as $A_{\mathfrak{p}}$-modules, so the map in question is isomorphic to $N'_{\mathfrak{p}}\to N_{\mathfrak{p}}$, which is injective (localization is exact). Thus $K_{\mathfrak{q}}=0$ for all $\mathfrak{q}$. A module all of whose localizations at primes vanish is zero (if $0\neq x\in K$, the ideal $\operatorname{Ann}(x)$ lies in a maximal ideal $\mathfrak{q}$ and $x/1\neq0$ in $K_{\mathfrak{q}}$). Hence $K=0$ and $B$ is flat over $A$. \hfill$\square$

\exercise{3.3}{Closure of a subset.

(a) Let $F$ be a subset of an affine scheme $\operatorname{Spec}A$. Show that the closure $\overline F$ of $F$ in (the underlying topological space of) $\operatorname{Spec}A$ is equal to $V(I)$, where $I=\cap_{\mathfrak{p}\in F}\mathfrak{p}$.

(b) Let $\varphi:A\to B$ be a ring homomorphism. Let $f:\operatorname{Spec}B\to\operatorname{Spec}A$ be the morphism of schemes associated to $\varphi$. Show that $\overline{\operatorname{Im}f}=V(\operatorname{Ker}\varphi)$. Study the situation where $B$ is the localization of $A$ at a prime ideal.}
\pf
(a) (For $F=\emptyset$ we have $I=A$ and $V(I)=\emptyset=\overline F$.) The set $V(I)$ is closed, and it contains $F$ because $I\subseteq\mathfrak{p}$ for every $\mathfrak{p}\in F$. Conversely let $V(J)$ be any closed set containing $F$. Then $J\subseteq\mathfrak{p}$ for all $\mathfrak{p}\in F$, so $J\subseteq I$ and $V(I)\subseteq V(J)$. Thus $V(I)$ is the smallest closed set containing $F$, i.e.\ $V(I)=\overline F$.

(b) We have $f(\mathfrak{q})=\varphi^{-1}(\mathfrak{q})$. By (a), $\overline{\operatorname{Im}f}=V(I)$ with
\[ I=\bigcap_{\mathfrak{q}\in\operatorname{Spec}B}\varphi^{-1}(\mathfrak{q})=\varphi^{-1}\Bigl(\bigcap_{\mathfrak{q}}\mathfrak{q}\Bigr)=\varphi^{-1}(\operatorname{Nil}B)=\sqrt{\operatorname{Ker}\varphi}, \]
because $\varphi(a)$ is nilpotent iff $\varphi(a^m)=0$ for some $m\geq1$ (Lemma~1.6: the intersection of all primes is the nilradical). Since $V(\sqrt J)=V(J)$ for every ideal $J$, we get $\overline{\operatorname{Im}f}=V(\operatorname{Ker}\varphi)$.

\emph{Remarks.} (i) The inclusion $\operatorname{Im}f\subseteq V(\operatorname{Ker}\varphi)$ is obvious, but equality can fail without the closure: for $\varphi:\mathbb{Z}\to\mathbb{Q}$ we have $\operatorname{Im}f=\{(0)\}$ while $V(\operatorname{Ker}\varphi)=\operatorname{Spec}\mathbb{Z}$. If $\varphi$ is surjective, then $\operatorname{Im}f=V(\operatorname{Ker}\varphi)$: this is Lemma~3.17, since $B\simeq A/\operatorname{Ker}\varphi$.

(ii) Let $\mathfrak{p}\in\operatorname{Spec}A$, $B=A_{\mathfrak{p}}$ and $\varphi:A\to A_{\mathfrak{p}}$ the localization map. A prime $\mathfrak{q}$ of $A$ is of the form $\varphi^{-1}(\mathfrak{Q})$ for a prime $\mathfrak{Q}$ of $A_{\mathfrak{p}}$ iff $\mathfrak{q}\cap(A\setminus\mathfrak{p})=\emptyset$, i.e.\ iff $\mathfrak{q}\subseteq\mathfrak{p}$ (the primes of $A_{\mathfrak{p}}$ correspond bijectively to these $\mathfrak{q}$, via $\mathfrak{q}A_{\mathfrak{p}}$). Hence
\[ \operatorname{Im}f=\{\mathfrak{q}\in\operatorname{Spec}A:\mathfrak{q}\subseteq\mathfrak{p}\}=\bigcap_{s\in A\setminus\mathfrak{p}}D(s), \]
the set of generizations of $\mathfrak{p}$. Moreover $\operatorname{Ker}\varphi=\{a\in A:sa=0\text{ for some }s\notin\mathfrak{p}\}$, and $\overline{\operatorname{Im}f}=V(\operatorname{Ker}\varphi)$ by (b). This set contains $\mathfrak{p}$ and hence the closure $V(\mathfrak{p})$ of $\{\mathfrak{p}\}$. Therefore $\operatorname{Im}f$ can be closed only if $V(\mathfrak{p})\subseteq\{\mathfrak{q}:\mathfrak{q}\subseteq\mathfrak{p}\}$, i.e.\ only if $\mathfrak{p}$ is a maximal ideal. In general $\operatorname{Im}f$ is not closed: for $A=\mathbb{Z}$ and $\mathfrak{p}=(p)$ the image is $\{(0),(p)\}$, which is not closed, with closure all of $\operatorname{Spec}\mathbb{Z}$ (here $\operatorname{Ker}\varphi=0$). If $A$ is a domain then $\operatorname{Ker}\varphi=0$ and $\operatorname{Im} f$ is dense in $\operatorname{Spec}A$ for every $\mathfrak{p}$. If $A$ is local with maximal ideal $\mathfrak{p}$, then $\varphi$ is an isomorphism and $\operatorname{Im}f=\operatorname{Spec}A$. \hfill$\square$

\exercise{3.4}{Let $X$ be a scheme and $f\in\mathcal{O}_X(X)$. Show that $U\mapsto f|_U\mathcal{O}_X(U)$ for every affine open subset $U$ defines a sheaf of ideals on $X$. We denote this sheaf by $f\mathcal{O}_X$. Show that $\operatorname{Supp}f\mathcal{O}_X$ is closed (see Exercise 2.5).}
\pf
Define, for every open $W\subseteq X$,
\[ \mathcal{I}(W)=\{s\in\mathcal{O}_X(W):\ s_x\in f_x\mathcal{O}_{X,x}\text{ for every }x\in W\}. \]
This is an $\mathcal{O}_X(W)$-submodule of $\mathcal{O}_X(W)$ (an ideal), stable under restriction. It is a sheaf: the condition on $s$ is a condition on the germs $s_x$, so if $s\in\mathcal{O}_X(W)$ restricts to elements of $\mathcal{I}(W_j)$ on an open covering $\{W_j\}$ of $W$, then $s\in\mathcal{I}(W)$; and the sections of $\mathcal{I}$ are sections of the sheaf $\mathcal{O}_X$, so the uniqueness and glueing conditions follow from those of $\mathcal{O}_X$ (a glued section of $\mathcal{O}_X$ has germs in $f_x\mathcal{O}_{X,x}$ by the same local property). Thus $\mathcal{I}$ is a sheaf of ideals of $\mathcal{O}_X$.

\emph{Claim: for every affine open $U$, $\mathcal{I}(U)=f|_U\mathcal{O}_X(U)$.} Write $U=\operatorname{Spec}A$ (we identify $U$ with $\operatorname{Spec}A$) and $a=f|_U\in A$. The inclusion $\supseteq$ is clear. Conversely let $s\in\mathcal{I}(U)$. For every $\mathfrak{p}\in U$ we have $s_{\mathfrak{p}}\in aA_{\mathfrak{p}}$ (Prop.~3.1(b)), i.e.\ $s/1=(b/t)a$ in $A_{\mathfrak{p}}$ for some $b\in A$, $t\in A\setminus\mathfrak{p}$, i.e.\ $t'(ts-ba)=0$ for some $t'\notin\mathfrak{p}$, so $t't\,s\in aA$ with $t't\notin\mathfrak{p}$. So the ideal $J=\{t\in A:ts\in aA\}$ is not contained in any prime $\mathfrak{p}$, hence $J=A$, $s\in aA$. This proves the claim.

Since the affine open subsets form a base of the topology and $\mathcal{I}$ is a sheaf, $\mathcal{I}$ is the unique sheaf of ideals whose sections over each affine open $U$ are $f|_U\mathcal{O}_X(U)$ (with the restriction maps of $\mathcal{O}_X$). This is the sheaf $f\mathcal{O}_X$ of the statement; its stalk at $x$ is $(f\mathcal{O}_X)_x=f_x\mathcal{O}_{X,x}$: indeed $f\in\mathcal{I}(X)$, so $f_x\mathcal{O}_{X,x}\subseteq\mathcal{I}_x$, and $\subseteq$ in the other direction holds by definition of $\mathcal{I}$.

\emph{Support.} $\operatorname{Supp}f\mathcal{O}_X=\{x\in X:f_x\mathcal{O}_{X,x}\neq0\}=\{x:f_x\neq0\}$. Its complement $\{x:f_x=0\}$ is open: if $f_x=0$, there is an open neighbourhood $V$ of $x$ with $f|_V=0$, hence $f_y=0$ for all $y\in V$. So the support is closed. (Contrast with Ex.~2.5: for a general sheaf the support need not be closed.) \hfill$\square$

\exercise{3.5}{Let $Y$ be a scheme that satisfies the conclusion of Proposition 3.25 for every affine scheme $X$. Show that $Y$ is affine.}
\pf
The hypothesis is that for every affine scheme $X$ the map $\rho_X:\operatorname{Mor}(X,Y)\to\operatorname{Hom}(\mathcal{O}_Y(Y),\mathcal{O}_X(X))$ is bijective. Put $A=\mathcal{O}_Y(Y)$ and $Z=\operatorname{Spec}A$, so $\mathcal{O}_Z(Z)=A$ (Prop.~3.1). We use (N) freely.

By Prop.~3.25 applied to the affine target $Z$ and the source $Y$, there is a unique morphism $u:Y\to Z$ with $\rho_Z(u)=u^{\#}(Z)=\operatorname{id}_A$. By the hypothesis applied to the affine scheme $Z$, there is a unique morphism $v:Z\to Y$ with $\rho_Y(v)=v^{\#}(Y)=\operatorname{id}_A\in\operatorname{Hom}(A,\mathcal{O}_Z(Z))$.

\emph{$u\circ v=\operatorname{id}_Z$.} By (N), $\rho_Z(u\circ v)=v^{\#}(Y)\circ\rho_Z(u)=\operatorname{id}_A=\rho_Z(\operatorname{id}_Z)$. As $\rho_Z:\operatorname{Mor}(Z,Z)\to\operatorname{Hom}(A,A)$ is injective (Lemma~3.23), $u\circ v=\operatorname{id}_Z$.

\emph{$v\circ u=\operatorname{id}_Y$.} Let $w:W\to Y$ be the inclusion of an affine open subset $W$ of $Y$. By (N),
\[ \rho_Y(v\circ u\circ w)=(u\circ w)^{\#}(Z)\circ\rho_Y(v)=\rho_Z(u\circ w)=w^{\#}(Y)\circ\rho_Z(u)=w^{\#}(Y)=\rho_Y(w), \]
where $w^{\#}(Y)$ is the restriction $A\to\mathcal{O}_Y(W)$ (Ex.~3.1). Since $W$ is affine, $\rho_Y:\operatorname{Mor}(W,Y)\to\operatorname{Hom}(A,\mathcal{O}_Y(W))$ is injective by hypothesis, so $v\circ u\circ w=w$. The affine open subsets $W$ cover $Y$, and two morphisms of schemes $Y\to Y$ that agree after composition with the inclusion of every member of an open covering are equal (they agree as maps, and the sheaf maps agree on a base; cf.\ Ex.~2.11). Hence $v\circ u=\operatorname{id}_Y$.

Therefore $u:Y\to\operatorname{Spec}A$ is an isomorphism and $Y$ is affine. \hfill$\square$

\exercise{3.6}{Generalize Example 3.32 to the case when $k$ is an arbitrary ring.}
\pf
Let $k$ be a ring, $I\subseteq k[T_1,\dots,T_n]$ an ideal, $C=k[T_1,\dots,T_n]/I$ and $X=\operatorname{Spec}C$, a $k$-scheme via $\iota:k\to C$ (Prop.~3.14). Let $X(k)$ be the set of sections of $X\to\operatorname{Spec}k$, i.e.\ of morphisms of $k$-schemes $\sigma:\operatorname{Spec}k\to X$ (Def.~3.28). Put
\[ Z=\{\alpha=(\alpha_1,\dots,\alpha_n)\in k^n:\ P(\alpha)=0\ \text{for all }P\in I\}. \]
\emph{Claim: there is a canonical bijection $\lambda:Z\to X(k)$.} (For a field $k$ this is Example~3.32; the point $\lambda(\alpha)\in X$ is the image of the section, namely the prime $\mathfrak{m}_\alpha=(T_1-\alpha_1,\dots,T_n-\alpha_n)/I$.)

Let $\operatorname{Hom}_{k\text{-alg}}(C,k)$ be the set of $k$-algebra homomorphisms. By Prop.~3.25 (target $X$ affine, source $\operatorname{Spec}k$), $\rho$ is a bijection $\operatorname{Mor}(\operatorname{Spec}k,X)\to\operatorname{Hom}_{\text{rings}}(C,k)$. A morphism $\sigma$ is a morphism of $k$-schemes iff $\pi\circ\sigma=\operatorname{id}_{\operatorname{Spec}k}$, where $\pi:X\to\operatorname{Spec}k$ is the structural morphism with $\rho(\pi)=\iota$. By (N), $\rho(\pi\circ\sigma)=\sigma^{\#}(X)\circ\rho(\pi)=\rho(\sigma)\circ\iota$, while $\rho(\operatorname{id})=\operatorname{id}_k$; as $\rho$ is bijective (Lemma~3.23), $\sigma$ is a section iff $\rho(\sigma)\circ\iota=\operatorname{id}_k$, i.e.\ iff $\rho(\sigma)$ is a $k$-algebra homomorphism $C\to k$. So
\[ X(k)\ \xrightarrow[\ \rho\ ]{\sim}\ \operatorname{Hom}_{k\text{-alg}}(C,k). \]
Finally, a $k$-algebra homomorphism $\psi:C\to k$ is the same as a $k$-algebra homomorphism $\tilde\psi:k[T_1,\dots,T_n]\to k$ vanishing on $I$. Such a $\tilde\psi$ is determined by $\alpha_i=\tilde\psi(T_i)$ and is the evaluation $P\mapsto P(\alpha)$, and any $\alpha\in k^n$ gives one. It vanishes on $I$ iff $P(\alpha)=0$ for all $P\in I$, i.e.\ iff $\alpha\in Z$. Hence $\alpha\mapsto(\text{evaluation at }\alpha)$ is a bijection from $Z$ onto $\operatorname{Hom}_{k\text{-alg}}(C,k)$, and composing with $\rho^{-1}$ we obtain the bijection $\lambda:Z\to X(k)$. \hfill$\square$

\exercise{3.7}{Let $X$ be a scheme over a field $k$. Let $\varphi:k[T_1,\dots,T_n]\to\mathcal{O}_X(X)$ be a homomorphism of $k$-algebras and $f:X\to\mathbb{A}^n_k$ the morphism induced by $\varphi$. Show that for any rational point $x\in X(k)$, via the identification $\mathbb{A}^n_k(k)=k^n$, we have $f(x)=(f_1(x),\dots,f_n(x))$, where $f_i=\varphi(T_i)$ and $f_i(x)$ is the image of $f_i$ in $k(x)=k$.}
\pf
Here $f$ is the morphism with $\rho(f)=f^{\#}(\mathbb{A}^n_k)=\varphi$ (Prop.~3.25). It is a morphism of $k$-schemes: with $\pi_X,\pi$ the structural morphisms of $X$ and $\mathbb{A}^n_k$, by (N) $\rho(\pi\circ f)=\varphi|_k=\rho(\pi_X)$ since $\varphi$ is a $k$-algebra map, so $\pi\circ f=\pi_X$ by injectivity of $\rho$ (Prop.~3.25).

Let $x\in X$ have $k(x)=k$ (Example~3.29) and $y=f(x)\in\mathbb{A}^n_k$. The local homomorphism $f^{\#}_x:\mathcal{O}_{\mathbb{A}^n_k,y}\to\mathcal{O}_{X,x}$ induces a field homomorphism $\bar f_x:k(y)\to k(x)=k$ which is a homomorphism of $k$-algebras (because $f$ is a $k$-morphism). Since $k(y)$ is a field containing $k$ (as the image of $k$ is $k$ itself, $k\to k(y)$ is injective), and $\bar f_x$ is a $k$-linear field embedding into $k$, we get $k(y)=k$ and $\bar f_x=\operatorname{id}_k$. So $y$ is a rational point of $\mathbb{A}^n_k$, and under Example~3.32 (with $I=0$) it corresponds to $\alpha=(\alpha_1,\dots,\alpha_n)$, where $\alpha_i$ is the image of $T_i$ in $k(y)=k$ (i.e.\ $y=\mathfrak{m}_\alpha$).

Now, for the global section $T_i$ of $\mathbb{A}^n_k$ we have $f^{\#}_x((T_i)_y)=(f^{\#}(\mathbb{A}^n_k)(T_i))_x=(\varphi(T_i))_x=(f_i)_x$. Passing to residue fields, $\bar f_x$ sends the class $\alpha_i$ of $(T_i)_y$ in $k(y)$ to the class $f_i(x)$ of $(f_i)_x$ in $k(x)$. Since $\bar f_x=\operatorname{id}_k$, we get $\alpha_i=f_i(x)$, i.e.\ $f(x)=(f_1(x),\dots,f_n(x))$. \hfill$\square$

\exercise{3.8}{Let $X$ be a scheme over a ring $A$. Let $f_0,\dots,f_n\in\mathcal{O}_X(X)$ be such that the $f_{i,x}$ generate the unit ideal of $\mathcal{O}_{X,x}$ for every $x\in X$. Show that $X$ is the union of the $X_{f_i}$ (see Definition 3.11), and that we have a morphism $f:X\to\mathbb{P}^n_A$ such that $f^{-1}(D_+(T_i))=X_{f_i}$ and that $f|_{X_{f_i}}$ is induced by the homomorphism $A[T_i^{-1}T_j]_j\to\mathcal{O}_X(X_{f_i})$ given by $T_i^{-1}T_j\mapsto f_i^{-1}f_j$. If $A=k$ is a field and $x\in X(k)$, determine $f(x)$ as in the preceding exercise.}
\pf
Write $B=A[T_0,\dots,T_n]$, $\mathbb{P}^n_A=\operatorname{Proj}B$ (Ex.~3.39), $U_i=X_{f_i}$ (open by Prop.~3.12), and $D_i=D_+(T_i)=\operatorname{Spec}B_{(T_i)}$ with $B_{(T_i)}=A[T_i^{-1}T_j]_j$.

\emph{A general fact.} (*) Let $g:U\to\operatorname{Spec}C$ be a morphism of schemes and $h\in C$. Then $g^{-1}(D(h))=U_{\rho(g)(h)}$. Indeed, for $x\in U$ and $\mathfrak{p}=g(x)$: $x\in g^{-1}(D(h))$ iff $h\notin\mathfrak{p}$ iff the image of $h$ in $C_{\mathfrak{p}}$ is a unit. As $g^{\#}_x:C_{\mathfrak{p}}\to\mathcal{O}_{U,x}$ is a local homomorphism, an element of $C_{\mathfrak{p}}$ is a unit iff its image in $\mathcal{O}_{U,x}$ is a unit (a non-unit lies in the maximal ideal, which maps into the maximal ideal). The image of $h$ in $\mathcal{O}_{U,x}$ is $(\rho(g)(h))_x$, whence (*).

\emph{$X=\bigcup U_i$.} Let $x\in X$. The ideal of the local ring $\mathcal{O}_{X,x}$ generated by the $f_{i,x}$ is the unit ideal, so not all $f_{i,x}$ lie in the maximal ideal. Hence some $f_{i,x}$ is a unit, i.e.\ $x\in U_i$. Moreover, as in the proof of Prop.~3.12, $f_i|_{U_i}$ is invertible in $\mathcal{O}_X(U_i)$ (glue local inverses); put $f_i^{-1}$ for its inverse.

\emph{Construction of $f$.} Note that $B_{(T_i)}$ is the polynomial ring over $A$ in the $n$ variables $S_j=T_i^{-1}T_j$, $j\neq i$ (the elements of degree $0$ in $A[T_0,\dots,T_n,T_i^{-1}]$ are the Laurent polynomials in $T_j$ with total degree $0$ and non-negative exponents for $j\neq i$; these are exactly the polynomials in the $S_j$). Since $\mathcal{O}_X(U_i)$ is an $A$-algebra (Cor.~3.26), there is a unique $A$-algebra homomorphism
\[ \psi_i:B_{(T_i)}\to\mathcal{O}_X(U_i),\qquad T_i^{-1}T_j\mapsto f_i^{-1}f_j\quad(j\neq i), \]
and then also $T_i^{-1}T_i=1\mapsto f_i^{-1}f_i=1$. By Prop.~3.25 (the target $D_i$ is affine) there is a unique morphism $g_i:U_i\to D_i\subseteq\mathbb{P}^n_A$ with $\rho(g_i)=\psi_i$; it is a morphism of $A$-schemes, by the argument of Ex.~3.7 (use (N) and injectivity of $\rho$ into $\operatorname{Spec}A$).

\emph{Compatibility.} Put $U_{ij}=U_i\cap U_j$ and $R_{ij}=B_{(T_iT_j)}=\mathcal{O}_{\mathbb{P}^n_A}(D_i\cap D_j)$. By Lemma~3.36(b),(c) with $g=T_iT_j$, $f=T_i$ (so $\alpha=T_j/T_i$), $D_i\cap D_j=D(T_i^{-1}T_j)\subseteq D_i$, and $R_{ij}=(B_{(T_i)})_{T_i^{-1}T_j}$, and similarly with $i,j$ exchanged. By (*), $g_i^{-1}(D_i\cap D_j)=(U_i)_{\psi_i(T_i^{-1}T_j)}=(U_i)_{f_i^{-1}f_j}=U_i\cap U_j$ (as $f_i$ is a unit on $U_i$, $f_i^{-1}f_j$ is a unit at $x\in U_i$ iff $f_{j,x}$ is). So $g_i$ and $g_j$ restrict to morphisms $U_{ij}\to\operatorname{Spec}R_{ij}$. Consider the ring homomorphism $\Theta:A[T_0,\dots,T_n]\to\mathcal{O}_X(U_{ij})[u,u^{-1}]$ given by $T_k\mapsto f_ku$. It sends $T_iT_j$ to the unit $f_if_ju^2$ (units on $U_{ij}$), so it extends to $B_{T_iT_j}$, and it maps the degree-$0$ part $R_{ij}$ into $\mathcal{O}_X(U_{ij})$ (the elements of degree $0$ go to elements of $u$-degree $0$). We obtain an $A$-algebra homomorphism $\theta_{ij}:R_{ij}\to\mathcal{O}_X(U_{ij})$ with $\theta_{ij}(T_k/T_i)=f_k/f_i$ and $\theta_{ij}(T_k/T_j)=f_k/f_j$ for all $k$. The restriction of $\theta_{ij}$ to $B_{(T_i)}$ is $\psi_i$ (followed by restriction to $U_{ij}$), and to $B_{(T_j)}$ it is $\psi_j$. Now $\rho(g_i|_{U_{ij}})$ also restricts to $\psi_i$ on $B_{(T_i)}$, and a homomorphism out of the localization $R_{ij}=(B_{(T_i)})_{T_i^{-1}T_j}$ is determined by its restriction to $B_{(T_i)}$; so $\rho(g_i|_{U_{ij}})=\theta_{ij}$, and likewise $\rho(g_j|_{U_{ij}})=\theta_{ij}$. By the injectivity in Prop.~3.25 (the target $\operatorname{Spec}R_{ij}$ is affine), $g_i|_{U_{ij}}=g_j|_{U_{ij}}$.

By glueing (Ex.~2.11) there is a morphism $f:X\to\mathbb{P}^n_A$ with $f|_{U_i}=g_i$, a morphism of $A$-schemes. Finally $f^{-1}(D_i)\cap U_j=g_j^{-1}(D_j\cap D_i)=U_i\cap U_j$ for all $j$ (computed above, also for $j=i$), so $f^{-1}(D_i)=\bigcup_j(U_i\cap U_j)=U_i=X_{f_i}$. By construction $f|_{X_{f_i}}$ is the morphism induced by $T_i^{-1}T_j\mapsto f_i^{-1}f_j$.

\emph{The case $A=k$ a field, $x\in X(k)$.} Then $f_j(x)\in k(x)=k$ is the image of $f_{j,x}$. Because $f_{i,x}$ is a unit iff its image $f_i(x)$ in the residue field is $\neq0$, the hypothesis says that $f_0(x),\dots,f_n(x)$ are not all $0$. Choose $i$ with $f_i(x)\neq0$; then $x\in U_i$ and $x\in U_i(k)$ (Remark~3.31). Apply Ex.~3.7 to $U_i$, the $k$-algebra homomorphism $\psi_i$, and the identification $D_i\simeq\mathbb{A}^n_k$: the point $f(x)=g_i(x)\in D_i$ is the rational point with coordinates $T_i^{-1}T_j\mapsto\psi_i(T_i^{-1}T_j)(x)=f_j(x)/f_i(x)$. Comparing with Lemma~3.43 (whose proof shows that the point of $\mathbb{P}^n_k(k)$ with homogeneous coordinates $(\beta_0,\dots,\beta_n)$, $\beta_i\neq0$, lies in $D_+(T_i)$ and corresponds there to the ideal generated by the $T_i^{-1}T_j-\beta_i^{-1}\beta_j$), we conclude
\[ f(x)=(f_0(x):f_1(x):\cdots:f_n(x))\in\mathbb{P}(k^{n+1})=\mathbb{P}^n_k(k). \]
This does not depend on the choice of $i$ since proportional coordinates give the same point. \hfill$\square$

\exercise{3.9}{Let $B$ be a graded ring. Let $f\in B$ be non-nilpotent, homogeneous of degree $0$. Show that $B_f$ possesses a natural grading, and that $D_+(f)\simeq\operatorname{Proj}B_f$.}
\pf
\emph{Grading.} Since $f\in B_0$, localization commutes with direct sums: $B_f=\bigoplus_{d\geq0}(B_d)_f$, where $(B_d)_f=\{b/f^m:b\in B_d,m\geq0\}$ is the localization of the $B_0$-module $B_d$. Put $(B_f)_d=(B_d)_f$ for $d\geq0$. Then $(B_f)_d(B_f)_e\subseteq(B_f)_{d+e}$ and $1\in(B_f)_0$, so $B_f$ is a graded ring and $B\to B_f$ is a homomorphism of graded rings; we have $(B_f)_0=(B_0)_f$. The hypothesis that $f$ is not nilpotent just says that $B_f\neq0$; the result below holds without it (if $f$ is nilpotent, both sides are empty).

\emph{The isomorphism.} Let $\iota:B\to B_f$. It is a graded homomorphism with $r=1$ in Lemma~3.40, and $M=\iota(B_+)B_f=(B_f)_+$, so $V_+(M)=\emptyset$. By Lemma~3.40 we have a morphism $\tilde f:\operatorname{Proj}B_f\to\operatorname{Proj}B$, $\mathfrak{q}\mapsto\iota^{-1}(\mathfrak{q})$, such that for every homogeneous $h\in B_+$, $\tilde f^{-1}(D_+(h))=D_+(h/1)$ and $\tilde f$ restricted to it is the morphism of affine schemes
\[ \operatorname{Spec}(B_f)_{(h/1)}\to\operatorname{Spec}B_{(h)} \]
induced by $B_{(h)}\to(B_f)_{(h)}$ (we write $h$ for $h/1$). Since $f/1$ is a unit of $B_f$, $f\notin\iota^{-1}(\mathfrak{q})$, so $\tilde f$ maps into $D_+(f)$.

Let $h\in B_+$ be homogeneous of degree $d>0$. Because $f\in B_0$ we have $f/1\in B_{(h)}$. The natural map $(B_{(h)})_f\to(B_f)_{(h)}$, $\ (a/h^k)/f^m\mapsto(a/f^m)/h^k$ ($a\in B_{kd}$), is an isomorphism: it is surjective because an element of $(B_f)_{(h)}$ has the form $(a/f^m)/h^k$ with $a\in B_{kd}$; it is injective because $a/(f^mh^k)=0$ in $B_{fh}$ means $(fh)^sa=0$ in $B$ for some $s$, i.e.\ $a/h^k$ is killed by a power of $f$ in $B_{(h)}$. Therefore $B_{(h)}\to(B_f)_{(h)}$ is the localization at $f$, and by Lemma~3.7 and Example~3.15 $\operatorname{Spec}(B_f)_{(h)}\to\operatorname{Spec}B_{(h)}$ is an isomorphism onto the open subscheme $D(f)\subseteq\operatorname{Spec}B_{(h)}$. By Lemma~3.36(b) (applied to $D_+(fh)\subseteq D_+(h)$: here $r=d=\deg(fh)$, so $\alpha=(fh)^dh^{-d}=f^d\in B_{(h)}$ and $D(\alpha)=D(f)$), $D(f)$ corresponds to $D_+(fh)=D_+(f)\cap D_+(h)$. Thus $\tilde f$ induces an isomorphism of schemes
\[ \tilde f^{-1}(D_+(h))\ \xrightarrow{\ \sim\ }\ D_+(f)\cap D_+(h). \]
Since $\tilde f$ lands in $D_+(f)$, we have $\tilde f^{-1}(D_+(f)\cap D_+(h))=\tilde f^{-1}(D_+(h))$. The open sets $D_+(f)\cap D_+(h)$, $h\in B_+$ homogeneous, cover $D_+(f)$ (the $D_+(h)$ cover $\operatorname{Proj}B$). So $\tilde f:\operatorname{Proj}B_f\to D_+(f)$ is an isomorphism over each member of an open covering of the target; hence it is an isomorphism (the local inverses glue). Thus $D_+(f)\simeq\operatorname{Proj}B_f$. \hfill$\square$

\exercise{3.10}{Let $A$ be a ring, $X=\mathbb{P}^n_A$. Show that $\mathcal{O}_X(X)=A$ (use the covering $\mathcal{U}=\{D_+(T_i)\}_i$ of $X$ and the complex associated to $\mathcal{U}$ and $\mathcal{O}_X$ as in Lemma 2.7). Deduce from this that $X$ is affine if and only if $n=0$.}
\pf
Let $B=A[T_0,\dots,T_n]$, $U_i=D_+(T_i)$, $U_{ij}=U_i\cap U_j=D_+(T_iT_j)$. By Prop.~3.38, $\mathcal{O}_X(U_i)=B_{(T_i)}$ and $\mathcal{O}_X(U_{ij})=B_{(T_iT_j)}$, and the restriction $B_{(T_i)}\to B_{(T_iT_j)}$ is the map of Lemma~3.36(c); with $g=T_iT_j=T_i\cdot T_j$ it sends $aT_i^{-N}$ to $(aT_j^N)(T_iT_j)^{-N}=aT_i^{-N}$, i.e.\ it is the inclusion of subrings of $B_{T_iT_j}$.

By Lemma~2.7, since $\mathcal{O}_X$ is a sheaf the sequence
\[ 0\to\mathcal{O}_X(X)\to\bigoplus_i\mathcal{O}_X(U_i)\to\bigoplus_{i,j}\mathcal{O}_X(U_{ij}),\qquad (s_i)\mapsto(s_i|_{U_{ij}}-s_j|_{U_{ij}}), \]
is exact. So $\mathcal{O}_X(X)=\{(s_i)_i:s_i\in B_{(T_i)},\ s_i=s_j\text{ in }B_{(T_iT_j)}\ \text{for all }i,j\}$.

All rings involved are free $A$-modules with monomial bases. Let $M_0=\{a\in\mathbb{Z}^{n+1}:a_0+\dots+a_n=0\}$ and $T^a=\prod T_k^{a_k}$. Then $B_{(T_i)}$ has $A$-basis $\{T^a:a\in S_i\}$ with $S_i=\{a\in M_0:a_k\geq0\ \text{for }k\neq i\}$ (the degree $0$ part of $B_{T_i}$), and $B_{(T_iT_j)}$ has basis $\{T^a:a\in S_{ij}\}$, $S_{ij}=\{a\in M_0:a_k\geq0\ \text{for }k\neq i,j\}\supseteq S_i\cup S_j$. Everything lies in the free $A$-module $A[M_0]$ with basis $\{T^a\}_{a\in M_0}$ (inside the Laurent polynomial ring $A[T_0^{\pm1},\dots,T_n^{\pm1}]$), and the restriction maps are inclusions. Write $s_i=\sum_{a\in S_i}c_i(a)T^a$. The condition $s_i=s_j$ in $A[M_0]$ says $c_i(a)=c_j(a)$ for all $a\in M_0$ (with $c_i(a)=0$ for $a\notin S_i$). So there is one family $c:M_0\to A$ with $c(a)=0$ unless $a\in\bigcap_iS_i$, and $s_i=\sum c(a)T^a$ for all $i$. Now $a\in\bigcap_iS_i$ means $a_k\geq0$ for all $k$ (for each $k$ there is an $i\neq k$ if $n\geq1$, and for $n=0$, $M_0=\{0\}$), and together with $\sum a_k=0$ this forces $a=0$. Hence $s_i=c(0)\in A$ for all $i$, and $\mathcal{O}_X(X)=A\cdot1$. Since $A\to B_{(T_0)}=A[T_0^{-1}T_j]$ is injective, the structure homomorphism $A\to\mathcal{O}_X(X)$ is an isomorphism.

\emph{Affineness.} If $n=0$, $\mathbb{P}^0_A=\operatorname{Spec}A[T_0]_{(T_0)}=\operatorname{Spec}A$ is affine. For $A=0$ we have $\mathbb{P}^n_A=\emptyset$, which is affine for every $n$, so the ``only if'' part requires $A\neq0$. Assume $A\neq0$, $n\geq1$ and, by contradiction, $X$ affine. Then $X\simeq\operatorname{Spec}C$ with $C=\mathcal{O}_X(X)\simeq A$, and the structural morphism $\pi:X\to\operatorname{Spec}A$ corresponds by Lemma~3.23 to the structure homomorphism $\iota:A\to\mathcal{O}_X(X)$, which is an isomorphism by what we just showed. Thus $\pi$ is the composition of an isomorphism $X\simeq\operatorname{Spec}\mathcal{O}_X(X)$ and $\operatorname{Spec}\iota$ (a homeomorphism); in particular $\pi$ is injective on points.

Restrict to $U_0=\operatorname{Spec}A[S_1,\dots,S_n]$, $S_j=T_0^{-1}T_j$. By definition of the $A$-scheme structure (Prop.~3.38), $\pi|_{U_0}$ is the morphism attached to the inclusion $A\to A[S_1,\dots,S_n]$, so on points $\mathfrak{q}\mapsto\mathfrak{q}\cap A$. Take a prime $\mathfrak{p}$ of $A$ ($A\neq0$) and put
\[ \mathfrak{q}_c=\mathfrak{p}A[S_1,\dots,S_n]+(S_1-c,S_2,\dots,S_n),\qquad c=0,1. \]
Then $A[S]/\mathfrak{q}_c\simeq A/\mathfrak{p}$ is a domain, so $\mathfrak{q}_c$ is prime, and $\mathfrak{q}_c\cap A=\mathfrak{p}$ for $c=0,1$. But $\mathfrak{q}_0\neq\mathfrak{q}_1$, since $S_1\in\mathfrak{q}_0$ whereas $S_1\notin\mathfrak{q}_1$ (otherwise $1=S_1-(S_1-1)\in\mathfrak{q}_1$). So $\pi$ is not injective, a contradiction. Hence $\mathbb{P}^n_A$ is not affine for $n\geq1$ and $A\neq0$. \hfill$\square$

\exercise{3.11}{Let $B=\bigoplus_{d\geq0}B_d$ be a graded ring. Let $e\geq1$ be an integer.

(a) Let us denote the graded ring $\bigoplus_{d\geq0}B_{de}$ by $C$ (so $C_d:=B_{de}$). Show that $\operatorname{Proj}B\simeq\operatorname{Proj}C$.

(b) Let us suppose that for every $d\geq1$, $B_{ed}$ is generated by $B_e^d$, and that $B_e$ is finitely generated over a ring $A$. Show that $\operatorname{Proj}B$ is a projective scheme over $A$.}
\pf
(a) The inclusion $\iota:C\hookrightarrow B$ is a homomorphism of graded rings in the sense of Lemma~3.40, with $r=e$ (it maps $C_d=B_{de}$ into $B_{ed}$). Let $M=\iota(C_+)B$. We claim $V_+(M)=\emptyset$. If $\mathfrak{p}\in\operatorname{Proj}B$ contains $M$, it contains every $B_{de}$, $d\geq1$. For a homogeneous $b\in B_k$, $k\geq1$, we have $b^e\in B_{ke}\subseteq\mathfrak{p}$, so $b\in\mathfrak{p}$, whence $B_+\subseteq\mathfrak{p}$, contradicting $\mathfrak{p}\in\operatorname{Proj}B$. By Lemma~3.40 we get a morphism $\phi:\operatorname{Proj}B\to\operatorname{Proj}C$, $\mathfrak{p}\mapsto\mathfrak{p}\cap C$, such that for every homogeneous $h\in C_+$ (say $h\in C_d=B_{de}$, $d\geq1$), $\phi^{-1}(D_+(h))=D_+(h)$ (here the right side is taken in $\operatorname{Proj}B$, $h$ viewed as an element of $B_{de}\subseteq B_+$), and $\phi$ restricted to it is induced by the homomorphism $\gamma:C_{(h)}\to B_{(h)}$, $c/h^N\mapsto c/h^N$.

We claim $\gamma$ is bijective. Elements of $B_{(h)}$ are $b/h^N$ with $b\in B_{Nde}=C_{Nd}$, so $b\in C$ and $\gamma$ is surjective. If $c/h^N\in C_{(h)}$ maps to $0$ in $B_h$, then $h^sc=0$ in $B$ for some $s$; this is an equation in $C$, so $c/h^N=0$ already in $C_h$. Hence $\gamma$ is an isomorphism, and $\phi$ restricts to an isomorphism of schemes $D_+(h)\to D_+(h)$ (the first in $\operatorname{Proj}B\simeq\operatorname{Spec}B_{(h)}$, the second in $\operatorname{Proj}C\simeq\operatorname{Spec}C_{(h)}$, by Prop.~3.38). The sets $D_+(h)$, $h\in C_+$ homogeneous, cover $\operatorname{Proj}C$. So $\phi$ is an isomorphism over each member of an open covering of the target, hence an isomorphism, and $\operatorname{Proj}B\simeq\operatorname{Proj}C$. This is an isomorphism of $A$-schemes for any ring homomorphism $A\to B_0=C_0$.

(b) Here $B$ is a graded $A$-algebra (so $A\to B_0$) and $\operatorname{Proj}B$ is an $A$-scheme. By (a), $\operatorname{Proj}B\simeq\operatorname{Proj}C$ as $A$-schemes, with $C_d=B_{de}$. Let $x_0,\dots,x_n\in B_e=C_1$ generate $B_e$ as an $A$-module (if $B_e=0$ take $n=0$, $x_0=0$). Let $R=A[T_0,\dots,T_n]$ with its usual grading and $\varphi:R\to C$ the $A$-algebra homomorphism with $T_i\mapsto x_i\in C_1$; it is a graded homomorphism ($r=1$). For $d\geq1$, $\varphi(R_d)$ is the $A$-span of the monomials of degree $d$ in the $x_i$, which is the $A$-span of all products of $d$ elements of $B_e$ (expand each factor in the $x_i$), and this is $B_{ed}=C_d$ by hypothesis. Thus
\[ \varphi(R_d)=C_d\quad\text{for all }d\geq1. \]
Hence $\varphi(R_+)C=C_+$, and Lemma~3.40 gives a morphism of $A$-schemes $g:\operatorname{Proj}C\to\operatorname{Proj}R=\mathbb{P}^n_A$ (the domain is everything because $V_+(\varphi(R_+)C)=V_+(C_+)=\emptyset$), with $g^{-1}(D_+(T_i))=D_+(x_i)$ and $g|_{D_+(x_i)}$ the morphism of affine schemes $\operatorname{Spec}C_{(x_i)}\to\operatorname{Spec}R_{(T_i)}$ induced by $\varphi_i:R_{(T_i)}\to C_{(x_i)}$, $P/T_i^N\mapsto\varphi(P)/x_i^N$. (If $x_i=0$ then $C_{(x_i)}=0$ and $D_+(x_i)=\emptyset$.)

\emph{$\varphi_i$ is surjective.} An element of $C_{(x_i)}$ has the form $c/x_i^N$ with $c\in C_N$. If $N\geq1$, then $c=\varphi(P)$ with $P\in R_N$ and $c/x_i^N=\varphi_i(P/T_i^N)$. If $N=0$, then $c\in C_0$, and $c/1=(cx_i)/x_i$ with $cx_i\in C_1=\varphi(R_1)$, say $cx_i=\varphi(P)$, so $c/1=\varphi_i(P/T_i)$. Hence $\varphi_i$ is surjective and, by Lemma~3.17, $g|_{D_+(x_i)}:D_+(x_i)\to D_+(T_i)$ is a closed immersion.

Since the $D_+(T_i)$ cover $\mathbb{P}^n_A$ and $g^{-1}(D_+(T_i))\to D_+(T_i)$ is a closed immersion for each $i$, $g$ is a closed immersion. (Being a closed immersion is local on the target: $g$ is injective, $g(\operatorname{Proj}C)\cap D_+(T_i)$ is closed in $D_+(T_i)$ for all $i$, so $g(\operatorname{Proj}C)$ is closed, $g$ is a homeomorphism onto its image because this holds over each $D_+(T_i)$, and surjectivity of $g^{\#}$ is checked on stalks, i.e.\ locally on the target (Ex.~2.7).) Therefore $\operatorname{Proj}B\simeq\operatorname{Proj}C$ is isomorphic to a closed subscheme of $\mathbb{P}^n_A$, i.e.\ it is a projective scheme over $A$ (Def.~3.42). \hfill$\square$

\emph{Notes for this part.} Results of the text are cited by number (``Prop.~3.25'', \dots). For a morphism $f:X\to Y$ of schemes and $V\subseteq Y$ open, $f^\#(V):\mathcal{O}_Y(V)\to\mathcal{O}_X(f^{-1}V)$ denotes the induced ring homomorphism, and $\rho:\operatorname{Mor}(X,Y)\to\operatorname{Hom}(\mathcal{O}_Y(Y),\mathcal{O}_X(X))$ is the bijection of Prop.~3.25 when $Y$ is affine. For $g:Z\to X$ the diagram (3.4) says that $\rho(f\circ g)=g^\#(X)\circ\rho(f)$. Ringed spaces are locally ringed and morphisms of ringed spaces have local stalk maps (Def.~2.19, 2.20). Ex.~3.13 uses the Veronese statement of Ex.~3.11(a), and its argument is included. Ex.~3.16 is used in Ex.~3.17(b). Ex.~3.19 uses, as black boxes, that $\mathcal{O}_K$ is a Dedekind domain and that the class group is finite. Ex.~3.21 uses Ex.~3.20.

\exercise{3.12}{Let $B$ be a Noetherian graded ring.
(a) Show that for any homogeneous $f\in B_+$, $B_{(f)}$ is Noetherian.
(b) Show that $\operatorname{Proj}B$ is Noetherian.}
\pf
Here $B=\bigoplus_{d\ge0}B_d$.

(a) Let $r=\deg f\ge1$. The localization $B_f$ is a $\mathbb{Z}$-graded ring: $(B_f)_n$ consists of the fractions $b/f^N$ with $b$ homogeneous of degree $Nr+n$. Then $B_{(f)}=(B_f)_0$. Let $I$ be an ideal of $B_{(f)}$. The ideal $IB_f$ of the Noetherian ring $B_f$ is finitely generated. Each generator is a finite $B_f$-combination of elements of $I$, so we can choose generators $a_1,\dots,a_m\in I$ of $IB_f$. We claim $I=(a_1,\dots,a_m)B_{(f)}$. Let $x\in I$. Then $x=\sum_i c_ia_i$ with $c_i\in B_f$. Taking degree-$0$ components and using $\deg a_i=0$ gives $x=\sum_i (c_i)_0a_i$ with $(c_i)_0\in B_{(f)}$. So $I$ is finitely generated, and $B_{(f)}$ is Noetherian.

(b) The ideal $B_+$ is homogeneous, and $B$ is Noetherian. So $B_+$ is generated by finitely many homogeneous elements $f_1,\dots,f_n$ of positive degree (it is generated by its homogeneous elements, and finitely many of them suffice). No point of $\operatorname{Proj}B$ contains $B_+$, so $\emptyset=V_+(B_+)=\bigcap_iV_+(f_i)$, that is $\operatorname{Proj}B=\bigcup_iD_+(f_i)$ (as in the text before Lemma~3.36). By Prop.~3.38, each $D_+(f_i)$ is an affine open subscheme isomorphic to $\operatorname{Spec}B_{(f_i)}$, and $B_{(f_i)}$ is Noetherian by (a). Hence $\operatorname{Proj}B$ is a finite union of affine opens with Noetherian rings, i.e.\ Noetherian (Def.~3.45). If $B_+=0$ the union is empty and $\operatorname{Proj}B=\emptyset$ is trivially Noetherian. \hfill$\square$

\exercise{3.13}{Let $B=A[X,Y,Z]$ be a polynomial ring. Let $B_d$ be the sub-$A$-module of $B$ generated by the elements of the form $X^aY^bZ^c$ with $a+2b+3c=d$.
(a) Show that the $B_d$ induce a grading on $B$.
(b) Determine $B_6$ and show that $B_{6d}$ is generated by the elements of $B_6^{\,d}$. Deduce from this that $\operatorname{Proj}B$ is isomorphic to a closed subscheme of $\mathbb{P}^6_A$.}
\pf
(a) The monomials $X^aY^bZ^c$ form an $A$-basis of $B$, and each monomial has exactly one weight $a+2b+3c$. So $B=\bigoplus_{d\ge0}B_d$ as $A$-modules, with $B_d$ free on the monomials of weight $d$. Weights add under multiplication of monomials, so $B_dB_e\subseteq B_{d+e}$. Also $1\in B_0=A$. Hence this is a grading of the $A$-algebra $B$ (with $B_0=A$).

(b) The monomials of weight $6$ are: for $c=0$, $X^6,\ X^4Y,\ X^2Y^2,\ Y^3$; for $c=1$ (so $a+2b=3$), $X^3Z,\ XYZ$; for $c=2$ (so $a=b=0$), $Z^2$. Thus $B_6$ is free over $A$ with basis
\[ m_0=X^6,\ m_1=X^4Y,\ m_2=X^2Y^2,\ m_3=Y^3,\ m_4=X^3Z,\ m_5=XYZ,\ m_6=Z^2. \]

\emph{Claim: for $d\ge1$, $B_{6d}$ is generated by the products of $d$ elements of $B_6$.} The products $B_6^{\,d}$ lie in $B_{6d}$. It suffices to show that every monomial $\mu=X^aY^bZ^c$ of weight $N=6d\ge6$ is divisible by one of $m_0,\dots,m_6$; then $\mu=m_i\mu'$ with $\mu'$ of weight $6(d-1)$, and induction on $d$ (the case $d=1$ being trivial) proves the claim. Note that $N$ is even.
\begin{itemize}
\item If $c\ge2$, then $Z^2=m_6$ divides $\mu$.
\item If $c=1$, then $a+2b=N-3$ is odd, so $a$ is odd and $a\ge1$. If $a\ge3$, $m_4=X^3Z$ divides $\mu$. If $a=1$, then $b=(N-4)/2\ge1$ and $m_5=XYZ$ divides $\mu$.
\item If $c=0$, then $a+2b=N$. If $b\ge3$, $m_3=Y^3$ divides $\mu$. If $b=2$, then $a=N-4\ge2$ and $m_2$ divides $\mu$. If $b=1$, then $a=N-2\ge4$ and $m_1$ divides $\mu$. If $b=0$, then $a=N\ge6$ and $m_0$ divides $\mu$.
\end{itemize}
This proves the claim.

\emph{Veronese subring.} Let $C=\bigoplus_{d\ge0}C_d$ with $C_d=B_{6d}$. This is a graded subring of $B$ (as $B_{6d}B_{6d'}\subseteq B_{6(d+d')}$), and $C_0=A$. We show that $\operatorname{Proj}B\simeq\operatorname{Proj}C$ (this is Ex.~3.11(a) for $e=6$; the argument works for any $e\ge1$). Let $\varphi:C\to B$ be the inclusion; it satisfies $\varphi(C_d)\subseteq B_{6d}$. Let $M=\varphi(C_+)B$ be the ideal of Lemma~3.40. Let $\mathfrak p\in\operatorname{Proj}B$ be a prime containing $M$. For every homogeneous $g\in B_+$ we have $g^6\in B_{6\deg g}\subseteq C_+\subseteq M\subseteq\mathfrak p$, so $g\in\mathfrak p$. Hence $B_+\subseteq\mathfrak p$, which is excluded in $\operatorname{Proj}B$. Therefore $V_+(M)=\emptyset$, and Lemma~3.40 gives a morphism of $A$-schemes $\pi:\operatorname{Proj}B\to\operatorname{Proj}C$, $\mathfrak p\mapsto\mathfrak p\cap C$, such that for every homogeneous $h\in C_+$ we have $\pi^{-1}(D_+(h))=D_+(h)$ and $\pi|_{D_+(h)}$ is the morphism of affine schemes induced by $C_{(h)}\to B_{(h)}$. If $h\in C_{\delta}=B_{6\delta}$, then
\[ C_{(h)}=\{a/h^N:\ a\in C_{N\delta}=B_{6N\delta}\}=\{a/h^N:\ a\in B_{6N\delta}\}=B_{(h)} . \]
More precisely, $C_h\to B_h$ is injective (localization is exact), and the map $C_{(h)}\to B_{(h)}$ is injective and surjective by the displayed equalities. So $\pi$ restricts to an isomorphism $D_+(h)\to D_+(h)$ over each member of the open cover $\{D_+(h)\}_h$ of $\operatorname{Proj}C$. Being an isomorphism is local on the target, so $\pi$ is an isomorphism.

\emph{Conclusion.} By the claim, $C_d=C_1^{\,d}$ for $d\ge1$, so $C$ is generated as an $A$-algebra by $C_1=B_6$. Hence the graded $A$-algebra homomorphism $A[T_0,\dots,T_6]\to C$, $T_i\mapsto m_i$, is surjective. Its kernel $I$ is a homogeneous ideal, and $C\simeq A[T_0,\dots,T_6]/I$. By Lemma~3.41, $\operatorname{Proj}C$ is isomorphic to a closed subscheme of $\mathbb{P}^6_A$. So is $\operatorname{Proj}B\simeq\operatorname{Proj}C$. \hfill$\square$

\exercise{3.14}{We say that a scheme $X$ is quasi-compact if it is quasi-compact as a topological space (Exercise 1.6). Show that a scheme is quasi-compact if and only if it is a finite union of affine schemes.}
\pf
($\Rightarrow$) Cover $X$ by affine open subsets. By quasi-compactness a finite subfamily already covers $X$.

($\Leftarrow$) Suppose $X=U_1\cup\dots\cup U_n$ with $U_i=\operatorname{Spec}A_i$ open. Each $\operatorname{Spec}A_i$ is quasi-compact (Ex.~1.6; in short, if $\operatorname{Spec}A=\bigcup_jD(g_j)$ then $1\in(g_j)_j$, so $1=\sum_{j\in F}c_jg_j$ for a finite set $F$, and $\operatorname{Spec}A=\bigcup_{j\in F}D(g_j)$; since the $D(g)$ form a base, every open cover has a finite subcover). Now let $X=\bigcup_{\lambda}W_\lambda$ be an open cover. For each $i$, finitely many $W_\lambda$ cover $U_i$. Taking all of them for $i=1,\dots,n$ gives a finite subcover of $X$. \hfill$\square$

\exercise{3.15}{Let $X$ be a quasi-compact scheme, $A=\mathcal{O}_X(X)$. Consider the morphism $f:X\to\operatorname{Spec}A$ induced by the identity on $A$ (Prop.~3.25). Show that $f(X)$ is dense in $\operatorname{Spec}A$.}
\pf
We prove a slightly more general statement, which is also used in Ex.~3.17(e).

\textbf{Lemma A.} Let $X$ be a quasi-compact scheme, $R$ a ring and $\psi:R\to\mathcal{O}_X(X)$ a ring homomorphism whose kernel consists of nilpotent elements. Let $h:X\to\operatorname{Spec}R$ be the morphism with $\rho(h)=\psi$ (Prop.~3.25). Then $h(X)$ is dense in $\operatorname{Spec}R$.

\emph{Proof.} Step 1: for $r\in R$, $h^{-1}(D(r))=X_{\psi(r)}$ (Def.~3.11). Let $x\in X$ and let $\mathfrak p$ be the prime $h(x)$. The stalk map $h^\#_x:R_{\mathfrak p}\to\mathcal{O}_{X,x}$ is local, and it sends $r/1$ to the germ $\psi(r)_x$ (because $h^\#(\operatorname{Spec}R)=\psi$). If $r\notin\mathfrak p$, then $r/1$ is a unit and so is $\psi(r)_x$. If $r\in\mathfrak p$, then $r/1$ lies in the maximal ideal, so $\psi(r)_x$ lies in $\mathfrak m_x$ and is not a unit. This proves Step 1.

Step 2: if $D(r)\cap h(X)=\emptyset$, then $r$ is nilpotent. Indeed, by Step 1, $X_{\psi(r)}=\emptyset$. Write $X=U_1\cup\dots\cup U_m$ with $U_i=\operatorname{Spec}B_i$ affine open (Ex.~3.14), and let $a_i\in B_i$ be the restriction of $\psi(r)$. Then $D(a_i)=(U_i)_{a_i}=X_{\psi(r)}\cap U_i=\emptyset$. So $a_i$ lies in every prime of $B_i$, hence is nilpotent. Choose $N$ with $a_i^N=0$ for all $i$. Then $\psi(r)^N$ restricts to $0$ on each $U_i$, so $\psi(r^N)=\psi(r)^N=0$ by the sheaf axiom. Hence $r^N\in\operatorname{Ker}\psi$ is nilpotent, so $r$ is nilpotent.

Step 3: a nonempty open subset of $\operatorname{Spec}R$ contains a nonempty $D(r)$, and $D(r)\neq\emptyset$ means $r$ is not nilpotent. By Step 2, $D(r)$ meets $h(X)$. So $h(X)$ is dense. $\square$

The exercise is the case $R=A$, $\psi=\operatorname{id}_A$, whose kernel is $0$. \hfill$\square$

\exercise{3.16}{Let $X$ be a scheme. Show that $X$ is locally Noetherian if and only if any affine open subscheme of $X$ is Noetherian.}
\pf
Here an affine open subscheme $U=\operatorname{Spec}A$ is called Noetherian when $A$ is a Noetherian ring (equivalently, $U$ is a Noetherian scheme, by Prop.~3.46(b) and Def.~3.45).

($\Leftarrow$) Every point has an affine open neighborhood $U=\operatorname{Spec}A$. By hypothesis $A$ is Noetherian, so $U$ is a Noetherian scheme (a one-term union in Def.~3.45). Hence $X$ is locally Noetherian.

($\Rightarrow$) Let $U=\operatorname{Spec}A$ be an affine open subscheme of $X$. Let $x\in U$. By assumption $x$ has an open neighborhood $W$ which is a Noetherian scheme. Choose $g\in A$ with $x\in D(g)\subseteq U\cap W$ (the $D(g)$ form a base of $U$). Then $D(g)$ is an affine open subset of $W$ (it is isomorphic to $\operatorname{Spec}A_g$ by Lemma~3.7). By Prop.~3.46(b) applied to $W$, the ring $\mathcal{O}_X(D(g))=A_g$ is Noetherian. So $U$ is covered by principal opens $D(g)$ with $A_g$ Noetherian. Since $U$ is quasi-compact (Ex.~1.6), finitely many of them, say $D(g_1),\dots,D(g_n)$, cover $U$. Thus $U$ is a finite union of affine opens with Noetherian rings, i.e.\ a Noetherian scheme. Applying Prop.~3.46(b) to the Noetherian scheme $U$ and its affine open subset $U$ itself, $A=\mathcal{O}_X(U)$ is Noetherian. \hfill$\square$

\exercise{3.17}{Let $f:X\to Y$ be a morphism of schemes. We say that $f$ is quasi-compact if the inverse image of any affine open subset is quasi-compact.
(a) Show that every closed immersion is quasi-compact.
(b) Show that an open immersion $f:X\to Y$ is quasi-compact if $Y$ is locally Noetherian.
Let us suppose in what follows that $f$ is quasi-compact.
(c) Let $\mathcal{J}=\operatorname{Ker}(f^\#:\mathcal{O}_Y\to f_*\mathcal{O}_X)$. Show that the ringed topological subspace $Z:=V(\mathcal{J})$ of $Y$ (see Lemma~2.23) is a scheme (reduce to the case when $Y$ is affine and draw inspiration from the proofs of Propositions 3.12 and 3.20).
(d) Let $j:Z\to Y$ be the closed immersion. Show that we have a morphism $g:X\to Z$ such that $f=j\circ g$, and that if $f$ decomposes into a morphism $g':X\to Z'$ followed by a closed immersion $j':Z'\to Y$, then $Z$ is a closed subscheme of $Z'$.
(e) Show that $f(X)$ is dense in $Z$ (reduce to the case when $Z$ is affine). We call $Z$ the scheme-theoretic closure of $f(X)$ in $Y$. If $X$ is reduced, show that $Z$ is reduced.}
\pf
(a) Let $j:Z\to Y$ be a closed immersion and $V\subseteq Y$ an affine open subset. The restriction $j^{-1}(V)\to V$ is again a closed immersion (the topological map is a closed embedding and the stalk maps are unchanged). By Prop.~3.20, $j^{-1}(V)$ is affine, hence quasi-compact (Ex.~3.14).

(b) Let $V=\operatorname{Spec}A\subseteq Y$ be affine open. By Ex.~3.16, $A$ is Noetherian. Since $f$ is an open immersion, $f^{-1}(V)$ is homeomorphic to the open subset $U=f(X)\cap V$ of $\operatorname{Spec}A$. Write $U=V\setminus V(I)$ with $I=(g_1,\dots,g_r)$ finitely generated. Then $U=\bigcup_sD(g_s)$ is a finite union of affine opens (Lemma~3.7), hence quasi-compact (Ex.~3.14).

For the rest, $f$ is quasi-compact. Recall that $\mathcal{J}$ is a sheaf of ideals of $\mathcal{O}_Y$ with $\mathcal{J}(V)=\operatorname{Ker}f^\#(V)$, and that $Z=V(\mathcal{J})=\{y:\mathcal{J}_y\ne\mathcal{O}_{Y,y}\}$ carries the sheaf $j^{-1}(\mathcal{O}_Y/\mathcal{J})$ (Lemma~2.23).

\emph{A computation in the affine case.} Let $V=\operatorname{Spec}A$ be an affine open subset of $Y$ and let $X'=f^{-1}(V)$, which is quasi-compact. Write $X'=U_1\cup\dots\cup U_n$ with $U_i=\operatorname{Spec}B_i$ affine open, and let $\varphi_i=(f|_{U_i})^\#(V):A\to B_i$. Let $\psi=f^\#(V):A\to\mathcal{O}_X(X')$ and $I=\operatorname{Ker}\psi$. Since $\mathcal{O}_X(X')\to\prod_i B_i$ is injective (sheaf axiom), $I=\bigcap_i\operatorname{Ker}\varphi_i$. We claim: $\mathcal{J}(D(g))=IA_g$ for every $g\in A$. By Prop.~3.14, $f^{-1}(D(g))\cap U_i=D(\varphi_i(g))$, and $f^\#(D(g)):A_g\to\mathcal{O}_X(f^{-1}D(g))\hookrightarrow\prod_i(B_i)_{\varphi_i(g)}$ is, componentwise, the localization of $\varphi_i$. So $a/g^k\in A_g$ is killed by $f^\#(D(g))$ iff for each $i$ there is $m_i$ with $\varphi_i(g^{m_i}a)=0$. Taking $m=\max m_i$ (finitely many $i$), this holds iff $g^ma\in I$ for some $m$, iff $a/g^k\in IA_g$. This proves the claim. Let $i:\operatorname{Spec}A/I\to V$ be the closed immersion of Lemma~3.17. It gives $(\operatorname{Ker}i^\#)(D(g))=I\otimes_AA_g=IA_g=\mathcal{J}(D(g))$. Two subsheaves of $\mathcal{O}_V$ which agree on all $D(g)$ are equal (a section over $U$ is determined by its restrictions to the $D(g)\subseteq U$). So $\operatorname{Ker}i^\#=\mathcal{J}|_V$. By Prop.~2.24 applied to $i$, $\operatorname{Spec}A/I\simeq V(\operatorname{Ker}i^\#)=Z\cap V$ as ringed spaces. $(\ast)$

(c) Cover $Y$ by affine opens $V$. Each $f^{-1}(V)$ is quasi-compact. By $(\ast)$, each $Z\cap V$ is isomorphic to an affine scheme, and the $Z\cap V$ form an open cover of $Z$ by affine schemes. So $Z$ is a scheme, and $j$ is a closed immersion of schemes by Lemma~2.23.

(d) \emph{Factorization.} Let $x\in X$, $y=f(x)$. If $s\in\mathcal{J}_y$ is represented by $s\in\mathcal{J}(V)$, then $f^\#(V)(s)=0$, so $f^\#_x(s)=0$. Hence $\mathcal{J}_y\subseteq\operatorname{Ker}f^\#_x$. This kernel is a proper ideal, since $f^\#_x$ is a ring homomorphism to a nonzero ring. So $\mathcal{J}_y\ne\mathcal{O}_{Y,y}$, i.e.\ $y\in Z$. Thus $f(X)\subseteq Z$, and $f=j\circ g$ for a continuous $g:X\to Z$ (the same map). The morphism $f^\#:\mathcal{O}_Y\to f_*\mathcal{O}_X$ kills $\mathcal{J}$, so it factors through $\mathcal{O}_Y/\mathcal{J}\to f_*\mathcal{O}_X=j_*g_*\mathcal{O}_X$. Applying $j^{-1}$ and using $j^{-1}j_*=\operatorname{id}$ for the closed subset $Z$ (Ex.~2.13, as in the proof of Prop.~2.24), we obtain $g^\#:\mathcal{O}_Z=j^{-1}(\mathcal{O}_Y/\mathcal{J})\to g_*\mathcal{O}_X$. Its stalk at $x$ is $\mathcal{O}_{Y,y}/\mathcal{J}_y\to\mathcal{O}_{X,x}$, induced by $f^\#_x$; it is local because $f^\#_x$ is and the maximal ideal of the quotient is the image of $\mathfrak m_y$. So $g$ is a morphism of schemes, and $j\circ g=f$ by construction.

\emph{Minimality.} Suppose $f=j'\circ g'$ with $j':Z'\to Y$ a closed immersion. By Prop.~2.24, we may identify $Z'$ with $V(\mathcal{J}')$, $\mathcal{J}'=\operatorname{Ker}j'^\#$, with its structure sheaf $\mathcal{O}_Y/\mathcal{J}'$ and the canonical $j'$. We have $f^\#=j'_*(g'^\#)\circ j'^\#$, hence $\mathcal{J}'\subseteq\mathcal{J}$. Consequently $\mathcal{J}'_y\ne\mathcal{O}_{Y,y}$ whenever $\mathcal{J}_y\neq\mathcal{O}_{Y,y}$, i.e.\ $Z\subseteq Z'$, and $Z$ is closed in $Z'$. The canonical surjection $\mathcal{O}_Y/\mathcal{J}'\to\mathcal{O}_Y/\mathcal{J}$, restricted to $Z'$ and using that the stalks of $\mathcal{O}_Y/\mathcal{J}$ vanish off $Z$ (so that its restriction to $Z'$ equals $h_*\mathcal{O}_Z$, $h:Z\to Z'$ the inclusion), gives $h^\#:\mathcal{O}_{Z'}\to h_*\mathcal{O}_Z$. Its stalks $\mathcal{O}_{Y,z}/\mathcal{J}'_z\to\mathcal{O}_{Y,z}/\mathcal{J}_z$ are surjective and local. So $h$ is a closed immersion with $j'\circ h=j$, which says that $Z$ is a closed subscheme of $Z'$.

(e) \emph{Density.} It suffices to show that $f(X)\cap Z\cap V$ is dense in $Z\cap V$ for every $V$ in an affine open cover of $Y$ (a subset which is dense in each member of an open cover is dense). Fix $V=\operatorname{Spec}A$ and use the notation of $(\ast)$. Then $Z\cap V\simeq\operatorname{Spec}A/I$, and $f(X)\cap V=f(X')=g(X')$. The morphism $X'\to Z\cap V\simeq\operatorname{Spec}(A/I)$ corresponds under $\rho$ to the map $\bar\psi:A/I\to\mathcal{O}_X(X')$ induced by $\psi$ (compare with (3.4): composing with $\operatorname{Spec}A/I\to\operatorname{Spec}A$ gives $\psi$). The map $\bar\psi$ is injective by the definition of $I$. By Lemma A (Ex.~3.15) applied to the quasi-compact scheme $X'$, $g(X')$ is dense in $\operatorname{Spec}(A/I)$.

\emph{Reducedness.} Assume $X$ is reduced. Then $\mathcal{O}_X(U)$ is reduced for every open $U$: if $s^n=0$, then $s_x$ is nilpotent in the reduced ring $\mathcal{O}_{X,x}$, so $s_x=0$ for all $x$ and $s=0$. In particular $\mathcal{O}_X(X')$ is reduced. The subring $A/I\subseteq\mathcal{O}_X(X')$ is reduced. Thus $Z\cap V\simeq\operatorname{Spec}A/I$ has reduced local rings (localizations of a reduced ring). As $V$ varies, these cover $Z$, so $Z$ is reduced. \hfill$\square$

\exercise{3.18}{Let $X$ be an affine algebraic variety over a field $k$. Show that there exists a projective variety $\overline X$ over $k$ such that $X$ is isomorphic to a dense open subscheme of $\overline X$.}
\pf
Write $X=\operatorname{Spec}A$ with $A=k[t_1,\dots,t_n]/I$ a finitely generated $k$-algebra (Def.~3.47). Let $S=k[T_0,\dots,T_n]$ with the usual grading. For $P\in k[t_1,\dots,t_n]$ nonzero of degree $\delta$ let $P^h=T_0^{\delta}P(T_1/T_0,\dots,T_n/T_0)\in S_\delta$ be its homogenization, and set $0^h=0$. Let $\bar I\subseteq S$ be the (homogeneous) ideal generated by all $P^h$, $P\in I$. Put $B=S/\bar I$ and $\overline X=\operatorname{Proj}B$. This is a closed subscheme of $\mathbb{P}^n_k$ (Lemma~3.41), hence a projective variety over $k$. Let $\epsilon:S\to k[t_1,\dots,t_n]$ be the $k$-algebra homomorphism $T_0\mapsto1$, $T_j\mapsto t_j$, and write $b(1,t)=\epsilon(b)$.

\emph{Step 1: for $b\in S$ homogeneous, $b\in\bar I\iff\epsilon(b)\in I$.} If $b=\sum_jc_jP_j^h$ with $c_j\in S$ and $P_j\in I$, then $\epsilon(b)=\sum_j\epsilon(c_j)P_j\in I$ because $\epsilon(P_j^h)=P_j$. Conversely let $\epsilon(b)=P\in I$, $b\ne0$. Write $b=T_0^mb'$ with $b'$ homogeneous and not divisible by $T_0$. Then $b'$ has a monomial free of $T_0$, so $\deg P=\deg b'$, and $\epsilon$ is injective on homogeneous polynomials of a fixed degree. Hence $b'=P^h$ and $b=T_0^mP^h\in\bar I$.

\emph{Step 2: $T_0$ is a non-zero-divisor on $B$.} Since $\bar I$ is homogeneous, it suffices to treat a homogeneous $b$ with $T_0b\in\bar I$. Then $\epsilon(b)=\epsilon(T_0b)\in I$ by Step 1, so $b\in\bar I$ by Step 1.

\emph{Step 3: $D_+(T_0)\simeq X$.} By Prop.~3.38, $D_+(T_0)\simeq\operatorname{Spec}B_{(T_0)}$. We have $S_{(T_0)}=k[T_j/T_0]_{j}\simeq k[t_1,\dots,t_n]$ via $b/T_0^N\mapsto\epsilon(b)$ for $b\in S_N$. Since $B_{T_0}=S_{T_0}/\bar IS_{T_0}$ is $\mathbb{Z}$-graded and $\bar IS_{T_0}$ is a homogeneous ideal, $B_{(T_0)}=S_{(T_0)}/(\bar IS_{T_0})_0$. The degree-$0$ part $(\bar IS_{T_0})_0$ consists of the sums of elements $b/T_0^N$ with $b\in\bar I$ homogeneous of degree $N$; under the isomorphism above these become the elements $\epsilon(b)$, which lie in $I$ by Step 1, and every $P\in I$ occurs (as $\epsilon(P^h)$). So the image of $(\bar IS_{T_0})_0$ is exactly $I$, and $B_{(T_0)}\simeq k[t_1,\dots,t_n]/I=A$. Hence $X=\operatorname{Spec}A$ is isomorphic, as a $k$-scheme, to the open subscheme $D_+(T_0)$ of $\overline X$.

\emph{Step 4: $D_+(T_0)$ is dense in $\overline X$.} A nonempty open subset of $\overline X$ contains a nonempty $D_+(g)$ with $g\in B_+$ homogeneous. Suppose $D_+(g)\cap D_+(T_0)=D_+(gT_0)=\emptyset$. Then $B_{(gT_0)}=0$ (Prop.~3.38), so $1=0$ in $B_{gT_0}$, i.e.\ $(gT_0)^m=0$ in $B$ for some $m$. By Step 2, $g^m=0$. Then $B_g=0$, so $B_{(g)}=0$ and $D_+(g)=\emptyset$, a contradiction. Thus every nonempty open set meets $D_+(T_0)$. \hfill$\square$

\exercise{3.19}{Let $K$ be a number field. Let $\mathcal{O}_K$ be the ring of integers of $K$ (i.e., the set of elements of $K$ which are integral over $\mathbb{Z}$). Using the finiteness theorem of the class group $\mathrm{cl}(K)$, show that every open subset of $\operatorname{Spec}\mathcal{O}_K$ is principal. Deduce from this that every open subscheme of $\operatorname{Spec}\mathcal{O}_K$ is affine.}
\pf
We use the standard facts: $\mathcal{O}_K$ is a Dedekind domain (so its nonzero prime ideals are maximal, and every nonzero ideal is a product of prime ideals, uniquely), and the class group $\operatorname{cl}(K)$ of invertible ideals modulo principal ones is finite. Let $h=|\operatorname{cl}(K)|$.

Let $U$ be an open subset of $X=\operatorname{Spec}\mathcal{O}_K$, $U=X\setminus V(J)$ for an ideal $J$. If $U=\emptyset$, then $U=D(0)$. Otherwise $J\neq0$ (as $V(0)=X$), and we factor $J=\mathfrak p_1^{e_1}\cdots\mathfrak p_r^{e_r}$ with distinct nonzero primes ($r\ge0$). A prime contains $J$ iff it contains some $\mathfrak p_i$, iff it equals $\mathfrak p_i$ (the $\mathfrak p_i$ are maximal and a prime containing a product contains a factor). So $V(J)=\{\mathfrak p_1,\dots,\mathfrak p_r\}$.

By Lagrange, the class of $\mathfrak p_i$ in the finite group $\operatorname{cl}(K)$ satisfies $[\mathfrak p_i]^h=1$, i.e.\ $\mathfrak p_i^h=\pi_i\mathcal{O}_K$ is principal. Put $g=\pi_1\cdots\pi_r$ ($g=1$ if $r=0$). Then $(g)=(\mathfrak p_1\cdots\mathfrak p_r)^h$, and the primes containing $g$ are exactly $\mathfrak p_1,\dots,\mathfrak p_r$. So $V(g)=V(J)$ and $U=D(g)$. Hence every open subset of $X$ is principal.

Consequently every open subscheme $U=D(g)$ of $\operatorname{Spec}\mathcal{O}_K$ is isomorphic to $\operatorname{Spec}(\mathcal{O}_K)_g$ by Lemma~3.7, hence is affine. \hfill$\square$

\exercise{3.20}{Let $A$ be a ring, $G$ a finite group of automorphisms of $A$, and $A^G$ the subring of elements of $A$ which are invariant under $G$. Let $p:\operatorname{Spec}A\to\operatorname{Spec}A^G$ denote the morphism induced by the inclusion $A^G\to A$.
(a) Show that $G$ acts naturally on $\operatorname{Spec}A$. Show that $p(x_1)=p(x_2)$ if and only if there exists $\sigma\in G$ such that $\sigma(x_1)=x_2$.
(b) Show that $A$ is integral over $A^G$ (let $a\in A$, and consider the polynomial $P(T):=\prod_{\sigma\in G}(T-\sigma a)\in A^G[T]$). Deduce from this that $p$ is surjective (Exercise 1.8).
(c) Let $a\in A$. Let $P(T)$ be as above and write $P(T)=T^d+\sum_{i\le d-1}b_iT^i$, $b_i\in A^G$. Show that $p(D(a))=\bigcup_iD(b_i)$. Show that $p$ is open (in the topological sense).
(d) Show that for any $b\in A^G$, we have $p^{-1}(D(b))=D(bA)$ and $(A^G)_b=(A_b)^G$. Let $V$ be an open subset of $\operatorname{Spec}A^G$. Show that $G$ acts on the scheme $p^{-1}(V)$ and that $\mathcal{O}_{\operatorname{Spec}A^G}(V)=\mathcal{O}_{\operatorname{Spec}A}(p^{-1}(V))^G$.}
\pf
Write $X=\operatorname{Spec}A$, $S=\operatorname{Spec}A^G$, $d=|G|$.

(a) For $\sigma\in G$ let $\varphi_\sigma=\operatorname{Spec}(\sigma^{-1}):X\to X$. It sends a prime $\mathfrak q$ to $(\sigma^{-1})^{-1}(\mathfrak q)=\sigma(\mathfrak q)$. Since $\operatorname{Spec}$ is contravariant, $\varphi_\sigma\varphi_\tau=\operatorname{Spec}(\tau^{-1}\sigma^{-1})=\varphi_{\sigma\tau}$. So $\sigma\mapsto\varphi_\sigma$ is a homomorphism $G\to\operatorname{Aut}(X)$ (automorphisms of schemes, Prop.~3.14), and we write $\sigma(x)=\varphi_\sigma(x)$. On global sections, $\varphi_\sigma^\#(X)=\sigma^{-1}$ (Prop.~3.14), so the invariants of $\mathcal{O}_X(X)$ are $A^G$. Since $\sigma^{-1}$ fixes $A^G$, we have $p\circ\varphi_\sigma=p$. So $p(\mathfrak q)=\mathfrak q\cap A^G$ is constant on $G$-orbits: $\sigma(\mathfrak q)\cap A^G=\sigma(\mathfrak q\cap A^G)=\mathfrak q\cap A^G$.

Conversely, let $\mathfrak q_1,\mathfrak q_2$ be primes with $\mathfrak q_1\cap A^G=\mathfrak q_2\cap A^G=:P$. Let $b\in\mathfrak q_2$. Then $N(b)=\prod_{\sigma\in G}\sigma(b)$ is $G$-invariant and lies in $\mathfrak q_2$ (the factor for $\sigma=1$), so $N(b)\in\mathfrak q_2\cap A^G=P\subseteq\mathfrak q_1$. As $\mathfrak q_1$ is prime, $\sigma(b)\in\mathfrak q_1$ for some $\sigma$, i.e.\ $b\in\sigma^{-1}(\mathfrak q_1)$. Therefore $\mathfrak q_2\subseteq\bigcup_{\sigma}\sigma(\mathfrak q_1)$, a finite union of primes. By prime avoidance, $\mathfrak q_2\subseteq\tau(\mathfrak q_1)$ for some $\tau$. By symmetry, $\mathfrak q_1\subseteq\rho(\mathfrak q_2)$ for some $\rho$. Then $\mathfrak q_2\subseteq\tau(\mathfrak q_1)\subseteq\tau\rho(\mathfrak q_2)$. Put $\sigma_0=\tau\rho$ and let $n$ be its order. Then $\mathfrak q_2\subseteq\sigma_0(\mathfrak q_2)\subseteq\sigma_0^2(\mathfrak q_2)\subseteq\dots\subseteq\sigma_0^n(\mathfrak q_2)=\mathfrak q_2$, so all these are equal and $\tau(\mathfrak q_1)=\mathfrak q_2$. This proves (a).

(b) For $a\in A$, the polynomial $P(T)=\prod_{\sigma\in G}(T-\sigma a)\in A[T]$ is monic of degree $d$ and $P(a)=0$ (factor $\sigma=1$). For $\tau\in G$, applying $\tau$ to the coefficients permutes the factors, so $\tau P=P$. Hence $P\in A^G[T]$ and $a$ is integral over $A^G$.

Surjectivity of $p$ (this is Ex.~1.8(c) for the injective integral map $A^G\to A$; here is the argument). Let $P_0$ be a prime of $R=A^G$ and $T=R\setminus P_0$. Since $R\subseteq A$ and $0\notin T$, the ring $T^{-1}A$ is nonzero. Let $M$ be a maximal ideal of $T^{-1}A$. The ring $T^{-1}A$ is integral over $T^{-1}R=R_{P_0}$. The domain $D=R_{P_0}/(M\cap R_{P_0})$ is contained in the field $F=T^{-1}A/M$ and $F$ is integral over $D$. Then $D$ is a field: if $0\neq x\in D$ and $x^{-1}$ satisfies $x^{-m}+c_1x^{-m+1}+\dots+c_m=0$ with $c_i\in D$, multiplying by $x^{m-1}$ gives $x^{-1}\in D$. So $M\cap R_{P_0}=P_0R_{P_0}$, and the prime $\mathfrak q=\{a\in A:a/1\in M\}$ satisfies $\mathfrak q\cap R=P_0$. Hence $p$ is surjective.

(c) Write $P(T)=T^d+\sum_{i<d}b_iT^i$. Here $b_i=\pm e_{d-i}(\sigma a)_{\sigma\in G}$ is, up to sign, the elementary symmetric polynomial of degree $d-i\ge1$ in the $\sigma a$.

($\subseteq$) Let $\mathfrak q\in D(a)$ and $P_0=\mathfrak q\cap A^G$. If all $b_i\in P_0\subseteq\mathfrak q$, then $a^d=-\sum_{i<d}b_ia^i\in\mathfrak q$, so $a\in\mathfrak q$, a contradiction. So some $b_i\notin P_0$, i.e.\ $p(\mathfrak q)\in D(b_i)$.

($\supseteq$) Let $P_0\in D(b_i)$ for some $i$. By (b) there is $\mathfrak q$ with $\mathfrak q\cap A^G=P_0$. If $\sigma(a)\in\mathfrak q$ for all $\sigma\in G$, then $b_i$, being a sum of products of $d-i\ge1$ of the $\sigma a$, lies in $\mathfrak q\cap A^G=P_0$, a contradiction. So $\sigma(a)\notin\mathfrak q$ for some $\sigma$, i.e.\ $a\notin\sigma^{-1}(\mathfrak q)$. The prime $\sigma^{-1}(\mathfrak q)$ lies in $D(a)$ and over $P_0$ by (a). Hence $P_0\in p(D(a))$.

Since the sets $D(a)$ form a base of $X$ and $p$ commutes with unions, $p$ maps open sets to open sets.

(d) For $b\in A^G$, $p^{-1}(D(b))=\{\mathfrak q:b\notin\mathfrak q\cap A^G\}=D(b)\subseteq X$, which is $D(bA)$ in the notation of the statement. $G$ acts on $A_b$ by $\sigma(a/b^n)=\sigma(a)/b^n$, and $(A^G)_b\to A_b$ is injective (localization is exact) with image in $(A_b)^G$. Conversely, let $a/b^n\in(A_b)^G$. For each $\sigma$, $\sigma(a)/b^n=a/b^n$ in $A_b$, so $b^{m_\sigma}(\sigma a-a)=0$ for some $m_\sigma$. With $m=\max_\sigma m_\sigma$, the element $b^ma$ is $G$-invariant, and $a/b^n=b^ma/b^{n+m}\in(A^G)_b$. So $(A^G)_b=(A_b)^G$.

Let $V\subseteq S$ be open and $W=p^{-1}(V)$. Since $p\varphi_\sigma=p$, $W$ is stable under every $\varphi_\sigma$, and $\varphi_\sigma$ restricts to an automorphism of the open subscheme $W$. So $G$ acts on $W$, and on the ring $\mathcal{O}_X(W)$ by the $\varphi_\sigma^\#(W)$; the invariants are the elements fixed by all of them. Moreover, the maps $\varphi_\sigma^\#$ commute with restriction to stable opens. Write $V=\bigcup_jD(b_j)$ with $b_j\in A^G$. Then $W=\bigcup_jD(b_j)$ (in $X$). By the first part, $\mathcal{O}_S(D(b_j))=(A^G)_{b_j}=(A_{b_j})^G=\mathcal{O}_X(D(b_j))^G$, and likewise for $D(b_jb_l)$. We show that $p^\#(V):\mathcal{O}_S(V)\to\mathcal{O}_X(W)^G$ is bijective. Injectivity: if $s\in\mathcal{O}_S(V)$ maps to $0$, then each $s|_{D(b_j)}\in(A^G)_{b_j}\subseteq A_{b_j}$ is $0$, so $s=0$. Surjectivity: let $t\in\mathcal{O}_X(W)^G$. Then each $t|_{D(b_j)}$ is $G$-invariant, hence lies in $\mathcal{O}_S(D(b_j))$. The sections $t_j=t|_{D(b_j)}$ agree on $D(b_jb_l)$ in $A_{b_jb_l}$, hence in the subring $(A^G)_{b_jb_l}$. By the sheaf property on $S$ they come from a unique $s\in\mathcal{O}_S(V)$, and $p^\#(V)(s)=t$ by the sheaf property on $X$. So $\mathcal{O}_S(V)=\mathcal{O}_X(p^{-1}V)^G$, via $p^\#(V)$. \hfill$\square$

\exercise{3.21}{Let $X$ be a scheme and $G$ a finite group acting on $X$ (i.e.\ $G$ is endowed with a group homomorphism $G\to\operatorname{Aut}(X)$). We define the quotient scheme $X/G$ by the universal property of Exercise 2.14, where we replace the ringed topological spaces by schemes. It does not always exist.
(a) Let $A$ be a ring and $G$ a finite group of automorphisms of $A$, which we identify with a group of automorphisms of the scheme $\operatorname{Spec}A$. Show that the quotient scheme $(\operatorname{Spec}A)/G$ exists and is isomorphic to $\operatorname{Spec}(A^G)$. Show that this is also the quotient as a ringed topological space (Exercise 2.14).
(b) Let $U$ be an open subscheme of $\operatorname{Spec}A$ that is stable under $G$. Preserving the notation of the preceding exercise, show that the quotient scheme $U/G$ exists and is isomorphic to $p(U)$.
(c) Let $G$ be a finite group acting on a scheme $X$. We suppose that every point $x\in X$ has an affine open neighborhood that is stable under $G$ (see Exercise 3.3.23 for examples of such $X$). Show that the quotient scheme $X/G$ exists.}
\pf
\emph{Setting.} A quotient scheme of $X$ by $G$ is a scheme $Y$ with a morphism $q:X\to Y$ such that $q\circ\sigma=q$ for all $\sigma\in G$, and such that every morphism of schemes $f:X\to Z$ with $f\circ\sigma=f$ for all $\sigma$ factors as $f=\tilde f\circ q$ for a unique morphism $\tilde f:Y\to Z$. It is unique up to unique isomorphism. Let $X=\operatorname{Spec}A$, $S=\operatorname{Spec}A^G$, $p:X\to S$, and use the notation $\varphi_\sigma$ of Ex.~3.20. The group $G$ is a finite group of automorphisms of $A$. (If $G$ is given as acting on $A$ through a homomorphism $G\to\operatorname{Aut}(A)$ with a kernel, replace $G$ by its image $H$: invariants, orbits and the condition $f\sigma=f$ only involve the image.)

\textbf{Lemma U.} Let $V\subseteq S$ be open and $W=p^{-1}(V)$. If $g,g':V\to Z$ are morphisms of schemes with $g\circ p|_W=g'\circ p|_W$, then $g=g'$.

\emph{Proof.} The map $p:W\to V$ is surjective (Ex.~3.20(b)), so $g=g'$ as maps of sets. Now $(g\circ p)^\#=g_*(p^\#)\circ g^\#$ and the same for $g'$. By Ex.~3.20(d), $p^\#(V'')$ is injective for every open $V''\subseteq V$ (it is the inclusion of $G$-invariants), so $g_*(p^\#)$ is injective. Hence $g^\#=g'^\#$. $\square$

\textbf{Lemma Q.} Let $V'\subseteq S$ be open and $W=p^{-1}(V')$, with $G$ acting on $W$ as in Ex.~3.20(d). Then $p|_W:W\to V'$ is a quotient scheme of $W$ by $G$.

\emph{Proof.} We have $p\varphi_\sigma=p$. Let $f:W\to Z$ be a morphism of schemes with $f\varphi_\sigma=f$ for all $\sigma$. Cover $Z$ by affine opens $Z_i=\operatorname{Spec}C_i$. Then $W_i=f^{-1}(Z_i)$ is open and $G$-stable (since $f\varphi_\sigma=f$). Let $V_i=p(W_i)$; it is open by Ex.~3.20(c). By Ex.~3.20(a), the fibres of $p$ are the $G$-orbits, so $p^{-1}(V_i)=\bigcup_\sigma\sigma(W_i)=W_i$. The $V_i$ cover $V'$ ($p$ is surjective and the $W_i$ cover $W$).

Let $\psi_i=\rho(f|_{W_i}):C_i\to\mathcal{O}_X(W_i)$ (Prop.~3.25). By (3.4), $\rho(f\varphi_\sigma|_{W_i})=\varphi_\sigma^\#(W_i)\circ\psi_i$, which equals $\psi_i$. So $\psi_i$ takes values in $\mathcal{O}_X(W_i)^G=\mathcal{O}_S(V_i)$ (Ex.~3.20(d)). Let $g_i:V_i\to Z_i\subseteq Z$ be the morphism with $\rho(g_i)=\psi_i$ (viewed as a map $C_i\to\mathcal{O}_S(V_i)$). Then $\rho(g_i\circ p|_{W_i})=p^\#(V_i)\circ\psi_i=\psi_i=\rho(f|_{W_i})$, so $g_i\circ p=f$ on $W_i$ by injectivity of $\rho$. On $V_i\cap V_j$, whose preimage is $W_i\cap W_j$, both $g_i$ and $g_j$ satisfy $g\circ p=f$. By Lemma~U they coincide. Gluing morphisms (Ex.~2.11) we get $\tilde f:V'\to Z$ with $\tilde f\circ p=f$. Uniqueness follows from Lemma~U. $\square$

(a) and (b). Taking $V'=S$ in Lemma Q shows that $(S,p)$ is a quotient scheme of $X$, so $(\operatorname{Spec}A)/G$ exists and is isomorphic to $\operatorname{Spec}A^G$. For (b), let $U\subseteq X$ be open and $G$-stable. Then $p(U)$ is open (Ex.~3.20(c)) and $p^{-1}(p(U))=\bigcup_\sigma\sigma(U)=U$. Lemma Q with $V'=p(U)$ shows that $p|_U:U\to p(U)$ is the quotient scheme of $U$ by $G$.

\emph{Ringed space quotient.} Let $Y=X/G$ be the set of orbits with the quotient topology and $q:X\to Y$, and let $\mathcal{O}_Y(V)=\mathcal{O}_X(q^{-1}V)^G$, as in Ex.~2.14. By Ex.~3.20(a), $p$ induces a bijection $\theta:Y\to S$ with $\theta q=p$. It is a homeomorphism: $p$ is continuous, open (Ex.~3.20(c)) and surjective, so a subset $V\subseteq S$ is open iff $p^{-1}(V)$ is open, which is the definition of the topology on $Y$. By Ex.~3.20(d), $p^\#$ identifies $\mathcal{O}_S(V)$ with $\mathcal{O}_X(p^{-1}V)^G=\mathcal{O}_Y(\theta^{-1}V)$. So $(\theta,\cdot):(Y,\mathcal{O}_Y)\to(S,\mathcal{O}_S)$ is an isomorphism of ringed spaces compatible with the maps from $X$. By Ex.~2.14, $(Y,\mathcal{O}_Y,q)$ is the quotient in the category of ringed topological spaces, and hence so is $(S,p)$.

(c) Let $\{U_i\}$ be a cover of $X$ by $G$-stable affine opens, $U_i=\operatorname{Spec}A_i$. The group $G$ acts on $U_i$ by scheme automorphisms, i.e.\ by ring automorphisms of $A_i$ (Prop.~3.25), through a finite group of automorphisms (as in the Setting). Let $S_i=\operatorname{Spec}A_i^G$ and $p_i:U_i\to S_i$.

Let $Y=X/G$ be the set of orbits with the quotient topology and $q:X\to Y$. It is open: $q^{-1}(q(W))=\bigcup_\sigma\sigma(W)$ is open for $W$ open. Define $\mathcal{O}_Y(V)=\mathcal{O}_X(q^{-1}V)^G$. This is a sheaf of rings: restriction preserves invariants, and if $s_a\in\mathcal{O}_Y(V_a)$ are compatible on overlaps for an open cover $V=\bigcup V_a$, then the glued section of $\mathcal{O}_X(q^{-1}V)$ is invariant, because $\sigma(s)$ and $s$ agree on every $q^{-1}V_a$. Let $V_i=q(U_i)$, which is open, with $q^{-1}(V_i)=U_i$ because $U_i$ is stable. The bijection $\theta_i:V_i\to S_i$ induced by $p_i$ (both are the set of orbits in $U_i$) is a homeomorphism: $V\subseteq V_i$ is open iff $q^{-1}V$ is open iff $p_i^{-1}(\theta_iV)$ is open iff $\theta_iV$ is open, as $p_i$ is open, continuous and surjective. By Ex.~3.20(d), $p_i^\#$ identifies $\mathcal{O}_{S_i}(\theta_iV)$ with $\mathcal{O}_X(q^{-1}V)^G=\mathcal{O}_Y(V)$. Hence $(V_i,\mathcal{O}_Y|_{V_i})\simeq S_i$ is an affine scheme. The $V_i$ cover $Y$, so $Y$ is a scheme. Moreover $q|_{U_i}$ corresponds to $p_i$ under $\theta_i$, so $q$ is a morphism of schemes (locality of stalk maps is checked on the cover $U_i$), and $q\varphi_\sigma=q$.

\emph{Universal property.} Let $f:X\to Z$ be a morphism of schemes with $f\varphi_\sigma=f$. Since $q^{-1}(V_i)=U_i\simeq p_i^{-1}(S_i)$, Lemma Q (for $S_i$ in place of $S$ and $V'=S_i$) yields a morphism $f_i:V_i\to Z$ with $f_i\circ q=f$ on $U_i$. On $V_i\cap V_j$, whose preimage is $U_i\cap U_j$, the morphisms $f_i$ and $f_j$ satisfy $g\circ q=f$; since $V_i\cap V_j$ is an open subset of $V_i\simeq S_i$, Lemma~U (for $S_i$) shows $f_i=f_j$ there. Gluing gives $\tilde f:Y\to Z$ with $\tilde f\circ q=f$. If $\tilde f'$ is another one, then $\tilde f|_{V_i}$ and $\tilde f'|_{V_i}$ agree by Lemma~U, so $\tilde f=\tilde f'$. Thus $(Y,q)$ is the quotient scheme $X/G$. \hfill$\square$

\exercise{3.22}{Let $k$ be a field of characteristic $0$. We let $G:=\mathbb{Z}$ act on the polynomial ring $k[T]$ by $n:T\mapsto T+n$ if $n\in\mathbb{Z}$.
(a) Show that $G$ is a subgroup of the group of automorphisms (of $k$-algebras) of $k[T]$. We identify it with a group of automorphisms of $\mathbb{A}^1_k$. Show that the only open subschemes of $\mathbb{A}^1_k$ that are stable under $G$ are $\emptyset$ and $\mathbb{A}^1_k$ itself.
(b) Show that the quotient scheme $\mathbb{A}^1_k/G$ exists and is equal to $\operatorname{Spec}k$.
(c) Show that $\mathbb{A}^1_k/G$ is not the quotient as a ringed topological space (Exercise 2.14).}
\pf
Here $G=\mathbb{Z}$ is infinite, so the finite-group results of Ex.~3.20--3.21 are not used. The definition of quotient scheme is the one in Ex.~3.21 (universal property, $G$ acting through $n\mapsto\varphi_n=\operatorname{Spec}(\tau_{-n})$).

(a) Let $\tau_n:k[T]\to k[T]$ be the $k$-algebra homomorphism with $T\mapsto T+n$ (where $n$ means $n\cdot1\in k$). Then $\tau_n\tau_m=\tau_{n+m}$ and $\tau_0=\operatorname{id}$, so each $\tau_n$ is an automorphism and $n\mapsto\tau_n$ is a homomorphism $\mathbb{Z}\to\operatorname{Aut}_k(k[T])$. It is injective since $\operatorname{char}k=0$: $\tau_n=\operatorname{id}$ means $T+n=T$, i.e.\ $n=0$ in $k$. So $G$ is identified with a subgroup. It acts on $X=\mathbb{A}^1_k=\operatorname{Spec}k[T]$ by $\varphi_n=\operatorname{Spec}(\tau_{-n})$, which sends a prime $\mathfrak q$ to $\tau_n(\mathfrak q)$.

Let $U\ne\emptyset$ be an open $G$-stable subset and $F=X\setminus U$, which is closed and $G$-stable. Since $k[T]$ is a principal ideal domain, a proper closed subset is a finite set of maximal ideals, so $F$ is either $X$ or a finite set of maximal ideals. Suppose $F\ne\emptyset$ is finite and contains $(\pi)$ with $\pi=T^m+c_{m-1}T^{m-1}+\dots$ monic irreducible, $m\ge1$. Then $\tau_n((\pi))=(\pi(T+n))$, and $\pi(T+n)=T^m+(c_{m-1}+nm)T^{m-1}+\dots$ is monic of degree $m$. Two monic polynomials generating the same ideal are equal, and for $n\ne n'$ the coefficients of $T^{m-1}$ differ, since $m(n-n')\ne0$ in characteristic $0$. So the orbit of $(\pi)$ is infinite, contradicting finiteness of $F$ and stability. Hence $F=\emptyset$ or $F=X$ (which is impossible for $U\neq\emptyset$). So $U=X$. The only stable open subsets are $\emptyset$ and $X$.

(b) Let $\pi:X\to\operatorname{Spec}k$ be the structural morphism. Since each $\tau_n$ is a $k$-algebra map, $\pi\varphi_n=\pi$.

\emph{Existence.} Let $f:X\to Z$ be a morphism of schemes with $f\varphi_n=f$ for all $n$. Pick $x\in X$ and an affine open $Z_0=\operatorname{Spec}C$ of $Z$ containing $f(x)$. Then $f^{-1}(Z_0)$ is a nonempty open set, stable under $G$ because $f\varphi_n=f$. By (a), $f^{-1}(Z_0)=X$, so $f$ is a morphism $X\to Z_0$. Let $\psi=\rho(f):C\to k[T]$ (Prop.~3.25). By (3.4), $\rho(f\varphi_n)=\tau_{-n}\circ\psi=\psi$ for all $n$. So the image of $\psi$ lies in $k[T]^G=\{P:P(T+n)=P(T)\ \forall n\in\mathbb{Z}\}$. For such $P$, the polynomial $P(T)-P(0)$ vanishes at $T=0,1,2,\dots$, which are pairwise distinct elements of $k$ as $\operatorname{char}k=0$. So $P=P(0)\in k$ is constant, and $k[T]^G=k$. Hence $\psi$ factors as $C\xrightarrow{\ \psi_0\ }k\hookrightarrow k[T]$. Let $\tilde f:\operatorname{Spec}k\to Z_0\subseteq Z$ be the morphism with $\rho(\tilde f)=\psi_0$. Then $\rho(\tilde f\circ\pi)=\psi=\rho(f)$, so $\tilde f\circ\pi=f$ by injectivity of $\rho$.

\emph{Uniqueness.} Let $\tilde f,\tilde f':\operatorname{Spec}k\to Z$ satisfy $\tilde f\pi=\tilde f'\pi$. Since $\pi$ is surjective, the images of the single point of $\operatorname{Spec}k$ coincide; choose an affine open $Z_0=\operatorname{Spec}C$ containing it. Both morphisms land in $Z_0$, and $\rho(\tilde f\pi)=\rho(\tilde f)$ followed by $k\hookrightarrow k[T]$ by (3.4). Since $k\to k[T]$ is injective, $\rho(\tilde f)=\rho(\tilde f')$, so $\tilde f=\tilde f'$ by Prop.~3.25.

So $(\operatorname{Spec}k,\pi)$ is the quotient scheme $\mathbb{A}^1_k/G$.

(c) We show that $\pi:X\to\operatorname{Spec}k$ does not have the universal property among ringed topological spaces. Let $\xi$ be the generic point of $X$ and let $Z=\{\eta,c\}$ be a two-point space with the indiscrete topology (the only open sets are $\emptyset$ and $Z$), with sheaf $\mathcal{O}_Z(Z)=k$, $\mathcal{O}_Z(\emptyset)=0$. This is a sheaf, and both stalks equal $k$, a local ring, so $Z$ is a ringed topological space. Define $f:X\to Z$ by $f(\xi)=\eta$ and $f(x)=c$ for every closed point $x$, and $f^\#(Z):k\to\mathcal{O}_X(X)=k[T]$ the inclusion. Then $f$ is continuous (the inverse images of $\emptyset$ and $Z$ are open), $f^\#$ is a morphism of sheaves of rings, and every stalk map $k\to\mathcal{O}_{X,x}$ is local (the maximal ideal of $k$ is $0$). Also $f\circ\varphi_n=f$: the maps $\varphi_n$ fix $\xi$ and permute closed points, and $\varphi_n^\#\circ f^\#=f^\#$ because $\tau_n$ is the identity on $k$. If $f=\tilde f\circ\pi$ for some map $\tilde f$ from the one-point space $\operatorname{Spec}k$, then $f$ would be constant, but it takes the two values $\eta$ and $c$. So no factorization exists, and $\operatorname{Spec}k$ is not the quotient in the category of ringed topological spaces. (Any other morphism $p':X\to\operatorname{Spec}k$ with $p'\varphi_n=p'$ is $\pi$ composed with an automorphism of $\operatorname{Spec}k$, since $\rho(p'):k\to k[T]$ lands in $k[T]^G=k$; so the same argument applies to it.)

\emph{Remark.} By Ex.~2.14 the quotient ringed space exists: it is the set $X/G$ with the quotient topology and $\mathcal{O}(V)=\mathcal{O}_X(q^{-1}V)^G$. By (a) its only open sets are $\emptyset$ and the whole space, and its ring of global sections is $k[T]^G=k$. So it is the indiscrete space on the infinite set $X/G$ (it contains the point $\xi$ and the $G$-orbits of closed points) with global sections $k$. This has many points and is not $\operatorname{Spec}k$, and $f$ above factors through it by $X\to X/G\to Z$. \hfill$\square$

\section*{2.4 Reduced schemes and integral schemes}

Conventions for this section. Schemes denoted $X$, $Y$ are non-empty whenever we speak of connectedness, irreducibility or integrality. For a scheme $X$, a point $x\in X$ and $s\in\mathcal O_X(U)$ with $x\in U$, $s_x$ denotes the germ; $\fm_x$ is the maximal ideal of $\mathcal O_{X,x}$. For an irreducible space, a \emph{generic point} is a point $\xi$ with $\overline{\{\xi\}}=X$; it is unique when it exists on a scheme, because the underlying space of a scheme is $T_0$ (Ex.~4.7 below). We say that $y$ \emph{specializes to} $x$ (or $x$ is a specialization of $y$) if $x\in\overline{\{y\}}$.

\exercise{4.1}{Let $k$ be a field and $P\in k[T_1,\dots,T_n]$. Show that $\Spec(k[T_1,\dots,T_n]/(P))$ is reduced (resp.\ irreducible; resp.\ integral) if and only if $P$ has no square factor (resp.\ admits only one irreducible factor; resp.\ is irreducible).}
\pf
\emph{Remark on the hypotheses.} The statement is meant for $P$ non-zero and non-constant. If $P=0$ the scheme is $\mathbb A^n_k$, which is integral although $0$ is not irreducible, and if $P$ is a unit the scheme is empty. So we assume $P\neq 0$ and $P$ not a unit. ``Only one irreducible factor'' means: $P=uQ^e$ with $u$ a unit, $Q$ irreducible, $e\ge1$.

Let $R=k[T_1,\dots,T_n]$ (a UFD), $A=R/(P)$ and $X=\Spec A$. Write $P=u\,P_1^{e_1}\cdots P_r^{e_r}$ with $r\ge1$, $P_i$ pairwise non-associate irreducible (hence prime) elements, $e_i\ge1$.

\emph{Radical of $(P)$.} We claim $\sqrt{(P)}=(P_1\cdots P_r)$. If $Q^m\in(P)$ then each $P_i$ divides $Q^m$, hence $Q$ ($P_i$ prime); since the $P_i$ are pairwise non-associate primes in a UFD, $P_1\cdots P_r\mid Q$. Conversely $(P_1\cdots P_r)^{m}$ is divisible by $P$ for $m=\max e_i$. Thus the nilradical of $A$ is $\sqrt{(P)}/(P)=(P_1\cdots P_r)/(P)$.

\emph{Reduced.} By Prop.~4.2(a),(b), the affine scheme $X$ is reduced iff $A$ is reduced, iff $\sqrt{(P)}=(P)$, iff $P\mid P_1\cdots P_r$, iff all $e_i=1$, i.e.\ $P$ has no square factor.

\emph{Irreducible.} By Prop.~4.7(a),(c) applied to $\Spec A$ and the ideal $0$, $X$ is irreducible iff the nilradical of $A$ is prime. Under the correspondence between ideals of $A$ and ideals of $R$ containing $(P)$, this holds iff $\sqrt{(P)}=(P_1\cdots P_r)$ is a prime ideal of $R$. If $r=1$ it is $(P_1)$, which is prime. If $r\ge2$, put $a=P_1$ and $b=P_2\cdots P_r$. Then $ab\in(P_1\cdots P_r)$, but $P_1\cdots P_r\nmid a$ (as $P_2\nmid P_1$) and $P_1\cdots P_r\nmid b$ (as $P_1\nmid b$), so the ideal is not prime. Hence $X$ is irreducible iff $r=1$, i.e.\ $P$ has only one irreducible factor.

\emph{Integral.} $X$ is integral iff $A$ is an integral domain (remark after Def.~4.16), iff $(P)$ is a prime ideal of $R$, iff $P$ is a prime element, iff $P$ is irreducible ($R$ is a UFD). This is consistent with the two previous cases: $r=1$ and $e_1=1$. \hfill$\square$

\exercise{4.2}{Let $X$ be a scheme and $x\in X$. Show that the image of the morphism $\Spec\mathcal O_{X,x}\to X$ is the set of points of $X$ that specialize to $x$.}
\pf
Recall (Example~3.16) that the morphism is the composition $\Spec \mathcal O_{X,x}\to U=\Spec A\hookrightarrow X$ for any affine open $U\ni x$, where $A=\mathcal O_X(U)$, the first map being induced by $A\to \mathcal O_{X,x}=A_\fp$ and $\fp$ the prime ideal corresponding to $x$. We compute the image in $U$.

\emph{Image of $\Spec A_\fp\to\Spec A$.} The primes of $A_\fp$ are the $\fq A_\fp$ with $\fq$ a prime of $A$ contained in $\fp$, and the contraction of $\fq A_\fp$ is $\fq$. So the image is $\{\fq\in\Spec A:\fq\subseteq\fp\}$.

\emph{Points specializing to $x$.} Let $y\in X$ with $x\in\overline{\{y\}}$. Since $U$ is an open neighbourhood of $x$ and $x$ lies in the closure of $\{y\}$, $U$ meets $\{y\}$, so $y\in U$. Let $y$ correspond to the prime $\fq$ of $A$. The closure of $\{y\}$ in $U$ is $\overline{\{y\}}\cap U=V(\fq)$ (Lemma~4.11(a)), so $x\in\overline{\{y\}}$ iff $\fp\in V(\fq)$ iff $\fq\subseteq\fp$ (Lemma~4.11(b)). Conversely a point in the image lies in $U$ and corresponds to some $\fq\subseteq\fp$, hence specializes to $x$. Therefore the image of $\Spec \mathcal O_{X,x}\to X$ is exactly $\{y\in X: x\in\overline{\{y\}}\}$. \hfill$\square$

\exercise{4.3}{Let $\mathcal O_K$ be a discrete valuation ring with field of fractions $K$ and uniformizing parameter $t$ (a generator of the maximal ideal). Show that $\Spec K[T]$ can be identified with an open subscheme of $\Spec \mathcal O_K[T]$. Determine the set of closed points of $\Spec K[T]$ which specialize to the point corresponding to the maximal ideal $(T,t)$ of $\mathcal O_K[T]$.}
\pf
Let $v$ be the normalized valuation of $\mathcal O_K$, $k=\mathcal O_K/t\mathcal O_K$ the residue field, $B=\mathcal O_K[T]$.

\emph{Open subscheme.} Every non-zero element of $\mathcal O_K$ is $ut^m$ with $u$ a unit and $m\ge0$, hence $K=\mathcal O_K[1/t]=(\mathcal O_K)_t$ and $K[T]=B_t$. By Example~3.15 (and Lemma~3.7), $\Spec K[T]=\Spec B_t\to\Spec B$ is an isomorphism onto the principal open subscheme $D(t)$. Its points are the primes of $B$ not containing $t$, corresponding to the primes of $K[T]$ by $\fq\mapsto\fq\cap B$.

\emph{Gauss lemma.} Call $G\in B$ \emph{primitive} if its image in $k[T]$ is non-zero, i.e.\ some coefficient is a unit. Since $k[T]$ is a domain, a product of primitive polynomials is primitive. Consequently, if $F'\in B$ is primitive, $H\in K[T]$ and $F'H\in B$, then $H\in B$: write $H=t^{-r}H_0$ with $r\in\mathbf Z$ and $H_0\in B$ primitive; then $F'H_0$ is primitive and $t^{-r}F'H_0\in B$ forces $r\le0$. Hence for primitive $F'$:
\[ F'K[T]\cap B=F'B. \]

\emph{Closed points.} The closed points of $\Spec K[T]$ are the maximal ideals $(F)$ with $F$ monic irreducible in $K[T]$ ($K[T]$ is a PID; closed points of an affine scheme are the maximal ideals). Write $F=T^d+a_{d-1}T^{d-1}+\dots+a_0$, $a_d:=1$. Let $s=-\min_{0\le i\le d}v(a_i)\ge0$ (the minimum is $\le0$ because $a_d=1$), and $F'=t^sF\in B$. Then $F'$ is primitive (some coefficient has valuation $0$ by choice of $s$). The prime of $B$ corresponding to $(F)$ is $\fq_F=FK[T]\cap B=F'K[T]\cap B=F'B$.

\emph{Specialization.} A point $y$ of $D(t)$ with prime $\fq$ specializes to the point $\fm=(T,t)$ of $\Spec B$ iff $\fm\in\overline{\{y\}}=V(\fq)$, i.e.\ iff $\fq\subseteq\fm$ (Lemma~4.11). Now $\fm=\{G\in B: G(0)\in t\mathcal O_K\}$ (the kernel of $B\to k$, $T\mapsto0$). So $\fq_F\subseteq\fm$ iff $F'\in\fm$ iff $v(t^sa_0)\ge1$ iff $v(a_0)>-s=\min_i v(a_i)$.

\emph{Answer.} The closed points of $\Spec K[T]$ specializing to $(T,t)$ are exactly the maximal ideals $(F)$ with $F=T^d+a_{d-1}T^{d-1}+\dots+a_0\in K[T]$ monic irreducible such that the constant term does not realize the minimal valuation of the coefficients:
\[ v(a_0)>\min_{0\le i\le d}v(a_i)\qquad(\text{with }v(0)=+\infty,\ a_d=1). \]
Equivalently, the primitive multiple $F'$ of $F$ in $\mathcal O_K[T]$ has constant term in $t\mathcal O_K$. Examples: $F=T$; $F=T-a$ with $a\in t\mathcal O_K$ (here $\min(v(a),0)=0<v(a)$ for $v(a)\ge1$ and the condition says exactly $v(a)\ge1$); every monic irreducible $F\in\mathcal O_K[T]$ with $F(0)\in t\mathcal O_K$; and also some $F$ with non-integral coefficients, e.g.\ for $\mathcal O_K=\mathbf Z_{(p)}$ and $F=T^2+p^{-1}T+1$ (irreducible over $\mathbf Q$ since its discriminant $p^{-2}-4=(1-4p^2)/p^2$ is negative, hence not a square), $F'=pT^2+T+p\in(T,p)$. \hfill$\square$

\exercise{4.4}{We say that a scheme $X$ is connected if the underlying topological space of $X$ is connected. For example, an irreducible scheme is connected. Let $X$ be a scheme having only a finite number of irreducible components $\{X_i\}_i$. Show that $X$ is connected if and only if for any pair $i,j$, there exist indices $i_0=i,i_1,\dots,i_r=j$ such that $X_{i_l}\cap X_{i_{l+1}}\neq\emptyset$ for every $l<r$. Also show that $X$ is integral if and only if $X$ is connected, and integral at every point $x\in X$.}
\pf
Throughout $X\ne\emptyset$. Every irreducible space is connected: if $Z=C_1\sqcup C_2$ with $C_i$ closed then $Z=C_1$ or $Z=C_2$. The $X_i$ are closed and cover $X$ (they are the irreducible components).

\emph{First equivalence.} Say $i\sim j$ if a chain as in the statement exists; this is an equivalence relation. Let $\mathcal C$ be the set of classes, and for $C\in\mathcal C$ put $Y_C=\bigcup_{i\in C}X_i$, a finite union of closed sets, hence closed. If $C\ne C'$ then $Y_C\cap Y_{C'}=\emptyset$ (a point in $X_i\cap X_j$ with $i\in C$, $j\in C'$ would give $i\sim j$). So $X=\bigsqcup_{C}Y_C$ is a finite disjoint union of closed sets, and each $Y_C$ is also open (its complement is the finite union of the other $Y_{C'}$). Moreover each $Y_C$ is connected: it is the union of connected sets $X_i$ ($i\in C$), any two of which are linked by a chain of pairwise meeting ones, and a union of connected sets, two of which meet, is connected. In summary, \emph{the $Y_C$ are exactly the connected components of $X$; they are open and closed, and finite in number} (we call this fact $(\ast)$).
Thus $X$ is connected iff there is only one class, i.e.\ iff any $i,j$ are linked by a chain. (If $X$ is connected and $X=\bigsqcup Y_C$ with clopen non-empty pieces, there is a single piece. Conversely, if there is one class, $X=Y_C$ is connected.)

\emph{Second equivalence.} ($\Rightarrow$) An integral scheme is irreducible, hence connected. It is integral at each point: pick an affine open $U\ni x$; $\mathcal O_X(U)$ is an integral domain by Prop.~4.17, and $\mathcal O_{X,x}$ is a localization of it, hence a domain.

($\Leftarrow$) Assume $X$ is connected and every $\mathcal O_{X,x}$ is a domain. Then $X$ is reduced. For $x\in X$, $\Spec\mathcal O_{X,x}$ is irreducible (Prop.~4.7(c), a domain has a unique minimal prime $(0)$), so by Prop.~4.12(b) exactly one irreducible component of $X$ passes through $x$. Hence distinct irreducible components of $X$ are disjoint. Now let $i,j$ be arbitrary; by the first part there is a chain $i=i_0,\dots,i_r=j$ with $X_{i_l}\cap X_{i_{l+1}}\ne\emptyset$; disjointness forces $X_{i_l}=X_{i_{l+1}}$ for all $l$, so $X_i=X_j$. Thus $X$ has a single irreducible component, namely $X$ itself, so $X$ is irreducible. Being reduced and irreducible, $X$ is integral. \hfill$\square$

\exercise{4.5}{Let $X$ be a scheme. Show that every irreducible component of $X$ is contained in a connected component (in the topological sense). Show that if $X$ is locally Noetherian, then the connected components are open. If $X$ is Noetherian, then there are only finitely many connected components.}
\pf
\emph{Irreducible components.} Let $Z$ be an irreducible component, $x\in Z$, and $C$ the connected component of $x$, the union of all connected subsets containing $x$. $Z$ is irreducible, hence connected (Ex.~4.4), and contains $x$, so $Z\subseteq C$.

\emph{Locally Noetherian case.} Let $C$ be a connected component of $X$ and $x\in C$. Choose an affine open $U=\Spec A\ni x$ with $A$ Noetherian (definition of locally Noetherian). The ring $A$ is Noetherian, so $U$ has finitely many irreducible components (Prop.~4.9, for the Noetherian scheme $U$). By the argument of $(\ast)$ in Ex.~4.4 (the classes of components under the chain relation give a partition of $U$ into finitely many clopen connected sets), $U$ is a finite disjoint union of connected clopen subsets of $U$. Let $W$ be the one containing $x$; $W$ is connected and open in $U$, hence open in $X$. As $W$ is a connected set containing $x$, $W\subseteq C$. So every point of $C$ has an open neighbourhood contained in $C$, and $C$ is open.

\emph{Noetherian case.} A Noetherian scheme is quasi-compact and locally Noetherian. The connected components are open by the above, pairwise disjoint and cover $X$. By quasi-compactness finitely many of them cover $X$; since they are disjoint and non-empty, no other component can exist. So there are finitely many connected components. \hfill$\square$

\exercise{4.6}{Let $A$ be a ring. We say that $e\in A$ is idempotent if $e\neq0$ and $e^2=e$. We will say that $e$ is indecomposable if it cannot be written as a sum of two idempotent elements. Let $X$ be a scheme.

(a) Show the equivalence of the following properties: (i) $X$ is connected; (ii) $\mathcal O_X(X)$ has no other idempotent elements than $1$; (iii) $\Spec\mathcal O_X(X)$ is connected.

(b) Show that any local scheme (i.e., the spectrum of a local ring) is connected.

(c) Let us suppose that the connected components of $X$ are open (e.g., $X$ locally Noetherian). Let $U$ be a connected component of $X$. Show that there exists a unique idempotent element $e$ of $\mathcal O_X(X)$ such that $e|_U=1$ and $e|_{X\setminus U}=0$. Show that this establishes a canonical bijection from the set of connected components of $X$ onto the set of indecomposable idempotent elements of $\mathcal O_X(X)$, the converse map being given by $e\mapsto V(1-e)$.}
\pf
\emph{Correction about ``indecomposable''.} Read literally (any two idempotents $e_1,e_2$ with $e=e_1+e_2$), no element is indecomposable in a ring of characteristic $2$: in $\mathbf F_2\times\mathbf F_2$ we have $(1,0)=(1,1)+(0,1)$. The intended notion, which we use, is: $e$ is \emph{indecomposable} if it is not a sum $e=e_1+e_2$ of two idempotents with $e_1e_2=0$ (orthogonal idempotents). When $2$ is invertible the orthogonality is automatic: $e_1+e_2$ idempotent gives $2e_1e_2=0$. Also $V(f)$ for a global section $f$ means $\{x\in X: f_x\in\fm_x\}$. We take $X\ne\emptyset$ in (a), and note that the zero ring has no idempotent in the sense above.

\emph{Preliminary: idempotents are locally $0$ or $1$.} If $R$ is a local ring and $\varepsilon^2=\varepsilon$, then $\varepsilon\in\{0,1\}$: if $\varepsilon\notin\fm$ it is a unit and $\varepsilon(\varepsilon-1)=0$ gives $\varepsilon=1$; if $\varepsilon\in\fm$ then $1-\varepsilon$ is a unit and $\varepsilon(1-\varepsilon)=0$ gives $\varepsilon=0$. Now let $e\in\mathcal O_X(X)$ with $e^2=e$ (we also allow $e=0$ here). Then $e_x\in\{0,1\}$ for all $x$. Put $U_e=\{x:e_x=1\}=V(1-e)$ (as $(1-e)_x\in\fm_x$ iff $e_x=1$) and $V_e=\{x:e_x=0\}$. Then $X=U_e\sqcup V_e$. If $e_x=1$ then $e=1$ on a neighbourhood of $x$ (a germ equality means equality of sections on a neighbourhood), so $U_e$ is open; likewise $V_e$ is open. So $U_e$ is open and closed. Moreover $U_e=\emptyset$ iff $e=0$, and $U_e=X$ iff $e=1$ (the sheaf property). Conversely for a clopen $W\subseteq X$ the section $1_W$ equal to $1$ on $W$ and $0$ on $X\setminus W$ (glued on the open cover $\{W,X\setminus W\}$) is an idempotent or zero, with $U_{1_W}=W$. So
\[ e\longmapsto U_e=V(1-e)\]
is a bijection from $\{e\in\mathcal O_X(X): e^2=e\}$ onto the set of clopen subsets of $X$, with inverse $W\mapsto 1_W$. Under it $e_1e_2=0$ and $e=e_1+e_2$ correspond to $U_{e_1}\cap U_{e_2}=\emptyset$ and $U_e=U_{e_1}\sqcup U_{e_2}$ (stalkwise, $e_{1,x}e_{2,x}=0$ means not both are $1$, and then the sum is the indicator of the union; conversely the indicator functions of a disjoint decomposition are orthogonal idempotents summing to $1_W$).

(a) (i)$\Leftrightarrow$(ii): By the bijection, $\mathcal O_X(X)$ has an idempotent $e\ne0,1$ iff $X$ has a clopen subset different from $\emptyset$ and $X$, iff $X$ is disconnected. For a non-empty $X$ the idempotent $1$ always exists. (ii)$\Leftrightarrow$(iii): apply the equivalence (i)$\Leftrightarrow$(ii), which holds for every non-empty scheme, to the affine scheme $\Spec\mathcal O_X(X)$, whose ring of global sections is $\mathcal O_X(X)$ itself.

(b) Let $R$ be local. By the preliminary, $R$ (with $X=\Spec R$, $\mathcal O_X(X)=R$) has no idempotents other than $1$ (its only idempotent in the sense of the exercise), so $\Spec R$ is connected by (a). (Alternatively: the closed point $\fm$ lies in every non-empty closed set $V(I)$, so two disjoint non-empty closed sets cannot exist.)

(c) Let $U$ be a connected component; it is closed, and by assumption open, so $X\setminus U$ is open too. Put $e=1_U$. Then $e^2=e$ and $e\ne0$, and the section $e$ is unique with $e|_U=1$, $e|_{X\setminus U}=0$ because $\mathcal O_X$ is a sheaf and $\{U,X\setminus U\}$ covers $X$. We have $U_e=U=V(1-e)$.

\emph{Bijection.} By the bijection above, indecomposable idempotents correspond to non-empty clopen subsets $W\subseteq X$ admitting no decomposition $W=W_1\sqcup W_2$ into two non-empty clopen subsets of $X$, i.e.\ to non-empty clopen subsets $W$ that are connected (a subset of $W$ is clopen in $X$ iff it is clopen in $W$, since $W$ is clopen). So we must show: the connected components of $X$ are exactly the non-empty connected clopen subsets. Each component $U$ is clopen (by assumption) and connected. Conversely, if $W$ is a non-empty connected clopen subset, $W$ lies in a connected component $C$ and $W$ is clopen in the connected $C$, so $W=C$. Therefore $U\mapsto e_U=1_U$ is a bijection from connected components onto indecomposable idempotents, whose inverse is $e\mapsto U_e=V(1-e)$. \hfill$\square$

\exercise{4.7}{We say that a topological space $X$ verifies the separation axiom $T_0$ if for every pair of points $x\neq y$, there exists an open subset which contains one of the points and not the other one. Show that the underlying topological space of a scheme verifies $T_0$.}
\pf
Let $X$ be a scheme and $x\neq y$ in $X$. Take an affine open $U=\Spec A$ containing $x$. If $y\notin U$, then $U$ is open, contains $x$ and not $y$. If $y\in U$, let $x,y$ correspond to distinct primes $\fp\ne\fq$ of $A$. Then $\fp\not\subseteq\fq$ or $\fq\not\subseteq\fp$; say $\fp\not\subseteq\fq$, and pick $a\in\fp\setminus\fq$. The principal open set $D(a)$ contains $\fq$ (i.e.\ $y$) and not $\fp$ (i.e.\ $x$). It is open in $U$, hence in $X$. In the other case, symmetrically, there is an open set containing $x$ and not $y$. \hfill$\square$

\exercise{4.8}{Let $X$ be a quasi-compact scheme (Exercise~3.14). Show that $X$ contains a closed point. See also Exercise~3.3.26 for a counterexample when $X$ is not quasi-compact.}
\pf
Assume $X\ne\emptyset$ (otherwise there is nothing to prove, or the statement is meant for non-empty $X$).

\emph{Step 1: a minimal non-empty closed subset exists.} Order the non-empty closed subsets of $X$ by inclusion. Let $(Z_\alpha)$ be a chain. If $\bigcap Z_\alpha=\emptyset$, the complements $X\setminus Z_\alpha$ form an open cover of $X$; by quasi-compactness, finitely many cover $X$, and since the chain is totally ordered the largest of these complements is $X$, so the corresponding $Z_\alpha=\emptyset$, a contradiction. So $\bigcap Z_\alpha$ is a non-empty closed lower bound. By Zorn's lemma there is a minimal non-empty closed subset $Z$.

\emph{Step 2: $\overline{\{z\}}=Z$ for every $z\in Z$.} Let $z\in Z$ and choose an affine open $V=\Spec B\ni z$. Then $Z\cap V$ is a non-empty closed subset of $V$, so it is $V(I)$ with $I\ne B$, and $V(I)$ contains a maximal ideal $\fm$ of $B$; call the corresponding point $w$. The point $w$ is closed in $V$, and $w\in Z$. Then $\overline{\{w\}}\subseteq Z$ is closed and non-empty, hence equals $Z$ by minimality. Therefore $Z\cap V=\overline{\{w\}}\cap V$ is the closure of $\{w\}$ in $V$, which is $\{w\}$. Since $z\in Z\cap V$, $z=w$, and so $\overline{\{z\}}=Z$.

\emph{Step 3: $Z$ is one point.} For $z,z'\in Z$, Step 2 gives $\overline{\{z\}}=Z=\overline{\{z'\}}$. An open set containing $z$ must meet $\{z'\}$ (because $z\in\overline{\{z'\}}$), so contains $z'$, and symmetrically. So $z,z'$ have the same open neighbourhoods, and $z=z'$ by Ex.~4.7. Thus $Z=\{z\}$ is a closed point of $X$. \hfill$\square$

\exercise{4.9}{Let $X$ be a Noetherian scheme. Show that the set of points $x\in X$ such that $\mathcal O_{X,x}$ is reduced (resp.\ is an integral domain) is open.}
\pf
Both properties of $x$ depend only on $\mathcal O_{X,x}=\mathcal O_{U,x}$ for any open $U\ni x$, and openness of a subset is local. $X$ is covered by affine opens $\Spec A$ with $A$ Noetherian, so it suffices to treat $X=\Spec A$, $A$ Noetherian.

\emph{Reduced locus.} Let $N$ be the nilradical of $A$. For a prime $\fp$, the nilradical of $A_\fp$ is $N_\fp$: if $a/s$ is nilpotent then $ta^m=0$ for some $m$ and $t\notin\fp$, so $(ta)^m=0$, $ta\in N$ and $a/s=ta/ts\in N_\fp$; the other inclusion is clear. As $A$ is Noetherian, $N=(n_1,\dots,n_m)$. Then $A_\fp$ is reduced iff $n_i/1=0$ in $A_\fp$ for all $i$ iff $\Ann(n_i)\not\subseteq\fp$ for all $i$. So the non-reduced locus is $\bigcup_{i=1}^mV(\Ann(n_i))$, a finite union of closed sets, so closed. The reduced locus is open.

\emph{Domain locus.} Since $A$ is Noetherian it has finitely many minimal primes $\fp_1,\dots,\fp_s$ (Prop.~4.9 and Prop.~4.7(b)). The minimal primes of $A_\fp$ are the $\fq A_\fp$ with $\fq\subseteq\fp$ minimal among primes of $A$ contained in $\fp$, that is (any prime inside $\fq$ is inside $\fp$) exactly the minimal primes $\fp_i$ of $A$ contained in $\fp$. A local ring is a domain iff it is reduced and has exactly one minimal prime (if reduced, $0$ is the intersection of its minimal primes). Hence $A_\fp$ is a domain iff $\fp\notin$ non-reduced locus and $\fp$ contains exactly one $\fp_i$, i.e.\ $\fp\notin\bigcup_{i\ne j}V(\fp_i)\cap V(\fp_j)$. The set of domain points is thus the complement of the closed set
\[ \bigcup_{i=1}^mV(\Ann(n_i))\ \cup\ \bigcup_{i\ne j}\bigl(V(\fp_i)\cap V(\fp_j)\bigr),\]
so it is open. \hfill$\square$

\exercise{4.10}{Let $f:X\to\Spec A$ be a quasi-compact morphism (Exercise~3.17). Let $I$ be an ideal of $A$. Show that $f(X)\subseteq V(I)$ if and only if $f^\#(\Spec A)(I)\subseteq\mathcal O_X(X)$ is nilpotent.}
\pf
\emph{Interpretation.} We write $\varphi=f^\#(\Spec A):A\to\mathcal O_X(X)$ and read ``$\varphi(I)$ is nilpotent'' as: every element of $\varphi(I)$ is nilpotent. If $I$ is finitely generated this is the same as nilpotence of the ideal $\varphi(I)\mathcal O_X(X)$ (a finitely generated ideal of nilpotent elements is nilpotent). For general $I$ the stronger reading (the ideal is nilpotent) is false: let $A=\mathcal O_X(X)=k[x_1,x_2,\dots]/(x_i^i)$, $X=\Spec A$, $f=\mathrm{id}$, $I=(x_1,x_2,\dots)$. Then $f(X)=V(I)$ since $I$ is in the nilradical, but $I^N\ni x_{N+1}^N\ne0$ for each $N$.

\emph{Key formula.} Let $x\in X$, $f(x)=\fp$. The local homomorphism $f^\#_x:A_\fp=\mathcal O_{\Spec A,\fp}\to\mathcal O_{X,x}$ satisfies $(f^\#_x)^{-1}(\fm_x)=\fp A_\fp$. For $a\in A$, the image of $a$ in $\mathcal O_{X,x}$ is $\varphi(a)_x$. Hence
\[ a\in\fp\iff a/1\in\fp A_\fp\iff \varphi(a)_x\in\fm_x. \]

($\Leftarrow$) Suppose each $\varphi(a)$, $a\in I$, is nilpotent. For $x\in X$, $\varphi(a)_x$ is nilpotent, hence in $\fm_x$, so $a\in f(x)$ for all $a\in I$, i.e.\ $f(x)\in V(I)$. (Quasi-compactness is not needed here.)

($\Rightarrow$) Suppose $f(X)\subseteq V(I)$ and let $a\in I$. For every $x\in X$, $a\in f(x)$, so $\varphi(a)_x\in\fm_x$. Since $\Spec A$ is affine and $f$ is quasi-compact, $X=f^{-1}(\Spec A)$ is quasi-compact (Ex.~3.14), hence a finite union of affine opens $U_j=\Spec B_j$, $1\le j\le m$. Let $b_j=\varphi(a)|_{U_j}\in B_j$. For every prime $\fq$ of $B_j$, corresponding to a point $x\in U_j$, we have $b_j/1\in\fq(B_j)_\fq=\fm_x$, hence $b_j\in\fq$. So $b_j$ lies in all primes of $B_j$, hence is nilpotent (the nilradical is the intersection of all primes): $b_j^{n_j}=0$. With $n=\max n_j$, $\varphi(a)^n$ restricts to $0$ on each $U_j$, so $\varphi(a)^n=0$. This holds for each $a\in I$ (with $n$ depending on $a$), proving the claim. If $I=(a_1,\dots,a_r)$ is finitely generated, the ideal generated by the nilpotent elements $\varphi(a_i)$ is nilpotent. \hfill$\square$

\exercise{4.11}{Let $f:X\to Y$ be a morphism of irreducible schemes with respective generic points $\xi_X,\xi_Y$. We say that $f$ is dominant if $f(X)$ is dense in $Y$. Let us suppose that $X$, $Y$ are integral. Show that the following properties are equivalent: (i) $f$ is dominant; (ii) $f^\#:\mathcal O_Y\to f_*\mathcal O_X$ is injective; (iii) for every open subset $V$ of $Y$ and every open subset $U\subseteq f^{-1}(V)$, the map $\mathcal O_Y(V)\to\mathcal O_X(U)$ is injective; (iv) $f(\xi_X)=\xi_Y$; (v) $\xi_Y\in f(X)$.}
\pf
\emph{Correction in (iii).} For $U=\emptyset$ we have $\mathcal O_X(U)=0$, so (iii) as printed fails trivially for $V\ne\emptyset$. We therefore read (iii) with $U$ and $V$ non-empty (then $V\neq\emptyset$ is automatic). Note that for non-empty open $V$ of $Y$ we have $\xi_Y\in V$ (since $\overline{\{\xi_Y\}}=Y$), and $\mathcal O_Y(V)$ is an integral domain, in particular non-zero.

We prove (i)$\Leftrightarrow$(iv)$\Leftrightarrow$(v), (iv)$\Rightarrow$(ii)$\Rightarrow$(i) and (iv)$\Rightarrow$(iii)$\Rightarrow$(v).

\emph{(iv)$\Rightarrow$(v)}: trivial. \emph{(v)$\Rightarrow$(i)}: $f(X)\ni\xi_Y$ and $\overline{\{\xi_Y\}}=Y$.

\emph{(i)$\Rightarrow$(iv)}: By continuity, $f(X)=f(\overline{\{\xi_X\}})\subseteq\overline{\{f(\xi_X)\}}$. Since $f(X)$ is dense, $Y\subseteq\overline{\{f(\xi_X)\}}$, so $f(\xi_X)$ is a generic point of $Y$. Two generic points $\eta,\eta'$ of $Y$ have the same open neighbourhoods (as above), hence are equal by Ex.~4.7. So $f(\xi_X)=\xi_Y$.

\emph{(iv)$\Rightarrow$(iii) and (ii).} Let $V\subseteq Y$ be a non-empty open subset and $U\subseteq f^{-1}(V)$ a non-empty open subset. Since $U\ne\emptyset$ we have $\xi_X\in U$, and $f(\xi_X)=\xi_Y$. Consider the maps
\[ \mathcal O_Y(V)\xrightarrow{f^\#(V)}\mathcal O_X(f^{-1}(V))\to\mathcal O_X(U)\to\mathcal O_{X,\xi_X},\]
whose composition equals $\mathcal O_Y(V)\to\mathcal O_{Y,\xi_Y}\xrightarrow{f^\#_{\xi_X}}\mathcal O_{X,\xi_X}$ by the definition of the stalk map. The first map $\mathcal O_Y(V)\to\mathcal O_{Y,\xi_Y}=K(Y)$ is injective by Prop.~4.18(b), and $f^\#_{\xi_X}:K(Y)\to K(X)$ is a ring homomorphism between fields, hence injective. So the composition is injective, and so is its first factor $\mathcal O_Y(V)\to\mathcal O_X(U)$. This is (iii). For (ii): let $V$ be any open subset of $Y$. If $V=\emptyset$ there is nothing to show. If $V\ne\emptyset$ then $\xi_Y\in V$, so $\xi_X\in f^{-1}(V)$ and $U=f^{-1}(V)$ is non-empty; by the above $\mathcal O_Y(V)\to\mathcal O_X(f^{-1}(V))=(f_*\mathcal O_X)(V)$ is injective. Hence $f^\#$ is injective.

\emph{(ii)$\Rightarrow$(i)}: Suppose $f(X)$ is not dense. Then $V=Y\setminus\overline{f(X)}$ is a non-empty open set with $f^{-1}(V)=\emptyset$, so $(f_*\mathcal O_X)(V)=\mathcal O_X(\emptyset)=0$ while $\mathcal O_Y(V)\ne0$. Hence $f^\#(V)$ is not injective, contradicting (ii).

\emph{(iii)$\Rightarrow$(v)}: Let $x\in X$, and choose an affine open $V=\Spec B\ni f(x)$ in $Y$ and an affine open $U=\Spec C\ni x$ of $X$ contained in $f^{-1}(V)$. As $X,Y$ are integral, $B$ and $C$ are integral domains (Prop.~4.17). The morphism $f|_U:U\to V$ is $\Spec\psi$ where $\psi:B\to C$ is the map $\mathcal O_Y(V)\to\mathcal O_X(U)$ of (iii) (Prop.~3.25), hence injective by assumption. The point $(0)$ of $\Spec C$ maps to $\psi^{-1}(0)=\Ker\psi=(0)\in\Spec B$. The point $(0)$ of $\Spec B$ is the generic point of $V$; its closure in $Y$ contains $V$, which is dense in the irreducible $Y$ (Prop.~4.5(a)); hence its closure is $Y$ and the point is $\xi_Y$ (uniqueness of the generic point, as in (i)$\Rightarrow$(iv)). Hence $\xi_Y\in f(U)\subseteq f(X)$. \hfill$\square$

\exercise{4.12}{Let $B$ be a graded ring. Let $Y$ be a reduced closed subscheme of $\operatorname{Proj}B$. Show that there exists a homogeneous ideal $I$ of $B$ such that $Y\simeq\operatorname{Proj}B/I$.}
\pf
Here $B=\bigoplus_{d\ge0}B_d$ (as in Subsection~2.3.3) and $Y$ is a closed subscheme with the closed immersion $j:Y\to\operatorname{Proj}B$. We prove more precisely that $I$ can be chosen so that $Y$ and $\operatorname{Proj}B/I$ are isomorphic \emph{as closed subschemes} of $\operatorname{Proj}B$.

\emph{Choice of $I$.} Let $Z=j(Y)$, a closed subset of $\operatorname{Proj}B$, so $Z=V_+(J)$ for a homogeneous ideal $J$. Let $I=\bigcap_{\fp\in Z}\fp$ ($I=B$ if $Z=\emptyset$). It is an intersection of homogeneous ideals, so it is homogeneous, and radical. Moreover $V_+(I)=Z$: every $\fp\in Z$ contains $I$, so $Z\subseteq V_+(I)$; conversely $I\supseteq J$ (each $\fp\in Z$ contains $J$), so $V_+(I)\subseteq V_+(J)=Z$.

\emph{The closed immersion $\operatorname{Proj}B/I\to\operatorname{Proj}B$.} Let $\phi:B\to B/I$ be the quotient map; it is a homomorphism of graded rings with $\phi(B_+)(B/I)=(B/I)_+$. By Lemma~3.40 (with $r=1$, $M=(B/I)_+$) we get a morphism $g:\operatorname{Proj}B/I\to\operatorname{Proj}B$ with $g^{-1}(D_+(h))=D_+(\phi(h))$ for homogeneous $h\in B_+$, and on this open set $g$ is $\Spec$ of $B_{(h)}\to(B/I)_{(\phi(h))}$. This ring map is surjective: an element of $(B/I)_{(\phi(h))}$ is $\phi(a)/\phi(h)^n$ with $a$ homogeneous of degree $n\deg h$, which is the image of $a/h^n$. Its kernel is $\{a/h^n:\ h^ma\in I\text{ for some }m\}=IB_h\cap B_{(h)}=I_{(h)}$. The open sets $D_+(h)$ cover $\operatorname{Proj}B$ (Subsection~2.3.3). By Lemma~3.17, on each $D_+(h)$ the morphism is a closed immersion with image $V(I_{(h)})$, where $I_{(h)}=IB_h\cap B_{(h)}$; by Lemma~3.36(d) this is $V_+(I)\cap D_+(h)$ under $D_+(h)\simeq\Spec B_{(h)}$. Being a homeomorphism onto a closed subset, and surjectivity of the map of sheaves, are local on the target; hence $g$ is a closed immersion with image $V_+(I)=Z$.

\emph{$\operatorname{Proj}B/I$ is reduced.} For homogeneous $h\in B_+$ the ring $(B/I)_{(\phi(h))}$ is the degree-$0$ part of the localization $(B/I)_{\phi(h)}$. The ring $B/I$ is reduced (as $I$ is radical), hence so is its localization, hence so is its subring $(B/I)_{(\phi(h))}$. Since these rings are the rings of sections of the affine opens $D_+(\phi(h))$ covering $\operatorname{Proj}B/I$, the scheme $\operatorname{Proj}B/I$ is reduced (Prop.~4.2(b)).

\emph{Conclusion.} $Y$ (via $j$) and $\operatorname{Proj}B/I$ (via $g$) are two reduced closed subschemes of $\operatorname{Proj}B$ with the same underlying closed subset $Z$. By the uniqueness in Prop.~4.2(e), they coincide: there is an isomorphism $Y\to\operatorname{Proj}B/I$ compatible with the immersions into $\operatorname{Proj}B$. In particular $Y\simeq\operatorname{Proj}B/I$. \hfill$\square$

\section*{2.5 Dimension}

Throughout, $k$ is a field, rings are commutative with $1$, and $\operatorname{ht}$ is the height. ``Algebraic variety'' is as in Def.~3.47, and a ``graded ring'' is $\mathbb N$-graded as in Section 2.3.2. Results of the text are cited by number. Honest notes: Ex.~5.5 is false as printed when the ideal contains non-zero homogeneous elements of degree $0$ (a counterexample is given), and it is proved for $I\subseteq A_+$, which is the case needed in practice; in Ex.~5.12(d) the set $X$ must be non-empty (and $n\ge1$), and the proof uses the separatedness of $\mathbb P^n_k$ from Chapter 3 (Prop.~3.3.6, Example 3.3.7); in Ex.~5.15 the ring $B$ is a homogeneous algebra over $k$ (so $B_0=k$, generated by $B_1$).

\exercise{5.1}{Let $X$ be a topological space and $\{X_i\}_i$ a covering of $X$ by closed subsets which is locally finite, that is, every $x\in X$ has an open neighborhood $U$ meeting only finitely many $X_i$. Show that $\dim X=\sup_i\dim X_i$.}
\pf
By Prop.~5.5(a), $\dim X_i\le\dim X$ for all $i$, so $\sup_i\dim X_i\le\dim X$.

Conversely, let $Z_0\subsetneq\dots\subsetneq Z_n$ be a chain of irreducible closed subsets of $X$. Pick $x\in Z_0$ (irreducible sets are non-empty) and an open neighborhood $U$ of $x$ meeting only $X_1,\dots,X_m$ (after renumbering). The set $Z_n\cap U$ is a non-empty open subset of the irreducible space $Z_n$ (it contains $x$), hence irreducible and dense in $Z_n$ (Prop.~4.5(a)). Since the $X_i$ cover $X$ and only $X_1,\dots,X_m$ meet $U$,
\[ Z_n\cap U=\bigcup_{i=1}^m\bigl((Z_n\cap U)\cap X_i\bigr), \]
a finite union of closed subsets of $Z_n\cap U$. By irreducibility $Z_n\cap U\subseteq X_i$ for some $i$. As $X_i$ is closed and $Z_n\cap U$ is dense in $Z_n$, we get $Z_n\subseteq X_i$. Then $Z_0\subsetneq\dots\subsetneq Z_n$ is a chain of irreducible closed subsets of $X_i$, so $n\le\dim X_i$. Hence $n\le\sup_i\dim X_i$ and, taking the supremum over chains, $\dim X\le\sup_i\dim X_i$. \hfill$\square$

\exercise{5.2}{Let $X$ be a scheme and $Z$ a closed subset of $X$. Show that for all $x\in X$ we have $\operatorname{codim}(\overline{\{x\}},X)=\dim\mathcal O_{X,x}$, and (for $Z\ne\emptyset$) $\operatorname{codim}(Z,X)=\min_{z\in Z}\dim\mathcal O_{X,z}$.}
\pf
(The codimension of the point $x$ means that of the irreducible closed set $\overline{\{x\}}$.)

\emph{First formula.} Choose an affine open $U=\Spec A\ni x$ and let $x$ correspond to the prime $\mathfrak p$, so $\mathcal O_{X,x}=A_{\mathfrak p}$. An irreducible closed subset $Z'$ of $X$ contains $\overline{\{x\}}$ iff $x\in Z'$ (as $Z'$ is closed). For such $Z'$ the set $Z'\cap U$ is a non-empty open subset of $Z'$, hence irreducible, dense in $Z'$, and closed in $U$. Conversely if $W$ is an irreducible closed subset of $U$ with $x\in W$, its closure $\overline W$ in $X$ is irreducible, closed, contains $x$, and $\overline W\cap U=W$. These two operations are mutually inverse and preserve inclusions: $Z'\mapsto Z'\cap U$ and $W\mapsto\overline W$ (for the first, $\overline{Z'\cap U}=Z'$ by density). So chains of irreducible closed subsets of $X$ containing $\overline{\{x\}}$ correspond to chains of irreducible closed subsets of $U$ containing $x$. By Prop.~4.7 (irreducible closed subsets of $\Spec A$ are the $V(\mathfrak q)$ with $\mathfrak q$ prime, and $x\in V(\mathfrak q)$ iff $\mathfrak q\subseteq\mathfrak p$) these are the chains of primes contained in $\mathfrak p$. Therefore $\operatorname{codim}(\overline{\{x\}},X)=\operatorname{ht}(\mathfrak p)=\dim A_{\mathfrak p}=\dim\mathcal O_{X,x}$ (Prop.~5.8(b)).

\emph{Second formula.} Let $Z_j$ ($j\in J$) be the irreducible components of $Z$, with generic points $z_j$ (Prop.~4.12(a)), so $Z_j=\overline{\{z_j\}}$. By definition and the first part, $\operatorname{codim}(Z,X)=\inf_j\operatorname{codim}(Z_j,X)=\inf_j\dim\mathcal O_{X,z_j}$. In particular $\operatorname{codim}(Z,X)\ge\inf_{z\in Z}\dim\mathcal O_{X,z}$. Conversely, let $z\in Z$. The irreducible closed set $\overline{\{z\}}\subseteq Z$ lies in some irreducible component $Z_j$ of $Z$ (every irreducible subset lies in a maximal one, by Zorn's lemma). Every chain of irreducible closed subsets of $X$ starting with a set containing $Z_j$ also starts with a set containing $\overline{\{z\}}$, so $\operatorname{codim}(Z_j,X)\le\operatorname{codim}(\overline{\{z\}},X)=\dim\mathcal O_{X,z}$. Hence $\operatorname{codim}(Z,X)\le\dim\mathcal O_{X,z}$ for every $z\in Z$. Thus $\operatorname{codim}(Z,X)=\inf_{z\in Z}\dim\mathcal O_{X,z}$, and the infimum is a minimum because the values lie in $\mathbb N\cup\{\infty\}$. \hfill$\square$

\exercise{5.3}{(a) Let $Z$ be a closed subset of a topological space $X$. Show that $\operatorname{codim}(Z,X)=0$ if and only if $Z$ contains an irreducible component of $X$. Give an example (with $X$ not irreducible) where $\operatorname{codim}(Z,X)=0$ and $\dim Z<\dim X$.

(b) Let $X=\Spec\mathcal O_K[T]$ where $\mathcal O_K$ is a discrete valuation ring with uniformizing parameter $t\ne0$. Let $f=tT-1$. Show that the ideal generated by $f$ is maximal. Let $x\in X$ be the corresponding point. Show that $X$ is irreducible, $\dim\mathcal O_{X,x}=1$, and $\operatorname{codim}(\{x\},X)+\dim\{x\}<\dim X$.}
\pf
(a) For an irreducible closed set $Y$, $\operatorname{codim}(Y,X)=0$ means there is no chain $Y\subseteq Z_0\subsetneq Z_1$ of irreducible closed sets, i.e. $Y$ is maximal among irreducible closed subsets of $X$, i.e. $Y$ is an irreducible component of $X$. (If $Y\subsetneq W$ with $W$ irreducible closed, take $Z_0=Y$, $Z_1=W$; conversely a chain gives such a $W$.) Since $\operatorname{codim}(Z,X)$ is the infimum of the $\operatorname{codim}(Z_j,X)$ over the components $Z_j$ of $Z$, and these are integers $\ge0$ or $\infty$, it is $0$ iff some component $Z_j$ of $Z$ is an irreducible component of $X$.

If $Z$ contains an irreducible component $C$ of $X$, then $C$ is an irreducible subset of $Z$, hence lies in a component $Z_j$ of $Z$. Since $Z_j$ is an irreducible closed subset of $X$ and $C$ is maximal, $C=Z_j$, so $\operatorname{codim}(Z,X)=0$. Conversely, if $\operatorname{codim}(Z,X)=0$, some component $Z_j$ of $Z$ is a component of $X$, so $Z$ contains one.

\emph{Example.} Let $X=\Spec k[x,y,z]/(xz,yz)=V(z)\cup V(x,y)\subseteq\mathbb A^3_k$, the union of the plane $P=V(z)\simeq\Spec k[x,y]$ and the line $L=V(x,y)\simeq\Spec k[z]$. These are irreducible closed (their ideals $(z)$, $(x,y)$ are prime), neither contains the other, so they are the irreducible components of $X$ (Prop.~4.5(c)). By Cor.~5.17, $\dim P=2$ and $\dim L=1$, so $\dim X=2$ (Prop.~5.5(c)). Take $Z=L$. It is an irreducible component of $X$, so $\operatorname{codim}(Z,X)=0$, while $\dim Z=1<2=\dim X$.

(b) Put $A=\mathcal O_K[T]$ and $K=\operatorname{Frac}\mathcal O_K$. Every non-zero element of $\mathcal O_K$ is a unit times a power of $t$, so $\mathcal O_K[1/t]=K$. The ring homomorphism $\psi:A\to K$ which is the inclusion on $\mathcal O_K$ and sends $T\mapsto1/t$ is surjective and has $f$ in its kernel. In $A/fA$ the element $t$ is invertible with inverse $\bar T$, so the structure map $\mathcal O_K\to A/fA$ extends to $K=\mathcal O_K[1/t]\to A/fA$; this map is surjective because $A/fA$ is generated over the image of $\mathcal O_K$ by $\bar T=t^{-1}$. As $A/fA$ surjects onto $K$ it is non-zero, and a ring homomorphism from the field $K$ to a non-zero ring is injective. So $K\simeq A/fA$ and $fA$ is maximal. Let $x$ be the corresponding point; it is a closed point.

$A$ is an integral domain, so $X$ is irreducible. The local ring $\mathcal O_{X,x}=A_{fA}$ has dimension $\operatorname{ht}(fA)$ (Prop.~5.8(b)). By Thm.~5.12 (Krull), $\operatorname{ht}(fA)\le1$; and $0\subsetneq fA$ since $f\neq0$ in the domain $A$, so $\operatorname{ht}(fA)=1$.

Finally $\dim\mathcal O_K=1$ (its primes are $0\subsetneq t\mathcal O_K$) and $\mathcal O_K$ is Noetherian, so $\dim X=\dim\mathcal O_K[T]=2$ (Cor.~5.17). Using Ex.~5.2, $\operatorname{codim}(\{x\},X)=\dim\mathcal O_{X,x}=1$, and $\dim\{x\}=0$ because it is a point. So $1+0<2$. \hfill$\square$

\exercise{5.4}{Let $A$ be a ring, $\mathfrak p_1,\dots,\mathfrak p_r$ prime ideals and $I$ an ideal contained in none of the $\mathfrak p_i$. Show $I\not\subseteq\mathfrak p_1\cup\dots\cup\mathfrak p_r$.

(a) Show the property for $r\le2$.

(b) Assume the property for $r-1$ and that $\mathfrak p_r$ does not contain any $\mathfrak p_i$, $i\le r-1$. Let $x\in I\setminus(\mathfrak p_1\cup\dots\cup\mathfrak p_{r-1})$. Show that there exists $y\in(I\mathfrak p_1\cdots\mathfrak p_{r-1})\setminus\mathfrak p_r$.

(c) Show that either $x$ or $x+y$ is not in $\mathfrak p_1\cup\dots\cup\mathfrak p_r$.}
\pf
(a) For $r=0$ the union is empty and $I\ni0$. For $r=1$ this is the hypothesis. For $r=2$: pick $a\in I\setminus\mathfrak p_1$ and $b\in I\setminus\mathfrak p_2$. If $a\notin\mathfrak p_2$ or $b\notin\mathfrak p_1$ we are done. Otherwise $a\in\mathfrak p_2\setminus\mathfrak p_1$ and $b\in\mathfrak p_1\setminus\mathfrak p_2$. Then $a+b\in I$, and $a+b\notin\mathfrak p_1$ (else $a\in\mathfrak p_1$) and $a+b\notin\mathfrak p_2$ (else $b\in\mathfrak p_2$). (Only the ideal property is used here.)

(b) The ideal $I\mathfrak p_1\cdots\mathfrak p_{r-1}$ is not contained in $\mathfrak p_r$: if it were, since $\mathfrak p_r$ is prime one of the factors would be contained in $\mathfrak p_r$, i.e. $I\subseteq\mathfrak p_r$ (excluded) or $\mathfrak p_i\subseteq\mathfrak p_r$ for some $i<r$ (excluded). Pick $y$ in the product but not in $\mathfrak p_r$. (The element $x$ exists by the property for $r-1$ applied to $\mathfrak p_1,\dots,\mathfrak p_{r-1}$.)

(c) Suppose $x\in\mathfrak p_r$ (otherwise $x\notin\mathfrak p_1\cup\dots\cup\mathfrak p_r$ and we are done). Then $x+y\in I$. For $i<r$ we have $y\in\mathfrak p_i$ and $x\notin\mathfrak p_i$, so $x+y\notin\mathfrak p_i$. For $i=r$ we have $x\in\mathfrak p_r$ and $y\notin\mathfrak p_r$, so $x+y\notin\mathfrak p_r$.

\emph{Conclusion (induction on $r$).} The case $r\le2$ is (a). Let $r\ge3$ and assume the result for fewer than $r$ primes (and every ideal). If $\mathfrak p_i\subseteq\mathfrak p_j$ for some $i\ne j$, we may delete $\mathfrak p_i$ without changing the union, and we conclude by induction. Otherwise no $\mathfrak p_i$ is contained in another, in particular $\mathfrak p_r$ contains no $\mathfrak p_i$ with $i<r$. By induction there is $x\in I\setminus(\mathfrak p_1\cup\dots\cup\mathfrak p_{r-1})$, and (b), (c) produce an element of $I$ outside $\mathfrak p_1\cup\dots\cup\mathfrak p_r$. \hfill$\square$

\exercise{5.5}{Let $A$ be a graded ring, $\mathfrak p_1,\dots,\mathfrak p_r$ homogeneous prime ideals and $I$ a homogeneous ideal contained in none of the $\mathfrak p_i$. Show that there is a homogeneous element of $I$ not in $\bigcup_i\mathfrak p_i$. (One may suppose $\mathfrak p_r$ contains no $\mathfrak p_i$, $i<r$.) (a) Show that there is a homogeneous element $a\in I\mathfrak p_1\cdots\mathfrak p_{r-1}\setminus\mathfrak p_r$. (b) Let $b\in I$ be homogeneous, $b\notin\bigcup_{i\le r-1}\mathfrak p_i$. Show that $b$ or $a^{\deg b}+b^{\deg a}$ is not in $\mathfrak p_1\cup\dots\cup\mathfrak p_r$. Conclude.}
\pf
\emph{The statement as printed is false if $I$ contains non-zero elements of degree $0$.} Let $A=k[s,t,x]$ with $\deg s=\deg t=0$, $\deg x=1$, and $\mathfrak p_1=(s)$, $\mathfrak p_2=(t,x)$ (homogeneous primes, neither contains the other), $I=(s,x)$. Then $I\not\subseteq\mathfrak p_1$ (as $x\notin(s)$) and $I\not\subseteq\mathfrak p_2$ (as $s\notin(t,x)$). But $I_0=s\,k[s,t]\subseteq\mathfrak p_1$ and $I_d=k[s,t]x^d\subseteq\mathfrak p_2$ for $d\ge1$, so every homogeneous element of $I$ lies in $\mathfrak p_1\cup\mathfrak p_2$.

The statement is true (and this is what the hint's proof gives) under the hypothesis $I\subseteq A_+=\bigoplus_{d>0}A_d$. This is automatic when $A_0$ is a field and $I\ne A$ (a proper homogeneous ideal has $I_0=0$), in particular for the homogeneous algebras $B$ over $k$ of Ex.~5.15; and if $I=A$ the element $1$ works. We prove it under $I\subseteq A_+$.

Induction on $r$; the cases $r=0$ (take $0$) and $r=1$ are clear (a homogeneous ideal not contained in $\mathfrak p_1$ has a homogeneous generator outside $\mathfrak p_1$). Let $r\ge2$. As in Ex.~5.4 we may delete any $\mathfrak p_i$ contained in another $\mathfrak p_j$, so we may assume no inclusions among the $\mathfrak p_i$ and, after renumbering, that $\mathfrak p_r$ contains no $\mathfrak p_i$ with $i<r$.

(a) The ideal $J=I\mathfrak p_1\cdots\mathfrak p_{r-1}$ is homogeneous, generated by products of homogeneous elements. As in Ex.~5.4(b), $J\not\subseteq\mathfrak p_r$. A homogeneous ideal is generated by its homogeneous elements, so some homogeneous $a\in J$ is not in $\mathfrak p_r$. Then $a\in I$, $a\in\mathfrak p_i$ for $i<r$, and $a\neq0$, so $\deg a\ge1$ since $I\subseteq A_+$.

(b) By the induction hypothesis (for $\mathfrak p_1,\dots,\mathfrak p_{r-1}$) there is a homogeneous $b\in I$ with $b\notin\mathfrak p_1\cup\dots\cup\mathfrak p_{r-1}$; $b\ne0$, so $\deg b\ge1$. If $b\notin\mathfrak p_r$ we are done. Suppose $b\in\mathfrak p_r$ and put $c=a^{\deg b}+b^{\deg a}$. Both summands are homogeneous of degree $\deg a\deg b$, and both lie in $I$ (exponents $\ge1$), so $c$ is a homogeneous element of $I$. For $i<r$: $a\in\mathfrak p_i$ gives $a^{\deg b}\in\mathfrak p_i$, while $b\notin\mathfrak p_i$ gives $b^{\deg a}\notin\mathfrak p_i$ (prime), so $c\notin\mathfrak p_i$. For $i=r$: $b\in\mathfrak p_r$ gives $b^{\deg a}\in\mathfrak p_r$, while $a\notin\mathfrak p_r$ gives $a^{\deg b}\notin\mathfrak p_r$, so $c\notin\mathfrak p_r$. Hence $b$ or $c$ is a homogeneous element of $I$ outside $\mathfrak p_1\cup\dots\cup\mathfrak p_r$. \hfill$\square$

\exercise{5.6}{Let $A$ be a Noetherian ring and $I$ an ideal.

(a) Let $x\in I$. Show $\operatorname{ht}(I/xA)+1\ge\operatorname{ht}(I)$, with equality if $x$ lies in no minimal prime ideal of $A$. Show that a prime minimal among those containing $x$ has height $1$ if $x$ is a regular element.

(b) Show $\operatorname{ht}(I)\ge1$ iff $I$ is not contained in the union of the minimal prime ideals of $A$.

(c) Let $r=\operatorname{ht}(I)$. Show that $I$ contains an ideal $J$ generated by $r$ elements with $\operatorname{ht}(J)=r$ and $\operatorname{ht}(I/J)=0$.}
\pf
We assume $I\ne A$ (for $I=A$ all heights are the infimum over the empty set). Recall $\operatorname{ht}(I)=\min\{\operatorname{ht}(P):P\supseteq I\text{ prime}\}$, and heights are finite in a Noetherian ring: $\operatorname{ht}(P)=\dim A_P<\infty$ by Cor.~5.14(b). Primes of $A/xA$ containing $I/xA$ are the $P/xA$ with $P\supseteq I$ prime (note $x\in I$), and $(A/xA)_{P/xA}=A_P/xA_P$.

(a) Choose a prime $P\supseteq I$ with $\operatorname{ht}(P/xA)=\operatorname{ht}(I/xA)$. Apply Thm.~5.15 to the Noetherian local ring $A_P$ and $x\in PA_P$: $\dim A_P/xA_P\ge\dim A_P-1$, that is $\operatorname{ht}(P/xA)\ge\operatorname{ht}(P)-1$. Hence $\operatorname{ht}(I/xA)+1\ge\operatorname{ht}(P)\ge\operatorname{ht}(I)$.

Suppose $x$ lies in no minimal prime of $A$. Choose a prime $P\supseteq I$ with $\operatorname{ht}(P)=\operatorname{ht}(I)$. The minimal primes of $A_P$ are the $QA_P$ with $Q\subseteq P$ a minimal prime of $A$, and $x\notin Q$, so $x/1$ lies in no minimal prime of $A_P$. By the equality case of Thm.~5.15, $\dim A_P/xA_P=\dim A_P-1$, i.e. $\operatorname{ht}(P/xA)=\operatorname{ht}(I)-1$. Thus $\operatorname{ht}(I/xA)\le\operatorname{ht}(P/xA)=\operatorname{ht}(I)-1$, and with the first part $\operatorname{ht}(I/xA)+1=\operatorname{ht}(I)$.

Now let $x$ be regular (a non-zero-divisor). First, a minimal prime $\mathfrak q$ of $A$ consists of zero divisors: $\mathfrak qA_{\mathfrak q}$ is the only prime of $A_{\mathfrak q}$, so it is the nilradical; for $u\in\mathfrak q$ there are $s\notin\mathfrak q$ and $n\ge1$ with $su^n=0$. Take $n$ minimal; then $s\ne0$ and $u\cdot(su^{n-1})=0$ with $su^{n-1}\ne0$ (for $n=1$ this is $s\ne0$), so $u$ is a zero divisor. Hence the regular element $x$ is in no minimal prime. Let $p$ be minimal among the primes containing $x$. By Thm.~5.12, $\operatorname{ht}(p)\le1$, and $\operatorname{ht}(p)\ne0$ since $p$ is not a minimal prime of $A$. So $\operatorname{ht}(p)=1$.

(b) $\operatorname{ht}(I)=0$ iff some prime $P\supseteq I$ has height $0$, i.e. is minimal, i.e. iff $I$ is contained in a minimal prime. By Prop.~4.9 there are finitely many minimal primes $\mathfrak q_1,\dots,\mathfrak q_s$. If $\operatorname{ht}(I)\ge1$, then $I\not\subseteq\mathfrak q_j$ for all $j$, so $I\not\subseteq\bigcup_j\mathfrak q_j$ by Ex.~5.4. Conversely if $I\not\subseteq\bigcup_j\mathfrak q_j$, then $I$ is in no minimal prime, so $\operatorname{ht}(I)\ge1$.

(c) We construct $x_1,\dots,x_r\in I$ such that $J_i=(x_1,\dots,x_i)$ has $\operatorname{ht}(J_i)=i$ for $0\le i\le r$, starting with $J_0=0$ (height $0$ since $A\neq0$). Let $0\le i<r$ and suppose $J_i$ is built. The ideal $J_i$ has finitely many minimal primes (Prop.~4.9 applied to $A/J_i$), all of height $\ge\operatorname{ht}(J_i)=i$; let $P_1,\dots,P_t$ be those of height exactly $i$ (there is at least one). Since $\operatorname{ht}(I)=r>i$, $I\not\subseteq P_j$. By Ex.~5.4 there is $x_{i+1}\in I\setminus\bigcup_jP_j$. Put $J_{i+1}=J_i+x_{i+1}A$. Let $Q\supseteq J_{i+1}$ be prime. Then $Q$ contains a minimal prime $P$ of $J_i$. If $P$ is one of the $P_j$, then $Q\ne P$ because $x_{i+1}\in Q\setminus P$, so $\operatorname{ht}(Q)\ge i+1$. Otherwise $\operatorname{ht}(Q)\ge\operatorname{ht}(P)\ge i+1$. So $\operatorname{ht}(J_{i+1})\ge i+1$, and $\le i+1$ by Cor.~5.14(a). Hence $\operatorname{ht}(J_{i+1})=i+1$.

Put $J=J_r\subseteq I$: generated by $r$ elements, $\operatorname{ht}(J)=r$. Let $P\supseteq I$ be a prime with $\operatorname{ht}(P)=r$. Then $P\supseteq J$, so it contains a minimal prime $Q$ of $J$ with $\operatorname{ht}(Q)\ge\operatorname{ht}(J)=r$. If $Q\subsetneq P$ then $\operatorname{ht}(P)\ge\operatorname{ht}(Q)+1>r$, impossible. So $P=Q$ is a minimal prime of $J$ containing $I$. Then $P/J$ is a minimal prime of $A/J$ containing $I/J$, i.e. $\operatorname{ht}(I/J)=0$. \hfill$\square$

\exercise{5.7}{Let $f:X\to Y$ be a morphism of algebraic varieties with $f(X)$ dense in $Y$. Show that $\dim X\ge\dim Y$ (use Prop.~5.19).}
\pf
Algebraic varieties are Noetherian (Rem.~3.48), so $X$ and $Y$ have finitely many irreducible components $X_1,\dots,X_m$ and $Y_1,\dots,Y_s$ (Prop.~4.9), with generic points $\xi_i$, $\eta_j$ (Prop.~4.12(a)). By Prop.~5.5(c) it suffices to show that for each $j$ there is an $i$ with $\dim Y_j\le\dim X_i$.

Since $f(X)=\bigcup_if(X_i)$, we have $Y=\overline{f(X)}=\bigcup_i\overline{f(X_i)}$, and each $\overline{f(X_i)}=\overline{\{f(\xi_i)\}}$ is irreducible and closed (as $f(X_i)=f(\overline{\{\xi_i\}})\subseteq\overline{\{f(\xi_i)\}}$ by continuity). The irreducible set $Y_j$ is covered by these finitely many closed sets, so $Y_j\subseteq\overline{\{f(\xi_i)\}}$ for some $i$; the latter is irreducible closed, hence contained in a component of $Y$, and by maximality of $Y_j$ we get $Y_j=\overline{\{f(\xi_i)\}}$. By uniqueness of the generic point (Prop.~4.12(a)), $f(\xi_i)=\eta_j$.

Give $X_i$ and $Y_j$ their reduced structures; they are integral algebraic varieties over $k$ (closed subschemes of varieties), with function fields the residue fields $k(\xi_i)$ and $k(\eta_j)$ at the generic points. The morphism $f$ is a morphism of $k$-schemes, so it induces a $k$-homomorphism of residue fields $k(\eta_j)=k(f(\xi_i))\to k(\xi_i)$, which is injective. Thus $k(\eta_j)$ is isomorphic over $k$ to a subfield of $k(\xi_i)$, and a transcendence basis of $k(\eta_j)$ over $k$ gives algebraically independent elements of $k(\xi_i)$. So $\operatorname{trdeg}_kk(\eta_j)\le\operatorname{trdeg}_kk(\xi_i)$. By Prop.~5.19, $\dim Y_j=\operatorname{trdeg}_kk(\eta_j)\le\operatorname{trdeg}_kk(\xi_i)=\dim X_i\le\dim X$. \hfill$\square$

\exercise{5.8}{Let $\mathcal O_K$ be as in Ex.~5.3, $K=\operatorname{Frac}(\mathcal O_K)$, $\kappa=\mathcal O_K/t\mathcal O_K$ the residue field. Let $A=K\times\kappa$ and $\varphi:\mathcal O_K\to A$ induced by $\mathcal O_K\to K$ and $\mathcal O_K\to\kappa$. Show that $f_\varphi:\Spec A\to\Spec\mathcal O_K$ is surjective and $\dim\mathcal O_K>\dim A$. Also show that $A$ is a finitely generated $\mathcal O_K$-algebra (a quotient of a polynomial ring over $\mathcal O_K$).}
\pf
The ring $A=K\times\kappa$ has exactly two prime ideals, $\mathfrak a_1=0\times\kappa$ (kernel of the projection to $K$) and $\mathfrak a_2=K\times0$ (kernel of the projection to $\kappa$). Indeed, if $\mathfrak a$ is prime then, with $e_1=(1,0)$, $e_2=(0,1)$, we have $e_1e_2=0\in\mathfrak a$ and $e_1+e_2=1\notin\mathfrak a$, so exactly one $e_i$ lies in $\mathfrak a$, and then $\mathfrak a$ is the preimage of a prime of the field $K$ or $\kappa$, i.e. $\mathfrak a_1$ or $\mathfrak a_2$. Neither contains the other, so $\dim A=0$. On the other hand $\dim\mathcal O_K=1$.

Now $\varphi^{-1}(\mathfrak a_1)=\ker(\mathcal O_K\to K)=0$ and $\varphi^{-1}(\mathfrak a_2)=\ker(\mathcal O_K\to\kappa)=t\mathcal O_K$. The spectrum of $\mathcal O_K$ is $\{0,t\mathcal O_K\}$, so $f_\varphi$ is surjective. (This shows that the integrality hypothesis in Prop.~5.10 cannot be dropped, and that Ex.~5.7 fails for schemes: $\Spec A\to\Spec\mathcal O_K$ has dense image, and $\dim\Spec A=0<1$.)

Finally, let $u=(t^{-1},0)\in A$ and consider the $\mathcal O_K$-algebra map $\theta:\mathcal O_K[T]\to A$, $T\mapsto u$. For $p=\sum_{i}c_iT^i$ we get $\theta(p)=\bigl(\sum_ic_it^{-i},\ \bar c_0\bigr)$. Given $(\alpha,\beta)\in K\times\kappa$, lift $\beta$ to $c_0\in\mathcal O_K$. Since $\mathcal O_K[1/t]=K$, write $\alpha-c_0=\sum_{i=0}^md_it^{-i}$ with $d_i\in\mathcal O_K$; then $d_0=(td_0)t^{-1}$, so $\alpha-c_0=\sum_{i=1}^me_it^{-i}$ with $e_1=d_1+td_0$ and $e_i=d_i$ for $i\ge2$ (all in $\mathcal O_K$). Hence $p=c_0+\sum_{i\ge1}e_iT^i$ satisfies $\theta(p)=(\alpha,\beta)$. So $\theta$ is surjective and $A$ is a finitely generated (indeed monogenic) $\mathcal O_K$-algebra. \hfill$\square$

\exercise{5.9}{Let $X$ be a scheme over a field $k$. Show that the points of $X(k)$ are closed in $X$. If $X$ is an algebraic variety over $k$, show that $x\in X$ is closed iff $k(x)$ is algebraic over $k$.}
\pf
Here $X(k)$ is the set of points $x$ with $k(x)=k$ (the images of the $k$-morphisms $\Spec k\to X$).

\emph{Claim: if $x\in X$ and $k(x)$ is algebraic over $k$, then $x$ is closed in $X$ (for any $k$-scheme $X$).} Let $V=\Spec B$ be an affine open subset of $X$ containing $x$, and let $x$ correspond to the prime $\mathfrak p$ of $B$. Then $k\subseteq B/\mathfrak p\subseteq k(x)$, so $B/\mathfrak p$ is an integral domain which is algebraic over $k$, hence a field: for $0\ne b\in B/\mathfrak p$, the ring $k[b]$ is a domain which is a finite dimensional $k$-vector space, so it is a field and $b$ is invertible in $k[b]$. Thus $\mathfrak p$ is maximal, and $\{x\}$ is closed in $V$. For an affine open $V$ not containing $x$, $\{x\}\cap V=\emptyset$ is closed in $V$. Since the affine opens cover $X$, the complement of $\{x\}$ is open, i.e. $\{x\}$ is closed.

Points of $X(k)$ have $k(x)=k$, which is algebraic over $k$, so they are closed. This also proves ($\Leftarrow$) in the variety case.

($\Rightarrow$) Let $X$ be an algebraic variety, $x$ closed. Choose an affine open $V=\Spec B\ni x$ with $B$ a finitely generated $k$-algebra (Def.~3.47), $x\leftrightarrow\mathfrak p$. Since $\{x\}$ is closed in $X$ it is closed in $V$, so $\mathfrak p=\mathfrak m$ is maximal (Lemma~4.11(a): $\overline{\{x\}}=V(\mathfrak p)$). Then $k(x)=B/\mathfrak m$ is a field which is a finitely generated $k$-algebra, hence a finite algebraic extension of $k$ (Cor.~1.12). \hfill$\square$

\exercise{5.10}{Let $k$ be an uncountable field. Let $X$ be an algebraic variety over $k$ with $\dim X\ge1$ and $(Y_n)_n$ a sequence of closed subsets with $\dim Y_n<\dim X$. Show $\bigcup_nY_n\ne X$.

(a) Show the result for $X=\mathbb A^1_k$, then for $\mathbb A^m_k$. (b) Using Noether normalization, show the result for an arbitrary affine variety. Deduce the general case.}
\pf
(a) \emph{Claim: for $m\ge1$ and non-zero polynomials $g_n\in k[T_1,\dots,T_m]$ ($n\in\mathbb N$) there is $a\in k^m$ with $g_n(a)\ne0$ for all $n$.} For $m=1$ each $g_n$ has finitely many roots, so the set of all roots of all $g_n$ is countable, while $k$ is not; pick $a\in k$ outside it. For $m\ge2$, write $g_n=\sum_jh_{n,j}(T_1,\dots,T_{m-1})T_m^j$ and choose $j_n$ with $h_n:=h_{n,j_n}\ne0$. By induction there is $a'\in k^{m-1}$ with $h_n(a')\ne0$ for all $n$. Then the one-variable polynomials $g_n(a',T_m)$ are non-zero, and the case $m=1$ gives $a_m\in k$ with $g_n(a',a_m)\ne0$ for all $n$.

Now let $Y_n\subseteq\mathbb A^m_k$ be closed with $\dim Y_n<m=\dim\mathbb A^m_k$ (Cor.~5.17), so $Y_n\ne\mathbb A^m_k$. Write $Y_n=V(J_n)$; $J_n\ne0$ (as $V(0)$ is everything), so choose $0\ne g_n\in J_n$; then $Y_n\subseteq V(g_n)$. With $a$ as in the claim, the maximal ideal $\mathfrak m_a=(T_i-a_i)_i$ is not in any $V(g_n)$, since $g_n\in\mathfrak m_a\iff g_n(a)=0$. So $\mathfrak m_a\notin\bigcup_nY_n$.

(b) \emph{Affine case.} Let $X=\Spec A$, $A$ a finitely generated $k$-algebra, $\dim X=d\ge1$ (so $A\ne0$). By Prop.~1.9 there is a finite injective homomorphism $\phi:A_0=k[T_1,\dots,T_e]\to A$. By Prop.~5.10(b) and Cor.~5.17, $\dim A=\dim A_0=e$, so $e=d$. Let $\pi:X\to\mathbb A^d_k$ be the induced morphism, $Y_n=V(J_n)$ and $Z_n=V(\phi^{-1}(J_n))\subseteq\mathbb A^d_k$. We claim $\pi(Y_n)=Z_n$. The inclusion $\subseteq$ is clear. For $\supseteq$: the induced map $A_0/\phi^{-1}(J_n)\to A/J_n$ is injective and finite, so by Prop.~5.10(b) $\Spec(A/J_n)\to\Spec(A_0/\phi^{-1}J_n)$ is surjective, i.e. every prime of $A_0$ containing $\phi^{-1}(J_n)$ is the inverse image of a prime of $A$ containing $J_n$. Moreover Prop.~5.10(b) and Prop.~5.8(a) give $\dim Z_n=\dim A_0/\phi^{-1}(J_n)=\dim A/J_n=\dim Y_n<d$. By (a) there is a point $z\in\mathbb A^d_k$ outside $\bigcup_nZ_n$. Since $\phi$ is injective and integral, $\pi$ is surjective (Prop.~5.10(b)); take $x\in X$ with $\pi(x)=z$. If $x\in Y_n$ then $z\in\pi(Y_n)=Z_n$, a contradiction. So $x\notin\bigcup_nY_n$.

\emph{General case.} $X$ is covered by finitely many affine open subvarieties $U_1,\dots,U_s$, and $\dim X=\max_i\dim U_i$ (Rem.~5.6), so there is an affine open $U=U_i$ with $\dim U=\dim X=:d\ge1$. The sets $Y_n\cap U$ are closed in $U$ with $\dim(Y_n\cap U)\le\dim Y_n<d$ (Prop.~5.5(a)). By the affine case some point of $U$ is outside all $Y_n\cap U$, hence outside $\bigcup_nY_n$. \hfill$\square$

\exercise{5.11}{(Schemes of dimension $0$.)

(a) Let $X$ be a scheme which is a (finite or not) disjoint union of open subschemes $X_i$. Show $\mathcal O_X(X)\simeq\prod_i\mathcal O_X(X_i)$.

(b) Show that any scheme of finite cardinality and dimension $0$ is affine.

(c) Let $X=\Spec A$ be a scheme of finite cardinality and dimension $0$. Show that every point $x\in X$ is open. Deduce that $A\simeq\bigoplus_{\mathfrak p\in\Spec A}A_{\mathfrak p}$.

(d) Show that (c) is false if we do not assume $\Spec A$ of dimension $0$.}
\pf
(a) The $X_i$ form an open cover of $X$ with $X_i\cap X_j=\emptyset$ for $i\ne j$, and $\mathcal O_X(\emptyset)=0$. By the sheaf axioms, $s\mapsto(s|_{X_i})_i$ is injective (a section with all restrictions zero is zero), and any family $(s_i)_i$ with $s_i\in\mathcal O_X(X_i)$ is compatible (restrictions to $X_i\cap X_j=\emptyset$ are both $0$), so it glues to a section of $\mathcal O_X(X)$. Hence the ring homomorphism $\mathcal O_X(X)\to\prod_i\mathcal O_X(X_i)$ is an isomorphism.

(b),(c) Let $X$ be a scheme with finitely many points and $\dim X=0$. \emph{Every point is closed:} let $y\in\overline{\{x\}}$. Then $\overline{\{y\}}\subseteq\overline{\{x\}}$ are irreducible closed sets, hence equal because $\dim X=0$. So $x$ and $y$ have the same closure, hence the same open neighborhoods (an open set contains $x$ iff it meets $\overline{\{x\}}$). By the $T_0$ property of schemes (Ex.~4.7), $x=y$. So $\overline{\{x\}}=\{x\}$. As $X$ is finite, every subset is a finite union of closed points, hence closed; so every subset is also open, and $X$ is discrete. This proves the first assertion of (c) (in particular for $X=\Spec A$).

Each one-point open subscheme $\{x\}$ is affine: choose an affine open $U=\Spec A'\ni x$; the non-empty open subset $\{x\}$ of $U$ contains some non-empty principal open $D(g)$, which must equal $\{x\}$; hence $\{x\}=D(g)\simeq\Spec A'_g$. So $X=\{x_1\}\sqcup\dots\sqcup\{x_n\}$ with $\{x_i\}=\Spec A_i$. Put $A=\prod_iA_i$ with idempotents $e_i$. A prime ideal of $A$ omits exactly one of the $e_i$: at most one, because $e_ie_j=0$ for $i\ne j$, and at least one, because $\sum_ie_i=1$. Hence $\Spec A$ is the disjoint union of the open sets $D(e_i)=\Spec A_{e_i}=\Spec A_i$. The morphism from $\bigsqcup_i\Spec A_i$ to $\Spec A$ which is, on each piece, the open immersion onto $D(e_i)$ is therefore an isomorphism. So $X\simeq\Spec A$ is affine, proving (b).

For the second assertion of (c), apply (a) to $X=\bigsqcup_{\mathfrak p}\{\mathfrak p\}$: $A=\mathcal O_X(X)\simeq\prod_{\mathfrak p}\mathcal O_X(\{\mathfrak p\})$. Since $\{\mathfrak p\}$ is the smallest open neighborhood of $\mathfrak p$, $\mathcal O_X(\{\mathfrak p\})=\mathcal O_{X,\mathfrak p}=A_{\mathfrak p}$, the restriction map being $a\mapsto a/1$. As the product is finite it equals the direct sum: $A\simeq\bigoplus_{\mathfrak p}A_{\mathfrak p}$, via $a\mapsto(a/1)_{\mathfrak p}$.

(d) Let $A=\mathcal O_K$ be a discrete valuation ring (e.g. $\mathbb Z_{(p)}$). Then $\Spec A=\{0,\mathfrak m\}$ is finite of dimension $1$. The map $A\to A_{(0)}\times A_{\mathfrak m}=K\times A$, $a\mapsto(a,a)$, is not surjective (its image is the diagonal, which does not contain $(1,0)$); indeed $A$ is a domain, so it is not a non-trivial product at all. Also the closed point $\mathfrak m$ is not open. So (c) fails. \hfill$\square$

\exercise{5.12}{Let $k$ be a field. We determine the affine open subsets of $\mathbb A^n_k$ and $\mathbb P^n_k$.

(a) Show that the principal open subsets of $\mathbb A^n_k$ and of $\mathbb P^n_k$ (not equal to $\mathbb P^n_k$) are affine.

(b) Let $X=\bigcup_iD(f_i)$ be a finite union of principal open subsets of $\mathbb A^n_k$. Show $\mathcal O_{\mathbb A^n_k}(X)=k[T_1,\dots,T_n]_f$, where $f=\gcd\{f_i\}$. Show that every affine open subset of $\mathbb A^n_k$ is principal.

(c) Show that every irreducible closed subset of $\mathbb P^n_k$ of dimension $n-1$ is principal.

(d) Let $X$ be an affine open subset of $\mathbb P^n_k$. Show that the irreducible components of $\mathbb P^n_k\setminus X$ have dimension $n-1$, and that $X$ is a principal open subset of $\mathbb P^n_k$.}
\pf
Let $R=k[T_1,\dots,T_n]$, a UFD with fraction field $K$. In $\mathbb P^n_k=\operatorname{Proj}B$, $B=k[T_0,\dots,T_n]$, the principal open sets are the $D_+(F)$ with $F$ homogeneous (those different from $\mathbb P^n_k$ have $\deg F\ge1$ or $F=0$), and a principal closed subset is read as a set $V_+(F)$ with $F$ homogeneous (global sections of $\mathcal O_{\mathbb P^n_k}$ are only the constants, so this is the meaning in $\mathbb P^n_k$).

(a) $D(f)\subseteq\mathbb A^n_k$ is isomorphic to $\Spec R_f$. For $F\in B$ homogeneous of positive degree, $D_+(F)\simeq\Spec B_{(F)}$ (Lemma~3.36 and the construction of the scheme structure on $\operatorname{Proj}B$, Prop.~3.38). Both are affine. (If $F=0$ the open set is empty, which is affine.)

(b) We may drop the $f_i$ equal to $0$ (as $D(0)=\emptyset$); if none is left, $X=\emptyset$, and $f=0$, $R_0=0$. So assume some $f_i\ne0$; then $f=\gcd\{f_i\}$ is defined up to a unit. Since $\mathbb A^n_k$ is integral, Prop.~4.18(c) gives, inside $K$,
\[ \mathcal O(X)=\bigcap_{x\in X}\mathcal O_{x}=\bigcap_i\ \bigcap_{x\in D(f_i)}\mathcal O_x=\bigcap_i\mathcal O(D(f_i))=\bigcap_iR_{f_i}. \]
We show $\bigcap_iR_{f_i}=R_f$. Since $f\mid f_i$, inverting $f_i$ inverts $f$, so $R_{f_i}\supseteq R_f$. Conversely let $u\in\bigcap_iR_{f_i}$ and write $u=a/b$ with $a,b\in R$ coprime. For each $i$, $u=c/f_i^m$ for some $c\in R$, $m\ge0$, so $af_i^m=bc$, hence $b\mid af_i^m$ and so $b\mid f_i^m$ (coprimality). Thus every prime factor $\pi$ of $b$ divides every $f_i$, hence divides $f$. Write $b=\epsilon\prod_j\pi_j^{e_j}$ with $\epsilon$ a unit; with $N=\max e_j$ we get $b\mid f^N$, so $u=a(f^N/b)/f^N\in R_f$. This proves $\mathcal O(X)=R_f$.

Now let $X\subseteq\mathbb A^n_k$ be an affine open subset. If $X=\emptyset$ then $X=D(0)$. Otherwise $X=\mathbb A^n_k\setminus V(I)$ with $I=(f_1,\dots,f_s)$ ($R$ is Noetherian), so $X=\bigcup D(f_i)$ and $\mathcal O(X)=R_f$ with $f=\gcd\{f_i\}$. Since $f\mid f_i$, $D(f_i)\subseteq D(f)$, so $X\subseteq D(f)$, and the restriction map $\mathcal O(D(f))=R_f\to\mathcal O(X)=R_f$ is the identity (both are subrings of $K$). The open immersion $j:X\to D(f)=\Spec R_f$ is a morphism of affine schemes, so by Prop.~3.25 it is the morphism induced by the ring homomorphism $\mathcal O(D(f))\to\mathcal O(X)$, which is an isomorphism; hence $j$ is an isomorphism, and in particular surjective. So $X=D(f)$ is principal.

(c) Let $Z\subseteq\mathbb P^n_k$ be irreducible closed of dimension $n-1$. Then $Z=V_+(\mathfrak p)$ for a homogeneous prime $\mathfrak p\not\supseteq B_+$ (the usual correspondence between irreducible closed subsets of $\operatorname{Proj}B$ and such primes), and $Z$ is homeomorphic to $\operatorname{Proj}(B/\mathfrak p)$. As $\mathfrak p\subseteq B_+$ the quotient $B/\mathfrak p$ is a homogeneous $k$-algebra, and Cor.~5.21 gives $\dim B/\mathfrak p=\dim Z+1=n$. Since $\dim B=n+1$ (Cor.~5.17) and $B$ is a finitely generated domain, Prop.~5.23(a) gives $\operatorname{ht}(\mathfrak p)=(n+1)-n=1$. Take a non-zero $g\in\mathfrak p$ and an irreducible factor $\pi$ of $g$ lying in $\mathfrak p$ ($\mathfrak p$ is prime). Then $(\pi)$ is a non-zero prime of the UFD $B$ with $0\subsetneq(\pi)\subseteq\mathfrak p$, so $\mathfrak p=(\pi)$ because $\operatorname{ht}(\mathfrak p)=1$. The generator $\pi$ is homogeneous: let $\pi_a\ne0$ be its lowest homogeneous component; since $\mathfrak p$ is homogeneous, $\pi_a\in\mathfrak p=(\pi)$, so $\pi_a=\pi h$ with $h\ne0$. Comparing total degrees in the domain $B$, $a=\deg\pi_a=\deg\pi+\deg h\ge\deg\pi\ge a$, so $\deg\pi=a$ and $\pi=\pi_a$. Hence $Z=V_+(\pi)$ with $\pi$ homogeneous, i.e. $Z$ is principal.

(d) Let $n\ge1$ and let $X\ne\emptyset$ be an affine open subset. (For $X=\emptyset$ the first assertion is false, since the complement is $\mathbb P^n_k$ of dimension $n$, but $X=D_+(0)$ is principal anyway; for $n=0$ everything is trivial.) Put $Z=\mathbb P^n_k\setminus X$. Then $Z\ne\emptyset$: otherwise $X=\mathbb P^n_k$ would be affine with $\mathcal O(\mathbb P^n_k)=k$, i.e. a point, which is absurd for $n\ge1$. Let $W_i=D_+(T_i)\simeq\mathbb A^n_k$. The scheme $\mathbb P^n_k$ is separated (Example~3.3.7, which works over $k$ as over $\mathbb Z$), so the intersection of the affine open sets $X$ and $W_i$ is affine (Prop.~3.3.6). By (b), $X\cap W_i=D(g_i)$ for some $g_i\in\mathcal O(W_i)$. As $\mathbb P^n_k$ is irreducible and $X\ne\emptyset$, $X\cap W_i\ne\emptyset$ (non-empty opens meet), so $g_i\ne0$, and $g_i$ is not nilpotent in the domain $\mathcal O(W_i)$. Hence $Z\cap W_i=V(g_i)$ is either empty or, by Cor.~5.26, every irreducible component of it has dimension $n-1$.

Let $Z_1$ be an irreducible component of $Z$. It meets some $W_i$, and $Z_1\cap W_i$ is an irreducible component of $Z\cap W_i$ (Prop.~4.5(b)). So $\dim(Z_1\cap W_i)=n-1$, and since a non-empty open subset of the irreducible variety $Z_1$ (reduced structure) has the same dimension (Prop.~5.19), $\dim Z_1=n-1$.

Let $Z_1,\dots,Z_s$ be the (finitely many) components of $Z$. By (c), $Z_j=V_+(F_j)$ with $F_j$ homogeneous of positive degree. With $F=F_1\cdots F_s$ we get $Z=\bigcup_jV_+(F_j)=V_+(F)$ and $X=\mathbb P^n_k\setminus V_+(F)=D_+(F)$, a principal open subset. \hfill$\square$

\exercise{5.13}{Let $k$ be a field. A function field in $n$ variables over $k$ is a field $K$ finitely generated over $k$ with $\operatorname{trdeg}_kK=n$. Show that an extension $K$ of $k$ is a function field in $n$ variables over $k$ if and only if $K=K(X)$ for an integral algebraic variety $X$ over $k$ of dimension $n$. If so, $X$ can be taken projective.}
\pf
($\Leftarrow$) Let $X$ be an integral algebraic variety of dimension $n$, with a non-empty affine open $U=\Spec A$, $A$ a finitely generated $k$-algebra which is a domain. Then $K(X)=\mathcal O_{X,\xi}=\operatorname{Frac}(A)$ (Prop.~4.18(a)), which is generated as a field by the finitely many generators of $A$ over $k$, and $\operatorname{trdeg}_kK(X)=\dim X=n$ by Prop.~5.19.

($\Rightarrow$) Let $K=k(f_1,\dots,f_r)$ with $\operatorname{trdeg}_kK=n$. Let $\mathfrak q$ be the kernel of $k[T_1,\dots,T_r]\to K$, $T_i\mapsto f_i$, and $A=k[T_1,\dots,T_r]/\mathfrak q\simeq k[f_1,\dots,f_r]\subseteq K$. It is a domain with $\operatorname{Frac}(A)=K$. Let $C=V(\mathfrak q)\simeq\Spec A$, a closed irreducible subset of $\mathbb A^r_k=D_+(T_0)\subseteq\mathbb P^r_k$. Let $X$ be the closure of $C$ in $\mathbb P^r_k$ with its reduced induced closed subscheme structure (Section 2.4): it is irreducible and reduced, hence integral, and it is a closed subscheme of $\mathbb P^r_k$, so a projective (hence algebraic) variety. The intersection $X\cap D_+(T_0)$ is the closure of $C$ in $D_+(T_0)$, which is $C$ (closed). The reduced closed subscheme of $\mathbb A^r_k$ with underlying set $C$ is $\Spec k[T]/\mathfrak q=\Spec A$ since $\mathfrak q$ is prime, so $X\cap D_+(T_0)\simeq\Spec A$ is a non-empty open subscheme of $X$. By Prop.~4.18(a), $K(X)=\operatorname{Frac}(A)=K$, and by Prop.~5.19, $\dim X=\operatorname{trdeg}_kK=n$. \hfill$\square$

\exercise{5.14}{Let $L/K$ be a finite field extension. For $x\in L$, multiplication by $x$ is a $K$-linear endomorphism of $L$; $\operatorname{Norm}_{L/K}(x)$ is its determinant.

(a) Show that $\operatorname{Norm}_{L/K}$ is a multiplicative map $L\to K$.

(b) Let $A\subseteq B$ be rings with $K=\operatorname{Frac}(A)$, $L=\operatorname{Frac}(B)$, $B$ integral over $A$ ($A,B$ domains, $L/K$ finite). Show that for any $b\in B$, $\operatorname{Norm}_{L/K}(b)$ is integral over $A$.

(c) Suppose moreover that $A$ is a polynomial ring over a field $k$. Show $\operatorname{Norm}_{L/K}(B)\subseteq A$.}
\pf
(a) Let $m_x:L\to L$ denote multiplication by $x$. Then $m_{xy}=m_x\circ m_y$, so $\operatorname{Norm}(xy)=\det(m_x)\det(m_y)=\operatorname{Norm}(x)\operatorname{Norm}(y)$; also $m_1=\operatorname{id}$ gives $\operatorname{Norm}(1)=1$.

(b) Let $b\in B$ and let $\mu\in K[T]$ be its minimal polynomial over $K$, of degree $d=[K(b):K]$, and $e=[L:K(b)]$. Choosing the $K$-basis $1,b,\dots,b^{d-1}$ of $K(b)$ and a $K(b)$-basis of $L$ with $e$ elements, the matrix of $m_b$ is block diagonal with $e$ copies of the companion matrix of $\mu$. So its characteristic polynomial is $\chi(T)=\mu(T)^e$, and $\operatorname{Norm}(b)=\det(m_b)=(-1)^{de}\chi(0)=(-1)^{de}\mu(0)^e$. Let $\Omega$ be an algebraic closure of $K$ and $\beta_1,\dots,\beta_d\in\Omega$ the roots of $\mu$ (with multiplicity), so $\mu(0)=(-1)^d\prod_j\beta_j$ and $\operatorname{Norm}(b)=\pm(\prod_j\beta_j)^e$. Since $b$ is integral over $A$, it is a root of a monic $P\in A[T]$; $\mu$ divides $P$ in $K[T]$, so each $\beta_j$ is a root of $P$ and thus integral over $A$. The elements of $\Omega$ integral over $A$ form a ring, so $\operatorname{Norm}(b)$, a product of $\beta_j$'s up to sign, is integral over $A$.

(c) A polynomial ring $A$ over $k$ is a UFD, hence integrally closed in $K$: if $u=p/q$ with $p,q$ coprime satisfies $u^m+a_{m-1}u^{m-1}+\dots+a_0=0$ with $a_i\in A$, then $p^m=-q(a_{m-1}p^{m-1}+\dots+a_0q^{m-1})$, so any prime factor of $q$ divides $p^m$, hence $p$, a contradiction; so $q$ is a unit and $u\in A$. By (b), and since $\operatorname{Norm}(b)\in K$, we conclude $\operatorname{Norm}_{L/K}(b)\in A$ for all $b\in B$. \hfill$\square$

\exercise{5.15}{Let $X=\operatorname{Proj}B$ be a projective variety over a field $k$.

(a) Let $f\in B_+$ be homogeneous and $\mathfrak p$ a prime ideal of $B$ minimal among those containing $f$. Show that $\mathfrak p$ is homogeneous, contained in $B_+$, and $\operatorname{ht}(\mathfrak p)\le1$. Show that if $V_+(f)=\emptyset$ then $B_+=\mathfrak p$ and $\dim X\le0$.

(b) Let $Y$ be a closed subvariety of $X$ of dimension $r$. Show that for any homogeneous $f_1,\dots,f_r\in B_+$, $V_+(f_1,\dots,f_r)\cap Y\ne\emptyset$.}
\pf
Throughout, $B$ is a homogeneous $k$-algebra: $B=\bigoplus_{d\ge0}B_d$, $B_0=k$, generated by finitely many elements of $B_1$; in particular $B$ is Noetherian and $B_+$ is a maximal ideal. Recall that the irreducible closed subsets of $\operatorname{Proj}B$ are the sets $V_+(\mathfrak q)$ with $\mathfrak q$ a homogeneous prime ideal not containing $B_+$, that $V_+(\mathfrak q)$ is homeomorphic to $\operatorname{Proj}(B/\mathfrak q)$, and that $B/\mathfrak q$ is again a homogeneous $k$-algebra and a domain.

(a) \emph{Homogeneity.} The ideal $\mathfrak p^h=\bigoplus_d(\mathfrak p\cap B_d)\subseteq\mathfrak p$ is prime (Lemma~3.35(a)) and contains $f$ (a homogeneous element of $\mathfrak p$). By minimality $\mathfrak p=\mathfrak p^h$ is homogeneous. \emph{Containment.} A proper homogeneous ideal $\mathfrak a$ has $\mathfrak a_0$ a proper ideal of the field $B_0=k$, so $\mathfrak a_0=0$ and $\mathfrak a\subseteq B_+$. Hence $\mathfrak p\subseteq B_+$. \emph{Height.} $f$ is not invertible (it lies in $B_+$), so $\operatorname{ht}(\mathfrak p)\le1$ by Thm.~5.12.

Suppose $V_+(f)=\emptyset$. If $\mathfrak p\ne B_+$ then $\mathfrak p\subsetneq B_+$ does not contain $B_+$, so $\mathfrak p\in\operatorname{Proj}B$ and contains $f$, i.e. $\mathfrak p\in V_+(f)$, a contradiction. Hence $\mathfrak p=B_+$, and $\operatorname{ht}(B_+)\le1$.

To show $\dim X\le0$, let $Z=V_+(\mathfrak q)$ be an irreducible component of $X$ and $D=B/\mathfrak q$, a homogeneous $k$-domain with $\operatorname{Proj}D\simeq Z$ non-empty. The image $\bar f\in D_+$ satisfies $V_+(\bar f)=Z\cap V_+(f)=\emptyset$. By what we just proved applied to $D$, any prime of $D$ minimal over $\bar fD$ equals $D_+$, and $\operatorname{ht}(D_+)\le1$. By Prop.~5.23(a) (for the finitely generated domain $D$), $\dim D=\operatorname{ht}(D_+)+\dim D/D_+=\operatorname{ht}(D_+)\le1$, since $D/D_+=k$. By Cor.~5.21, $\dim Z=\dim D-1\le0$. Since $\dim X$ is the supremum of the dimensions of its irreducible components (Prop.~5.5(c)), $\dim X\le0$.

(b) \emph{Lemma.} Let $W\subseteq X$ be an irreducible closed subset with $\dim W\ge1$ and let $g\in B_+$ be homogeneous. Then $W\cap V_+(g)$ contains a non-empty irreducible closed subset of dimension $\ge\dim W-1$.

\emph{Proof of the Lemma.} Write $W=V_+(\mathfrak q)\simeq\operatorname{Proj}D$ with $D=B/\mathfrak q$ a homogeneous domain over $k$, and let $\bar g\in D_+$ be the image of $g$. Then $W\cap V_+(g)\simeq V_+(\bar g)\subseteq\operatorname{Proj}D$. By Cor.~5.21, $\dim D=\dim W+1\ge2$. If $\bar g=0$ then $W\cap V_+(g)=W$ and we are done. Otherwise, let $\mathfrak p_1,\dots,\mathfrak p_t$ be the minimal primes of $\bar gD$ (finitely many, Prop.~4.9); they are homogeneous and contained in $D_+$ by (a) applied to $D$. Each has height exactly $1$ (height $\le1$ by Thm.~5.12, and $\ne0$ because $D$ is a domain and $\bar g\ne0$). If all $\mathfrak p_j$ were equal to $D_+$ then $V_+(\bar g)=\emptyset$ (a homogeneous prime containing $\bar g$ contains some $\mathfrak p_j=D_+$), and part (a) applied to $D$ would give $\dim W\le0$, contrary to the hypothesis. So some $\mathfrak p_j\ne D_+$; then $\mathfrak p_j\subsetneq D_+$, so $\mathfrak p_j\in\operatorname{Proj}D$ and $V_+(\mathfrak p_j)\simeq\operatorname{Proj}(D/\mathfrak p_j)$ is a non-empty irreducible closed subset of $V_+(\bar g)$. By Prop.~5.23(a), $\dim D/\mathfrak p_j=\dim D-\operatorname{ht}(\mathfrak p_j)=\dim D-1$, and by Cor.~5.21 (as $D/\mathfrak p_j$ is a homogeneous algebra with non-empty $\operatorname{Proj}$), $\dim V_+(\mathfrak p_j)=\dim D/\mathfrak p_j-1=\dim D-2=\dim W-1$. This proves the Lemma.

Now let $Y$ be a closed subset of $X$ of dimension $r$ and $f_1,\dots,f_r\in B_+$ homogeneous. If $r=0$ the empty sequence gives $V_+()=X$ and $Y\neq\emptyset$, so the claim holds. Let $r\ge1$. $Y$ has finitely many irreducible components and $\dim Y$ is the maximum of their dimensions, so $Y$ contains an irreducible closed subset $W_0$ with $\dim W_0=r$. Inductively for $i=1,\dots,r$, given an irreducible closed $W_{i-1}\subseteq Y\cap V_+(f_1,\dots,f_{i-1})$ with $\dim W_{i-1}\ge r-i+1\ge1$, the Lemma applied to $W_{i-1}$ and $f_i$ gives a non-empty irreducible closed $W_i\subseteq W_{i-1}\cap V_+(f_i)$ with $\dim W_i\ge r-i$. Finally $W_r$ is non-empty and contained in $Y\cap V_+(f_1,\dots,f_r)$. Hence $V_+(f_1,\dots,f_r)\cap Y\ne\emptyset$. \hfill$\square$

\end{document}
